%% abtex2-modelo-trabalho-academico.tex, v<VERSION> laurocesar
%% Copyright 2012-2015 by abnTeX2 group at http://www.abntex.net.br/ 
%%
%% This work may be distributed and/or modified under the
%% conditions of the LaTeX Project Public License, either version 1.3
%% of this license or (at your option) any later version.
%% The latest version of this license is in
%%   http://www.latex-project.org/lppl.txt
%% and version 1.3 or later is part of all distributions of LaTeX
%% version 2005/12/01 or later.
%%
%% This work has the LPPL maintenance status `maintained'.
%% 
%% The Current Maintainer of this work is the abnTeX2 team, led
%% by Lauro César Araujo. Further information are available on 
%% http://www.abntex.net.br/
%%
%% This work consists of the files abntex2-modelo-trabalho-academico.tex,
%% abntex2-modelo-include-comandos and abntex2-modelo-references.bib
%%

% ------------------------------------------------------------------------
% ------------------------------------------------------------------------
% abnTeX2: Modelo de Trabalho Academico (tese de doutorado, dissertacao de
% mestrado e trabalhos monograficos em geral) em conformidade com 
% ABNT NBR 14724:2011: Informacao e documentacao - Trabalhos academicos -
% Apresentacao
% ------------------------------------------------------------------------
% ------------------------------------------------------------------------

\documentclass[
	% -- opções da classe memoir --
	12pt,				% tamanho da fonte
	openright,			% capítulos começam em pág ímpar (insere página vazia caso preciso)
	%twoside,			% para impressão em verso e anverso. Oposto a oneside
  oneside,     % para impressão em verso e anverso. Oposto a twoside
	a4paper,			% tamanho do papel. 
	% -- opções da classe abntex2 --
	%chapter=TITLE,		% títulos de capítulos convertidos em letras maiúsculas
	%section=TITLE,		% títulos de seções convertidos em letras maiúsculas
	%subsection=TITLE,	% títulos de subseções convertidos em letras maiúsculas
	%subsubsection=TITLE,% títulos de subsubseções convertidos em letras maiúsculas
	% -- opções do pacote babel --
	english,			% idioma adicional para hifenização
	french,				% idioma adicional para hifenização
	spanish,			% idioma adicional para hifenização
	brazil				% o último idioma é o principal do documento
	]{abntex2}

% ---
% Pacotes básicos 
% ---
%\usepackage{lmodern}			% Usa a fonte Latin Modern
\usepackage{cfr-lm}      % Technically, lmodern does since the fonts it installs are identical to those used by cfr-lm		
\usepackage[T1]{fontenc}		% Selecao de codigos de fonte.
\usepackage[utf8]{inputenc}		% Codificacao do documento (conversão automática dos acentos)
\usepackage{lastpage}			% Usado pela Ficha catalográfica
\usepackage{indentfirst}		% Indenta o primeiro parágrafo de cada seção.
\usepackage{color}				% Controle das cores
\usepackage{graphicx}			% Inclusão de gráficos
\usepackage{microtype} 			% para melhorias de justificação

%%%%%%%%%%%%%%%%%%%%%%
% pacotes incluídos
\usepackage{longtable}
\usepackage{lscape}
\usepackage{supertabular}

%\usepackage{fontspec}
%\usepackage{fontawesome}
\usepackage{hyperref}
%%%%%%%%%%%%%%%%%%%%%%
% ---
		
% ---
% Pacotes adicionais, usados apenas no âmbito do Modelo Canônico do abnteX2
% ---
\usepackage{lipsum}				% para geração de dummy text
% ---

% ---
% Pacotes de citações
% ---
\usepackage[brazilian,hyperpageref]{backref}	 % Paginas com as citações na bibl
\usepackage[alf]{abntex2cite}	% Citações padrão ABNT

% --- 
% CONFIGURAÇÕES DE PACOTES
% --- 

% ---
% Configurações do pacote backref
% Usado sem a opção hyperpageref de backref
\renewcommand{\backrefpagesname}{Citado na(s) página(s):~}
% Texto padrão antes do número das páginas
\renewcommand{\backref}{}
% Define os textos da citação
\renewcommand*{\backrefalt}[4]{
	\ifcase #1 %
		Nenhuma citação no texto.%
	\or
		Citado na página #2.%
	\else
		Citado #1 vezes nas páginas #2.%
	\fi}%
% ---

% ---
% Configurações de aparência do PDF final

% alterando o aspecto da cor azul
\definecolor{blue}{RGB}{41,5,195}

% definindo comandos para fontawesome

% \newfontfamily{\FA}{FontAwesome Regular}
% \def\twitter{{\FA \faTwitter}}


% informações do PDF
\makeatletter
\hypersetup{
     	%pagebackref=true,
		pdftitle={\@title}, 
		pdfauthor={\@author},
    	pdfsubject={\imprimirpreambulo},
	    pdfcreator={LaTeX with abnTeX2},
		pdfkeywords={abnt}{latex}{abntex}{abntex2}{trabalho acadêmico}, 
		colorlinks=true,       		% false: boxed links; true: colored links
    	linkcolor=blue,          	% color of internal links
    	citecolor=blue,        		% color of links to bibliography
    	filecolor=magenta,      		% color of file links
		urlcolor=blue,
		bookmarksdepth=4
}
\makeatother
% --- 

% --- 
% Espaçamentos entre linhas e parágrafos 
% --- 

% O tamanho do parágrafo é dado por:
\setlength{\parindent}{1.3cm}

% Controle do espaçamento entre um parágrafo e outro:
\setlength{\parskip}{0.2cm}  % tente também \onelineskip

% ---
% compila o indice
% ---
\makeindex
% ---

% ----
% Início do documento
% ----
\begin{document}

%%%%%%%%%%%%%%%%%%%%%%%%%%%%%%%%%%%%%%%%%%%%%%
%%%
%%% Capa e Contra-Capa
%%%
%%%%%%%%%%%%%%%%%%%%%%%%%%%%%%%%%%%%%%%%%%%%%%


%%%%%%%%%%%%%%%%%%%%%%%%%%%%%%%%%%%%%%%%%%%%%%
%%%
%%% Capa e Contra-Capa
%%%
%%%%%%%%%%%%%%%%%%%%%%%%%%%%%%%%%%%%%%%%%%%%%%


\begin{center} 
    \includegraphics[width=2.57cm]{./images/brasao_da_republica.jpeg} \\
    
    {\sffamily
    Ministério da Educação\\
    Universidade Federal do Agreste de Pernambuco\\
    Pró-Reitoria de Ensino de Graduação
    }
          
    \vfill
    % \vspace{4.5cm}
    {{\Large Projeto Pedagógico de Curso}}\\
    \vspace{.5cm}
    {{\Large \textbf{Bacharelado em Ciência da Computação}}}\\
    \vspace{1cm}
    {{\large Perfil: 2.1.0}}\\
    \vfill
    
    %\includegraphics[width=1cm]{images/uagUFRPE.eps} \\
    
    % {\scshape Pró-Reitoria de Ensino de Graduação\\}
    % {\scshape Unidade Acadêmica de Garanhuns\\}
    % {\scshape Bacharelado em Ciência da Computação\\}
    % Av. Bom Pastor, s/n, Boa Vista, CEP 55292-270, Garanhuns, PE\\
    % Telefone: (87) 3764-5500\\
    % \url{http://bcc.uag.ufrpe.br}\\ \vskip\baselineskip
    
    \vspace{1cm}
    \vskip\baselineskip
    Garanhuns, 2022.
\end{center}

\newpage

%%%%%%%%%%%%%%%%%%%%%%%%%%%%%%%%%%%%%%%%%%%%%%%%%%%%%%%%
%%%
%%% Comissão de Elaboração da Reforma Curricular do Curso
%%%
%%%%%%%%%%%%%%%%%%%%%%%%%%%%%%%%%%%%%%%%%%%%%%%%%%%%%%%%

%%%%%%%%%%%%%%%%%%%%%%%%%%%%%%%%%%%%%%%%%%%%%%%%%%%%%%%%
%%%
%%% Comissão de Elaboração da Reforma Curricular do Curso
%%%
%%%%%%%%%%%%%%%%%%%%%%%%%%%%%%%%%%%%%%%%%%%%%%%%%%%%%%%%

\begin{center}
  \textbf{PROJETO PEDAGÓGICO DO CURSO DE\\BACHARELADO EM CIÊNCIA DA COMPUTAÇÃO}

  \vspace{36pt}

  \textbf{REITOR} \\
  Airon Aparecido Silva de Melo
  
  \vspace{16pt}

  \textbf{PRÓ-REITORA DE ENSINO E GRADUAÇÃO} \\
  Emanuelle Camila Moraes de Melo Albuquerque Lima
  
  \vspace{16pt}
  
  \textbf{PRÓ-REITOR DE PESQUISA, PÓS-GRADUAÇÃO E INOVAÇÃO}\\
  José Romualdo de Sousa Lima
  
  \vspace{16pt}
  
  \textbf{PRÓ-REITOR DE EXTENSÃO E CULTURA}\\
  Marcos Pinheiro Franque
  
  \vspace{16pt}

  \textbf{PRÓ-REITOR DE ADMINISTRAÇÃO}\\
  José Renato Correia Ferro
  
  \vspace{16pt}

  \textbf{PRÓ-REITOR DE PLANEJAMENTO}\\
  Victor Netto Maia
  
  \vspace{16pt}

  \textbf{COORDENADORA DO CURSO}\\
  Thaís Alves Burity Rocha

  \vspace{16pt}

  \textbf{PEDAGOGA}\\
  Maria Gorete Rodrigues de Siqueira

  \vspace{16pt}

  \textbf{COORDENADORIA DE REGULAÇÃO DOS CURSOS}\\
  Amanda Maria Rodrigues Diniz

  \vspace{16pt}

\newpage

  \textbf{REFORMULAÇÃO E ATUALIZAÇÃO DO PPC}

  \vspace{16pt}

  Thaís Alves Burity Rocha\\
  Presidente do NDE

  \vspace{16pt}

  Tiago Buarque Assunção de Carvalho\\
  Vice-presidente do NDE
  
  \vspace{16pt}

  Dimas Cassimiro do Nascimento Filho\\
  Membro do NDE
  
  \vspace{16pt}

  Hélder Fernando de Araújo Oliveira\\
  Membro do NDE
  
  \vspace{16pt}

  Ícaro Lins Leitão da Cunha\\
  Membro do NDE

  \vspace{16pt}

  Rodrigo Gusmão de Carvalho Rocha\\
  Membro do NDE

  \vspace{16pt}

  Sérgio Francisco Tavares de Oliveira Mendonça\\
  Membro do NDE
\end{center}

\newpage


%%%%%%%%%%%%%%%%%%%%%%%%%%%%%%%%%%%%%%%%%%%%%%%%%%%%%%%%
%%%
%%% Dados da Criação/Autorização/Reconhecimento do Curso
%%%
%%%%%%%%%%%%%%%%%%%%%%%%%%%%%%%%%%%%%%%%%%%%%%%%%%%%%%%%

%%%%%%%%%%%%%%%%%%%%%%%%%%%%%%%%%%%%%%%%%%%%%%
%%%
%%% Identificação do Curso
%%%
%%%%%%%%%%%%%%%%%%%%%%%%%%%%%%%%%%%%%%%%%%%%%%

\chapter*[Dados de Identificação do Curso]{DADOS DE IDENTIFICAÇÃO DO CURSO}
\addcontentsline{toc}{chapter}{Dados de Identificação do Curso}


\begin{center}
\begin{footnotesize} 
\begin{longtable}{lp{9.4cm}}
  \caption{\label{quadro:sintese-do-curso}Síntese do curso.}\\
  \toprule
  \multicolumn{2}{c}{\bfseries SÍNTESE DO CURSO} \\
  \midrule
  \textbf{Modalidade} & Presencial \\ \midrule
  \textbf{Denominação do Curso}	& Ciência da Computação \\ \midrule
  \textbf{Habilitação}	& Bacharelado \\ \midrule
  \textbf{Local de Oferta}	& Universidade Federal do Agreste de Pernambuco \newline Av. Bom Pastor, s/n, Boa Vista \newline CEP 55292-270 - Garanhuns/PE \\ \midrule
  \textbf{Turno(s) de Funcionamento}	 & Noturno \\ \midrule
  \textbf{Número de Vagas}	& 80 vagas anuais \\ \midrule
  \textbf{Periodicidade de Oferta}	& Semestral \\ \midrule
  \textbf{Carga Horária Total} & 3.200 horas \\ \midrule
  \textbf{Período de Integralização}	& 4,5 anos (9 semestres) \\ \midrule
  \textbf{Período Máximo de Integralização} 	& 4,5 anos + 70\% desse tempo (15 semestres)\footnote{Resolução UFRPE/CEPE nº 154/2001.} \\ \midrule
  \textbf{Ato Regulatório do Curso} &	Portaria nº 9.222, de 27 de dezembro de 2018. \\ \midrule
  \textbf{Mantida}	& Universidade Federal Rural de Pernambuco UFRPE \newline Pessoa Jurídica de Direito Público - Federal \newline R. Dom Manoel de Medeiros, s/n - Dois Irmãos. Recife - PE \\ \midrule
  \textbf{Corpo Dirigente da UFAPE}	& Nome: Airon Aparecido de Melo \newline Cargo: Reitor \newline Telefone: 0xx(87) 3764-5505 \newline E-mail: airon.melo@ufape.edu.br \\
\bottomrule
\end{longtable}
\end{footnotesize}
\end{center}

\newpage

% Seleciona o idioma do documento (conforme pacotes do babel)
%\selectlanguage{english}
\selectlanguage{brazil}

% Retira espaço extra obsoleto entre as frases.
\frenchspacing 

% ---
% inserir lista de ilustrações
% ---
\pdfbookmark[0]{\listfigurename}{lof}
\listoffigures*
\cleardoublepage
% ---

% ---
% inserir lista de tabelas
% ---
\pdfbookmark[0]{\listtablename}{lot}
\listoftables*
\cleardoublepage
% ---

% ---
% inserir lista de abreviaturas e siglas
% ---
\begin{siglas}
  \item[UFRPE] Universidade Federal Rural de Pernambuco
  \item[UAG] Unidade Acadêmica de Garanhuns
  \item[BCC] Bacharelado em Ciência da Computação
\end{siglas}
% ---

% ---
% inserir lista de símbolos
% ---
% \begin{simbolos}
%   \item[$ \Gamma $] Letra grega Gama
%   \item[$ \Lambda $] Lambda
%   \item[$ \zeta $] Letra grega minúscula zeta
%   \item[$ \in $] Pertence
% \end{simbolos}
% ---

% ---
% inserir o sumario
% ---
\pdfbookmark[0]{\contentsname}{toc}
\tableofcontents*
\cleardoublepage
% ---



% ----------------------------------------------------------
% ELEMENTOS TEXTUAIS
% ----------------------------------------------------------
\textual

% ----------------------------------------------------------
% Introdução (exemplo de capítulo sem numeração, mas presente no Sumário)
% ----------------------------------------------------------
\chapter*[Introdução]{Introdução}
\addcontentsline{toc}{chapter}{Introdução}
% ----------------------------------------------------------

\section*{A Universidade Federal Rural de Pernambuco}

A história da Universidade Federal Rural de Pernambuco (UFRPE) tem início no 
dia 3 de novembro de 1912, quando foi lançada a pedra fundamental do
edifício que abrigaria as Escolas Superiores de Agricultura e Medicina
Veterinária, na cidade de Olinda, pelo Revmo. Abade do Mosteiro de São Bento,
Dom Pedro Roeser. A trajetória da Universidade vem, desde então, caminhando
para a formação do que é a UFRPE hoje, instituição que engloba nove campi
em todo o estado de Pernambuco que oferecem 19 (dezenove) cursos de Graduação,
20 (vinte) de Pós-Graduação, sendo 6 (seis) de Doutorado e 14 (quatorze) de
Mestrado, além dos cursos de Especialização, Aperfeiçoamento e Extensão. Através do CODAI -
Colégio Agrícola Dom Agostinho Ikas, a Universidade oferece ainda o curso Técnico Agrícola em nível médio. Recentemente, foi incorporado à oferta da
UFRPE o Curso de Licenciatura em Física, Licenciatura em Computação e Sistemas
de Informação na modalidade de  Educação à Distância.

Tradicionalmente a UFRPE tem a missão da interiorização de cursos de Graduação, e também 
no Estado de Pernambuco com os seus \textit{Campi} Avançados,  sendo esta a primeira Universidade Pública Federal do Brasil a ter uma extensão universitária.

\subsection*{A Unidade Acadêmica de Garanhuns}

A Unidade Acadêmica de Garanhuns (UAG) foi a primeira extensão universitária a ser instalada no país, tendo suas atividades iniciadas no segundo semestre de 2005 com os cursos de Agronomia, Licenciatura Licenciatura em Pedagogia, Medicina Veterinária e Zootecnia. A escolha dos cursos oferecidos foi feita tendo em vista as características do entorno da UAG, levando em consideração que o papel central da universidade é de influir no desenvolvimento regional.

\subsection*{Organização da estrutura acadêmica}

A administração do curso está estruturada da seguinte forma:

\begin{enumerate}
\item Departamento (Unidade Acadêmica de Garanhuns): unidade básica da estrutura da Universidade para efeito de organização didático-científica e 
  administrativa, que integra as atividades de Ensino, Pesquisa e Extensão e
  áreas afins do conhecimento e respectivos docentes e discentes, sendo
  administrado pelo Diretor Geral/Acadêmico e Diretor Administrativo.
\item Conselho de Ensino, Pesquisa e Extensão (CEPE): órgão deliberativo em
matéria de Ensino, Pesquisa e Extensão, em sua área de conhecimento, lotado na
Sede em Recife-PE, Campus de Dois Irmãos.
\item Chefia de Departamento (Unidade): A
  Unidade possui um Diretor Geral eleito na forma das normas eleitorais da
  UFRPE, pelos professores, técnicos e estudantes da Unidade, empossado pelo
  Reitor para um mandato de 4 (quatro) anos, conforme estatuto da UFRPE e
  Regimento da Unidade, o qual nomeará o Coordenador Geral de Cursos. O diretor
  administrativo será nomeado pelo Reitor.
\item Colegiado de Curso: órgão deliberativo em matéria de Ensino, Pesquisa e
Extensão,  no âmbito do curso, sendo composto pelos docentes e representantes
estudantis, no limite máximo do Estatuto da UFRPE.
\item Coordenador de Curso: responsável pela supervisão das
  atividades acadêmicas do curso, eleito na forma das normas eleitorais da
  Universidade, empossado pelo Reitor para um mandato de 2 (dois) anos.
\end{enumerate}

\section*{A Informática e Computação na UFRPE}

Em Informática e Computação a história começa em 1999, quando o Curso de
Licenciatura em Computação da UFRPE foi instituído sendo a primeira oferta
vestibular em 2000. Legalmente, o Curso encontra-se Autorizado segundo Resolução
CEPE 265/1999, implantado segundo Resolução CUNI no. 181/99. Foi Reconhecido
junto ao MEC/INEP em novembro de 2005 e encontra-se em aguardo da Publicação da
Portaria em D.O.U. O Projeto Pedagógico em vigor encontra-se homologado segundo
Resolução 90/2004 do CUNI. Recentemente, dezembro de 2006, foi avaliado pelo
INEP/MEC, o qual recebeu o conceito CONDIÇÕES BOAS, apesar de recém-criado.
O curso de Licenciatura em Computação foi criado pois havia na UFRPE poucos
docentes da Área de Computação mas havia um consolidado Departamento de
Educação. Assim, naquele momento, o perfil mais adequado era realmente o de
Licenciatura em Computação. Neste cenário, há na UFRPE desde 2000 a
competência em Informática e Computação instalada e em 2005 seu quadro
docente completou 6 Doutores em Computação, referendando a competência
instalada em 2000. Além disso, a UFRPE conta hoje com o núcleo de Educação à Distância e Inclusão Digital e Social com projetos de pesquisa aprovados na
FACEPE. Sendo o curso de Licenciatura em Física seu 1. Projeto de Curso na
modalidade a Distância, aprovado segundo Portaria no.3.726  de  21/10/2005,
publicada em 24/10/05 do D.O.U. Desta forma, a UFRPE sentiu-se confortável
em sugerir a criação do Curso de Bacharelado em Ciência da Computação no
Município de Garanhuns.



\section*{Contexto de inserção do curso na região}

A ideia de criar um curso na área de computação existe desde a concepção da
Unidade Acadêmica de Garanhuns em setembro de 2005, quando começaram a funcionar
4 cursos de graduação: Agronomia, Licenciatura Normal Superior (transformada no
curso de Licenciatura em Pedagogia), Veterinária e Zootecnia. Atualmente a UAG
funciona em sede própria e oferece além desses cursos o curso de especialização
em Educação e Desenvolvimento Sustentável. Em reunião geral datada de dezembro
de 2007, ficou decidido, em processo de votação, que seria proposta a criação de
3 (três) novos cursos, dentro do processo de Reestruturação Universitária
(REUNI). Entre eles foi indicado o curso de Ciência da Computação no turno noturno com o
objetivo tanto de suprir a necessidade de um curso na área de informática,
quanto visando o desenvolvimento acadêmico da Universidade e uma forte interação
com os demais cursos de graduação da UAG.

Além de interagir com as demais áreas da UAG, o curso de Bacharelado em Ciência
da Computação, vem atender a uma demanda regional identificada tanto junto ao
poder público local e população. Portanto o curso de BCC se
inserirá dentro do contexto dos demais cursos da área de computação da UFRPE, de
forma a contribuir com o desenvolvimento da UAG e dentro da realidade local.
Para tanto foram definidas áreas de atuação dos profissionais do curso bem como
áreas de conhecimento que trouxessem essas características ao curso.

Pode-se constatar que o uso do computador deixou de ser um diferencial para se
tornar necessidade fundamental, tanto no contexto profissional quanto no dia a
dia das pessoas. O advento da \textit{Internet} transformou o computador em elemento
chave na construção da chamada sociedade da informação, modificando inclusive a
forma de relacionamento na sociedade moderna. Dados de setembro de 2008
demonstram que existem no mundo cerca de 1,5 bilhão de usuários da
\textit{Internet}, no Brasil a rede atende a aproximadamente 26\% da população,
índice muito aquém da média dos países com melhor nível de desenvolvimento onde a internet atende a
mais de 60\% da população (dados de março de 2008).

O Brasil sofre com graves problemas tanto no acesso da população aos recursos
computacionais quanto nas desigualdades regionais. Juntamente com a \textit{Internet}
surgem novas oportunidades  de desenvolvimento ligadas à produção de conteúdo
para a rede o desenvolvimento de sistemas que usem grande quantidade de dados. Neste aspecto, é urgente a
formação de profissionais ligados ao desenvolvimento de \textit{Software}. Em
2006 a Sociedade Brasileira de Computação (SBC) definiu cinco grandes desafios atuais
da computação:

\begin{enumerate}
  \item Gestão da informação em grandes volumes de dados multimídia distribuídos;
  \item Modelagem computacional de sistemas complexos artificiais, naturais,
  socioculturais e da interação homem-natureza;
  \item Impactos na computação da transição do silício para novas tecnologias;
  \item Acesso participativo universal do cidadão brasileiro ao conhecimento; e,
  \item Sistemas disponíveis, corretos, seguros, escalonáveis, persistentes e ubíquos.
\end{enumerate}


\section*{Motivação (e Justificativa)}

Dentro desses desafios podemos contextualizar o curso de BCC da UAG em uma
região carente de profissionais na área de desenvolvimento de software.

A região onde se encontra a UAG tem uma economia com base na agropecuária, e o
município de Garanhuns tem uma forte atuação no setor de serviços, com forte
apelo para o uso da informática. 

O curso também irá se inserir de forma bastante  ativa no projeto intitulado
Garanhuns Digital que tem por objetivo implantar  \textit{Internet} sem fio em
escolas, empresas e repartições públicas. 

Assim, a informática  é também um dos eixos norteadores do desenvolvimento
municipal.  O programa de expansão das universidades federais centra-se na
possibilidade de  responder às demandas regionais sem, no entanto,
restringir-se apenas à região,  mas produzindo e  transferindo conhecimentos
além desta, função inerente a toda Universidade.

Portanto, o curso de Bacharelado em Ciência da Computação foi projetado com
eixos fundamentados em áreas do conhecimento que viessem a contribuir no
desenvolvimento regional.

\section*{Objetivo do curso}

O principal objetivo do curso de Bacharelado em Ciência da Computação, da
Unidade Acadêmica de Garanhuns, é o desenvolvimento científico, cultural e
econômico da região, contribuindo assim com a sua inserção em
contexto nacional e internacional. Os objetivos específicos do curso são:

\begin{enumerate}
  \item Desenvolver nos estudantes o perfil  científico de pesquisador tanto
  para atuação na área acadêmica quanto para atuação em outros ramos de atividade;
  \item Desenvolver nos estudantes um espírito  empreendedor, incentivando e
  motivando a sua independência e criatividade;
  \item Promover a interdisciplinaridade buscando  atualização constante na área
  de computação;
  \item Motivar e orientar o estudante para que ele  tenha uma postura ativa
  diante da necessidade de um aprendizado contínuo e autônomo;
  \item Promover uma postura ética e socialmente  comprometida com o papel do
  estudante no desenvolvimento científico, tecnológico,  social e econômico da
  sua região e do País.
  \item Promover interação constante com escolas do  ensino fundamental e médio
  local de forma a estimular vocações e colaborar ativamente com a melhoria da educação;
  \item Estabelecer metas e realizar um processo de avaliação constante do BCC
  da UAG\-UFRPE.
\end{enumerate}

\section*{Planejamento Estratégico para o Curso na UAG, UFRPE, no Estado e
no País}

Aqui apresentam-se as metas e estratégias do grupo de profissionais envolvidos com a área de
computação na UAG, UFRPE. Tais objetivos para o futuro do Curso de Bacharelado em
Ciência da Computação, para a UAG, UFRPE e para o Estado de Pernambuco serão tomados como
guias e linhas-mestre para as ações destes indivíduos nos próximos 10 anos.

\subsection*{Metas e Estratégias para 2 anos}

\begin{itemize}
  \item Graduação forte e consolidada. Formação de corpo docente adequado,
reestruturação curricular para acompanhar as evoluções pedagógicas e 
tecnológicas, infra-estutura de trabalho para os docentes, discentes e 
técnico-administrativos envolvidos no curso e isto inclui apoio a criação de
 empresa júnior e centro acadêmico pelos discentes do curso;
  \item Implantação de curso de Especialização. Criação de curso de
especialização visando a capacitação  de  alunos do curso de  Bacharelado em Ciência da Computação e afins, nos diversos núcleos do conhecimento da área de Computação;
  \item Consolidação de grupos de pesquisa, nas áreas de competência dos membros do colegiado do curso.
  \item Consolidação de infraestrutura operacional. Para suportar convênios e
alavancar novos   investimentos, uma infraestrutura operacional deve estar completamente
definida  e em operação para instalação de laboratórios de pesquisa, treinamento
de educadores,  convênios. Envolvimento dos professores de computação no
Departamento de Tecnologia da Informação, UAG, UFRPE, para  prover a instituição
de políticas para aquisição de hardware, software, instalação  e manutenção de equipamentos e sistemas, automação de
serviços administrativos e acadêmicos.
\end{itemize}

\subsection*{Metas para 5 anos}

\begin{itemize}
  \item Corpo docente altamente qualificado com 100\% dos professores com
doutorado ou afastados em programas de doutoramento;
  \item Corpo docente altamente qualificado com 100\% dos professores com
doutorado engajados em projetos de pesquisa e/ou extensão;
  \item Implantação de Mestrado em Ciência da Computação;
  \item Curso de Graduação entre os 25 melhores do Brasil;
  \item Estudo de viabilidade para captação de recursos, além de convênios, para
investimentos do grupo;
  \item Interação com outros centros de excelência em Computação e áreas relacionadas.
\end{itemize}

\subsection*{Metas para 10 anos}

\begin{itemize}
  \item Implantação de Doutorado em Ciência da Computação;
  \item Curso de Graduação entre os 15 melhores do Brasil;
  \item Autonomia financeira para investimentos no curso e no campus da
UAG/UFRPE;
\end{itemize}


%%%%%%%%%%%%%%%%%%%%%%%%%%%%%%%%%%%%%%%%%%%%%%%%%%%%%%%%%%%%%%%%%%%%%%%%%%%%%%%%%%%%%
%
% chapter A área da ciência da computação
%
%%%%%%%%%%%%%%%%%%%%%%%%%%%%%%%%%%%%%%%%%%%%%%%%%%%%%%%%%%%%%%%%%%%%%%%%%%%%%%%%%%%%%%


\chapter{A área da Ciência da Computação}

De um modo geral, nós podemos definir computação como uma atividade que usa o
computador para atingir seu objetivo ou meta. Assim, computação inclui
construção e implementação de projetos de hardware e software para uma extensa gama de propósitos; processando, estruturando, e administrando diversos tipos de sistemas de informação, que consiste em usar o computador para estudos científicos, desenvolvendo sistemas
inteligentes; comunicação e mídia de
entretenimento, colhendo informações pertinentes. Computação também pode ter outros significados que são mais específicos, baseados no contexto no qual o termo é empregado.

\section{Categorias de cursos da área de computação e informática}

Com as diretrizes curriculares de 1999 foi criada a denominação da área de
computação e informática orientando a elaboração do projeto político pedagógico
dentro do tipo de curso escolhido. Assim, foram limitadas as possibilidades de
nomes de cursos dessa área a cinco tipos: Bacharelado em Ciência da Computação,
Engenharia da Computação, Bacharelado em Sistemas da Informação, Cursos de
Licenciatura em Computação e Cursos Superiores Tecnológicos. Segundo as
diretrizes esses cursos se enquadram em quatro categorias básicas:

\begin{enumerate}
  \item Cursos que têm a computação como atividade-fim: Ciência da Computação e
  Engenharia da Computação;
  \item Cursos que têm a informática como atividade-meio: Sistema da Informação;
  \item Cursos voltados para o ensino da informática: Licenciatura em
 Computação; e,
  \item Cursos Tecnológicos e sequênciais.
\end{enumerate}

Para que o curso escolhido se inserisse melhor dentro do desenvolvimento da UAG
optou-se pelo curso de Bacharelado em Ciência da Computação (BCC) que se
enquadra na categoria de curso com a computação como atividade-fim.

O curso de BCC da UAG foi idealizado a partir do currículo de referência
formulado em documento de 2005 pela IEEE \textit{Computer Society} e levando em conta as
tendências e desafios para a área de informática descrita em publicação recente
sobre a trajetória dos cursos de graduação da área de computação e informática
publicada pela Sociedade Brasileira de Computação. A matriz curricular
construída a partir do estudo de projetos de cursos de outras Instituições de
Ensino Superior, de alto nível, seguindo as recomendações do Currículo de
Referência da Sociedade Brasileira de Computação (1999).
	
\section{Áreas de formação}

De acordo com as diretrizes curriculares para os cursos das áreas de computação
e informática publicadas pelo MEC-SESU, esses cursos devem abranger quatro
grandes áreas de formação:
\begin{enumerate}
  \item Formação Básica;
  \item Formação Tecnológica;
  \item Formação Complementar; e,
  \item Formação Humanística.
\end{enumerate}


%%%%%%%%%%%%%%%%%%%%%%%%%%%%%%%%%%%%%%%%%%%%%%%%%%%%%%%%%%%%%%%%%%%%%%%%%%%%%%%%%%%%%
%
% chapter Fundamentos norteadores do curso
%
%%%%%%%%%%%%%%%%%%%%%%%%%%%%%%%%%%%%%%%%%%%%%%%%%%%%%%%%%%%%%%%%%%%%%%%%%%%%%%%%%%%%%%


\chapter{Fundamentos norteadores do curso}

A humanidade é caracterizada no aprender e no expressar os aprendizados
adquiridos para a sociedade sem, contudo, seguir modelos e fórmulas rígidas.
Neste sentido, a educação faz parte da construção e do cerne da vida para o
crescimento de um grupo socialmente construído a partir de crenças e idéias.


\section{Fundamentos ético-políticos}


Sob este pressuposto, a formação do Bacharel em Ciência da Computação da UFRPE
norteia-se na qualidade de ser cidadão íntegro e emancipado politicamente, capaz
de conduzir e posicionar-se diante de fatos, de forma coerente diante de uma
sociedade complexa e competitiva.

As proposições didático-pedagógicas para a efetivação dos pressupostos
ético-políticos fundamentam-se na justiça, respeito mútuo, participação,
diálogo, reflexão, responsabilidade, solidariedade, dignidade humana e ética com
a natureza.

Neste sentido, o Curso foi estruturado para que o estudante, como cidadão, além
de estar apto a atuar na sua profissão, seja capaz de refletir, entender e
valorizar a dimensão humana bem como da capacidade da natureza relacionada com a
Ciência e Tecnologia.

O Bacharel em Ciência da Computação não deverá apresentar apenas uma formação
voltada para o atendimento das demandas do exercício profissional específico,
mas saber fazer do uso de seu conhecimento, transformando-o em ações
responsáveis socialmente.

\section{Fundamentos epistemológicos}


A função do Curso na sociedade expressa-se em afirmações, tais como: ``para
preparar os estudantes para o exercício profissional e o ser cidadão'', ``para
passar aos estudantes os conhecimentos construídos pelas sociedades ao longo dos
anos''. Tais assertivas sugerem um denominador comum, ou seja, no Curso
trabalha-se com o conhecimento. Disso derivam outros questionamentos: O que é o
conhecimento? Como ele se produz? Como as pessoas dele se apropriam? As
respostas resultantes, no decurso da história da humanidade deram origem às
várias correntes epistemológicas, isto é, às diferentes teorias ligadas ao
conhecimento.

Uma destas teorias está ligada ao processo de construção do conhecimento. Nela,
o conhecimento é visto como resultado de uma interação entre o sujeito que quer
conhecer o objeto a ser conhecido. Trata-se de uma interação dinâmica, pois à
medida que o sujeito age sobre o objeto do conhecimento, ele o transforma e se
transforma. Refere-se a objetos do conhecimento, isto é, conceitos, idéias e
definições que são construções sociais, existentes nas mentes das pessoas e que
possibilitam identificar características e formas de porções da matéria do mundo
físico e de fenômenos do mundo social.

Por sua vez, ao interagir com os objetos de conhecimento já socialmente
construídos, o sujeito lhes atribui interpretação própria, modifica-se
refletindo sobre suas características, e em consequência, modifica-se a si
mesmo. Esse processo dinâmico e contínuo de transformações, ou seja, de
construção e reconstrução do conhecimento, define o princípio básico do
Construtivismo, no qual o estudante tem o papel de construir e reconstruir seu
conhecimento.

Por outro lado, a interação se dá no interior do sujeito que quer aprender, por
meio das habilidades e estruturas mentais desenvolvidas pelo mesmo. Trata-se de
uma interação realizada por intermédio do meio físico e social.

Nesta perspectiva, o sujeito da aprendizagem é histórico e social e o objeto do
conhecimento é cultural. O primeiro, porque considera o sujeito inscrito nos
valores e no momento histórico de sua comunidade. O segundo, porque é construído
pela cultura de cada grupo social. Deste pressuposto, o Construtivismo é
considerado interracionista.

Neste sentido, a base epistemológica do Curso se dará no exercício da construção
de conhecimentos que, além de gerar desenvolvimento, também esteja voltado para
a satisfação das necessidades sociais e o respeito com os recursos da natureza,
tendo em vista as gerações futuras.

\section{Fundamentos didático-pedagógicos}

A consistência da proposta pedagógica reside em orientar e fornecer os meios
específicos à elaboração e à conservação de produtos alimentícios para a
industrialização, via conhecimento e aplicação de técnicas e operações de
natureza física, química e biológica, não deixando de lado a percepção de que o
processo econômico deve servir-se da natureza, mas, de forma mais duradoura,
levando em conta a limitação ecológica imposta pela natureza no processo
econômico de produção. Neste sentido, para que uma proposta pedagógica ganhe
maior sentido, deve colocar o desenvolvimento promovido dentro da moldura da
ecosfera.

Tendo-se por concepção que o conhecimento é o elo capaz de firmar com qualidade
os significados das ações realizadas pelo homem no mundo, este deve ocorrer de
forma efetiva, e para que isto seja possível, a ação reflexiva deve permear as
atividades didático-pedagógicas na formação acadêmica.

Para isto, a formação do Bacharel em Ciência da Computação da UFRPE tem por base
uma política de desenvolvimento que permite uma formação ética para elevar as
condições de vida das pessoas, sem perturbar funções ecossistêmicas essenciais,
levando-se em conta o fazer, e o que é possível em face do que é desejável. A
proposta pedagógica do Curso norteia-se em fundamentos que abrangem as dimensões
ético-políticas, epistemológicas e didático-pedagógicas.

Para direcionar as atividades de ensino, propriamente ditas, é mister uma opção
epistemológica. O ecletismo pode falsear a relação entre professor e estudante.
Assim, se a opção for pela visão de conhecimento como uma representação do
mundo, o professor buscará, em primeiro lugar, as concepções prévias do
estudante sobre o assunto que será estudado para, depois, propor uma explicação
escolar.

Nesta abordagem, o professor está mais interessando em saber o que o estudante
já sabe, para servir-lhe de âncora ao conhecimento acadêmico a ser proposto. Por
sua vez, os estudantes ressignificarão as mensagens do professor dentro do
contexto de suas estruturas cognitivas, construídas no cotidiano individual.

Considerando que as relações entre o estudante e o professor na escola são
orientadas pela Pedagogia, cujo foco de trabalho é a educação, o Curso de
Bacharelado em Ciência da Computação da UFRPE fundamenta-se nos princípios do
Construtivismo Sócio-Interacionista.

Nesta abordagem, o conhecimento passa a ser visto como um conjunto de verdades
relativas, que correspondem à uma interpretação que o homem dá ao mundo físico e
social.

O professor exerce o papel de catalisador (mediador) do processo de interação
que ocorre entre o sujeito da aprendizagem (o estudante) e o objeto do
conhecimento social (o conhecimento social compartilhado). Ensinar, nesta visão,
é preparar as melhores condições para que possa haver aprendizagem. Em
consequência, cabe ao professor conhecer seus estudantes, interagir com eles,
buscando sua história e permitir-lhes que manifestem suas concepções prévias
diante dos assuntos a serem estudados.

Por sua vez, o estudante deixa de ser um mero receptor de informações, passando
a ser um construtor, numa Pedagogia inspirada nos princípios da construção do
conhecimento, própria do sujeito que pauta o seu fazer pela pesquisa, pela
interrogação e pela problematização. Em outros termos, o conhecimento se
constitui nas relações que cada sujeito estabelece, frente às interpretações que
o professor lhe faz de um saber construído e aceito socialmente. Assim, o
processo de aprendizagem ocorre de dentro para fora, ou seja, é o próprio
estudante que, a partir de sua experiência de vida, de seu próprio universo
simbólico fará uma interpretação do ``saber oficial'', interpretação esta que
deverá compartilhar ao máximo com outros membros da sociedade.

Uma Pedagogia inspirada nos pressupostos ora apresentados faz do Curso de
Bacharelado em Ciência da Computação da UFRPE um processo formativo com visão
reflexiva, natural do processo educativo, objetivando oferecer as condições para
o desenvolvimento harmonioso dos estudantes nos domínios cognitivos, afetivos,
psicológicos, biológicos e sociais.




%%%%%%%%%%%%%%%%%%%%%%%%%%%%%%%%%%%%%%%%%%%%%%%%%%%%%%%%%%%%%%%%%%%%%%%%%%%%%%%%%%%%%
%
% chapter Pressupostos metodológicos
%
%%%%%%%%%%%%%%%%%%%%%%%%%%%%%%%%%%%%%%%%%%%%%%%%%%%%%%%%%%%%%%%%%%%%%%%%%%%%%%%%%%%%%%




\chapter{Pressupostos metodológicos}

O Curso de Bacharelado em Ciência da Computação da UFRPE estabelece como
pressupostos metodológicos: 

\section{Relação teoria-prática}

A relação teoria-prática pode ser entendida como eixo articulador da produção do
conhecimento, servindo para o acadêmico vislumbrar possibilidades futuras de engajamento no mercado de
trabalho bem como potencializando o aprendizado teórico em si. Abandona-se aqui
a idéia de que primeiro o estudante precisa dominar a teoria para depois
entender a prática e a realidade, resultando em um aprendizado memorístico.

Busca-se a construção do conhecimento de forma ampla, muitas vezes integrando,
numa mesma situação teoria e prática. Além disso, sustenta-se a idéia de que
relacionar teoria e prática não consiste em atividade exclusiva de sala de aula,
devendo-se proporcionar ao acadêmico, desde o primeiro semestre, atividades
incluídas na carga horária semanal das diferentes disciplinas que compõem a
matriz curricular bem como atividades complementares que contribuam
indiretamente à compreensão do Curso e de sua contribuição na sociedade como um
todo. Desta forma, além das atividades apresentadas na matriz curricular, as
atividades complementares definidas para os acadêmicos do Curso de Bacharelado
em Ciência da Computação da UFRPE servirão de meio para atingir a desejada
capacidade de relacionar teoria e prática. 

\section{Interdisciplinaridade}

Considera-se que para se atingir o perfil de Bacharel em Ciência da Computação com sólida
formação generalista necessita-se a realização de estudos disciplinares que
permitam a sistematização e o aprofundamento de conceitos e relações, cujo
domínio é imprescindível na construção da competência profissional desejada. No
entanto, sabe-se que a construção de um conhecimento sólido transpõe o conteúdo
de uma única disciplina, necessitando que o acadêmico primeiramente tenha
conhecimento da contextualização da disciplina específica no todo e que, num
segundo momento, desenvolva atividades que necessitem dos conteúdos expostos em
várias disciplinas, tornando possível aplicar conhecimentos adquiridos ao longo
de todo o Curso no desenvolvimento de uma atividade específica. Desta forma,
além de aprofundar conhecimentos disciplinares, a matriz curricular contempla
estudos e atividades interdisciplinares, propostas ao longo do curso por
diferentes disciplinas. Além das atividades interdisciplinares formais, várias
atividades são desenvolvidas por disciplinas afins concomitantemente,
proporcionando o aprendizado não intencional e aplicação de conceitos
complementares, transcendendo, desta forma, os limites de sala de aula. 

\section{Pesquisa enquanto princípio educativo}

A pesquisa, compreendida como qprocesso formador, é um elemento constitutivo e
fundamental do processo de aprender a aprender aprendendo e, portanto,
prevalente em vários momentos curriculares. A familiaridade com a teoria só pode
se dar por meio do conhecimento das pesquisas que lhe dão sustentação. De modo
semelhante, a atuação prática possui uma dimensão investigativa e constitui uma
forma não de simples reprodução mas de criação ou, pelo menos, de recriação do
conhecimento. A familiaridade com os procedimentos de investigação e com o
processo histórico de produção e disseminação de conhecimentos apresenta grande
relevância. No Curso, a pesquisa se constitui em instrumento de ensino e em
conteúdo de aprendizagem na formação: para que a atitude de investigação e a
relação de autonomia se concretizem, o bacharel em ciência da computação
necessita conhecer e saber usar os procedimentos de investigação científica. Tal
atividade é proporcionada aos acadêmicos através dos Programas de Iniciação
Científica e Programa de Estágios Voluntários (atividades complementares) e do
Trabalho de Conclusão de Curso (atividade obrigatória).

\section{Integração com pesquisa e extensão}

O Projeto Pedagógico do Bacharelado em  Ciência da Computação prevê a inserção
dos alunos em projetos de pesquisa e extensão  universitária na região onde se
encontra o curso. Além disso, prevê a modernização  do processo empreendedor em
Tecnologias da Informação e Comunicação (TIC) da  região no qual está inserido.

Para acelerar este processo nas suas diversas Unidades  Acadêmicas, a UFRPE
deverá prover mecanismos de intercâmbio entre seus Campi, possibilitando que
tanto alunos quanto docentes participem de editais de  fomento e bolsas de apoio
tecnológico e apoio social. Também deverá modernizar  seu regimento e
legislações internas para possibilitar que docentes e  discentes das várias
unidades acadêmicas da UFRPE possam usufruir dos mesmos direitos e deveres.

\section{Ensino problematizado e contextualizado}

Entende-se que o sucesso do processo ensino-aprendizagem está
relacionado diretamente à capacidade de colocar de forma ampla o problema a ser
resolvido e contextualizá-lo no âmbito do curso como um todo, assegurando, para
garantir tal objetivo, a indissociabilidade entre ensino, pesquisa e extensão. 

A articulação entre ensino, pesquisa e extensão é fundamental no processo de
produção do conhecimento, pois permite estabelecer um diálogo entre a Ciência da
Computação e as demais áreas, relacionando o conhecimento científico à realidade
social. 

Além das atividades contempladas nas disciplinas que proporcionam a
problematização e contextualização do ensino, entendendo ser o docente um agente
indispensável na execução desta atividade, o Trabalho de Conclusão de Curso, o
Estágio Curricular Supervisionado e as Atividades Complementares focarão,
prioritariamente, a interdisciplinaridade e contextualização do ensino. 

\section{Flexibilidade curricular}

O ensino de graduação, voltado para a construção do
conhecimento, não pode pautar-se por uma estrutura curricular rígida, baseada
num enfoque unicamente disciplinar e sequenciada a partir de uma hierarquização
artificial dos conteúdos, quando a realidade se apresenta em uma multiplicidade
interdependente e a dinâmica de transformação desta coloca a necessidade de um
aprender permanente. 

Desta forma, a flexibilidade desponta como elemento
indispensável à estruturação curricular de modo a atender tanto às demandas da
sociedade moderna quanto àquelas que se direcionam a uma dimensão criativa e
libertária para a existência humana, constituindo-se não apenas em
possibilidade, mas em condição necessária à efetivação de uma formação
profissional de qualidade. 

No curso, a flexibilidade curricular será garantida
através de uma série de ações tomadas visando oportunizar os acadêmicos
vivenciar oportunidades no âmbito da Universidade, tais como: redução da carga
horária total do Curso (tempo livre como importante componente curricular,
oportunizando a ocorrência de momentos formativos inovadores e profícuos),
oferecimento de disciplinas optativas nas diferentes ênfases do Curso
(oportunidade de escolha por parte do acadêmico, respeitando suas competências e
habilidades), atividades complementares (flexíveis e diversas, com carga horária
mínima estabelecida, mas definidas pelos acadêmicos, conforme seus anseios). 

\section{Integração Universidade-Empresa}

O desafio de formar um Bacharel em Ciência  da Computação preparado para
enfrentar o mercado de trabalho altamente  competitivo que exige profissionais
altamente qualificados. O próprio conceito  de qualificação profissional vem se
alterando, com a presença cada vez maior de  componentes associadas às
capacidades de coordenar informações, interagir  com pessoas, interpretar de
maneira dinâmica a realidade.

\section{Estímulo à capacidade de trabalho autônomo}

Tendo consciência do diferencial na formação  profissional relacionado à
capacidade de desenvolver atividades de forma  autônoma, o curso visa estimular,
ao longo de toda a sua duração, a capacidade de  trabalho de forma autônoma,
onde o estudante se converte em protagonista de  sua própria aprendizagem e
desenvolve sua capacidade de ``aprender a aprender''.  A realidade mostra que
este diferencial pode ser conseguido com treinamento, de forma que o processo
ensino-aprendizagem contribuirá, perpassando todas as disciplinas, para o
alcance desta capacidade ao final do curso.

\section{Desenvolvimento de habilidades para o trabalho em equipe}

Da mesma forma que o  explicitado no item anterior, tem-se a consciência para o
diferencial do profissional com  habilidade de trabalho em equipe. Busca-se,
desta forma, ao longo de todo o  Curso, promover atividades que promovam a
possibilidade de desenvolver trabalhos  em equipe, inclusive nas formas de
avaliação das disciplinas.

%%%%%%%%%%%%%%%%%%%%%%%%%%%%%%%%%%%%%%%%%%%%%%%%%%%%%%%%%%%%%%%%%%%%%%%%%%%%%%%%%%%%%
%
% chapter Estágio Supervisionado Obrigatório
%
%%%%%%%%%%%%%%%%%%%%%%%%%%%%%%%%%%%%%%%%%%%%%%%%%%%%%%%%%%%%%%%%%%%%%%%%%%%%%%%%%%%%%%

\section{Estágio Supervisionado Obrigatório}
\label{cap_eso}

De acordo com a Lei nº 11.788/2008, o estágio é um ``ato educativo escolar supervisionado, desenvolvido no ambiente de trabalho'' que tem o propósito de garantir o ``aprendizado de competências próprias da atividade profissional e a contextualização curricular, objetivando o desenvolvimento do educando para a vida cidadã e para o trabalho''. O Estágio Supervisionado Obrigatório (ESO), fazendo parte da matriz curricular, constitui-se num espaço de aprendizagem concreta de vivência prática. O objetivo central se direciona na aplicação dos conhecimentos científicos adquiridos durante a realização do curso e a vivência profissional. Este projeto pedagógico se fundamenta também na Resolução UFRPE/CONSEPE nº 678/2008\footnote{Resolução UFRPE/CONSEPE nº 678/2008: \url{http://www.ufape.edu.br/estagios}}, especificamente, sobre o Estágio Supervisionado Obrigatório.

Os mecanismos de acompanhamento e de cumprimento serão estabelecidos e acompanhados pelo professor da disciplina de estágio supervisionado obrigatório e em consenso com a Comissão de Estágio Supervisionado (COE). Esta comissão tem o objetivo de planejar, coordenar, validar e avaliar o estágio dos acadêmicos do curso. As atividades de estágio devem estar previstas nas seguintes áreas de conhecimento:

\begin{itemize}
    \item Gestão da Tecnologia da Informação;
    \item Análise de Desenvolvimento de Sistemas;
    \item Banco de Dados;
    \item  Computação Gráfica;
    \item  Engenharia de Software;
    \item  Engenharia da Computação;
    \item  Redes de Computadores e Sistemas Distribuídos;
    \item  Outras áreas aprovadas pela COE.
\end{itemize}

As etapas necessárias para iniciar o processo de ESO envolvem outras sub etapas, estas estão bem descritas e apresentadas no site institucional do setor de estágio e de forma resumida, abaixo:

\begin{enumerate}
    \item Matrícula na disciplina de ESO;
    \item Solicitação do seguro junto a esta IES;
    \item Entrega do Termo de Compromisso;
    \item Realização do ESO;
    \item Escrita do Relatório Técnico do ESO;
    \item Defesa do Relatório Técnico do ESO;
    \item Avaliação do Relatório Técnico do ESO;
    \item Entrega do Relatório Técnico do ESO.
\end{enumerate}

As explanações para as etapas apresentadas estão a seguir:

\begin{enumerate}
    \item \textbf{Matrícula na disciplina de ESO} – Somente poderão se matricular na disciplina de ESO os alunos que foram aprovados nas disciplinas de Banco de Dados, Redes de Computadores e Engenharia de Software. Esta decisão leva em conta que os estagiários devem estar com aproximadamente mais da metade da carga total do curso concretizada e já possuem um grau de conhecimento adequado para estagiar na área, tornando assim o estágio melhor desenvolvido e mais bem aproveitado para um futuro vínculo empregatício.
    \item \textbf{Solicitação do seguro junto a esta IES} – Segundo a Lei nº 11.788 a contratação de seguro de vida contra acidentes pessoais em favor do estagiário é obrigatória. Como o estágio é obrigatório para obtenção do diploma no curso de BCC da UFAPE, o seguro fica a cargo dessa instituição de ensino.
    \item \textbf{Entrega do Termo de Compromisso} – O Termo de Compromisso de Estágio é um acordo tripartite celebrado entre o educando, a parte concedente do estágio e a instituição de ensino, prevendo as condições de adequação do estágio à proposta pedagógica do  curso, à etapa e modalidade da formação escolar  do  estudante  e ao horário e calendário escolar. O termo deve ser entregue impresso em três vias, assinadas e carimbadas pela parte concedente, pelo  Supervisor,  pelo  Orientador, pelo Estagiário e pela Instituição.
    \item \textbf{Realização do ESO} – Em hipótese alguma o estágio pode ser iniciado sem a concretização das etapas 1, 2 e 3 apresentadas anteriormente. O estágio deve ser realizado sob supervisão de alguém formado na área de TI ou que possua no mínimo dois anos de experiência na área (comprovada via diploma ou declaração) e orientado por algum professor da UFAPE (dando preferência aos professores do curso de BCC/UFAPE). O estagiário não deverá ultrapassar 06 (seis) horas diárias e 30 (trinta) horas semanais para as atividades do estágio, assim sendo, é preciso estipular 10 semanas para concretização das 300 horas necessárias para o ESO em BCC, considerando que este tempo deve estar dentro do prazo para finalização do período corrente e da data limite para a defesa do relatório (dadas no calendário para cada semestre letivo).
    \item \textbf{Escrita do Relatório Técnico do ESO} – Após a realização das atividades do estágio e integralização da carga horária total o estagiário deve escrever o relatório técnico do estágio, apresentando as atividades realizadas, seguindo o modelo disponibilizado pelo curso. Somente após as correções sugeridas pelo professor orientador e o aval do mesmo para defesa, o estagiário deverá entregar o relatório ao professor da disciplina de ESO, que por sua vez, repassa para a COE, responsável por marcar a data de defesa do relatório e sugerir as melhorias.
    \item \textbf{Defesa do Relatório Técnico do ESO} – O estagiário deverá realizar uma apresentação oral do Relatório de Estágio para o professor presidente da COE. A defesa visa avaliação e composição da nota final de ESO.
    \item \textbf{Avaliação do Relatório Técnico do ESO} – A nota final do ESO (média na disciplina) será composta da seguinte forma: 25\% da nota é baseada na avaliação realizada pelo supervisor do estagiário na empresa, através de preenchimento de formulário padrão encaminhado pela Coordenação do Curso; 75\% da nota é baseada na média das avaliações realizadas pelo presidente da COE (conforme o etapa 6) e pelo professor-orientador do estagiário.
    \item \textbf{Entrega do Relatório Técnico do ESO} – O acadêmico deverá depositar, após a correção final do relatório (sugeridas pelo presidente da COE após defesa), junto à Biblioteca, arquivo com o Relatório do ESO, em sua versão final, de acordo com as normas da ABNT, em formato PDF aberto, conforme Manual de Normalização, conforme orientações disponibilizadas pela Biblioteca Ariano Suassuna.
\end{enumerate}

Uma outra alternativa para o atendimento do componente curricular Estágio Supervisionado Obrigatório é a Equiparação de Atividades ao ESO, procedimento que permite o aproveitamento de atividades realizadas em projetos de ensino, pesquisa e extensão. A Resolução UFRPE/CEPE n.º 425/2010 regulamenta a previsão nos Projetos Pedagógicos de curso da equiparação das atividades de Extensão, Monitoria e Iniciação Científica como estágios curriculares. Este procedimento está descrito na subseção seguinte.

\subsection{Da Equiparação de Atividades com o ESO}

Com base na Lei n.º 11.788/2008, que dispõe o estágio de estudantes e na Resolução UFRPE/CEPE n.º 678/2008, que estabelece normas para organização e regulamentação do Estágio Supervisionado Obrigatório para os estudantes dos Cursos de Graduação da UFRPE e dá outras providências; e, na Resolução UFRPE/CEPE n.º 425/2010, que regulamenta a equiparação ao Estágio Supervisionado Obrigatório, das atividades de Extensão, Monitoria e Iniciação Científica dos Cursos de Graduação da UFRPE, devem ser compatíveis com a formação acadêmica do estudante requerente, comprovadas a carga horária de 300 (trezentas) horas, e seguidas as orientações\footnote{\href{http://www.ufape.edu.br/estagios}{http://www.ufape.edu.br/estagios}} da Coordenadoria de Estágios, fica estabelecido que:

\begin{itemize}
    \item O discente deve manifestar interesse em equiparar suas atividades, exclusivamente, de Extensão, Monitoria ou Iniciação Científica ao Estágio Supervisionado Obrigatório, e deve requerer a equiparação à Coordenação do Curso, através da abertura de processo administrativo, no início do semestre anterior ao da conclusão do curso, conforme estabelecido na Resolução UFRPE/CEPE n.º 425/2010. O processo administrativo deve conter os documentos:

    \begin{enumerate}
        \item Modelo de Requerimento de Equiparação;
	\item Modelo de Declaração de autorização de Equiparação-Orientador;
	\item Relatório das Atividades do Projeto;
	\item Certificado ou declaração de realização da atividade a ser equiparada, constando a carga horária da mesma;
	\item Relatório Final de ESO, conforme orientações do curso e do Manual de Normalização, disponibilizado pela Biblioteca Ariano Suassuna.
    \end{enumerate}

    É possível que o processo seja aberto contendo apenas os quatro primeiros documentos, enquanto o relatório final segue em construção sob as orientações do Professor da Disciplina e do Professor Orientador.
    
    \item O processo administrativo de equiparação das atividades de Extensão, Monitoria ou Iniciação Científica será encaminhado para relatoria, e uma vez aprovado em CCD, o discente será submetido aos procedimentos de avaliação na disciplina de ESO.
    \item Os casos omissos a essas orientações serão submetidos à apreciação do CCD do curso, Pró-Reitoria de Ensino de Graduação e Coordenadoria de Estágios.
\end{itemize}

\subsection{Do Estágio Não Obrigatório}

O Estágio Não Obrigatório é uma atividade que tem como objetivo proporcionar ao aluno a oportunidade de aplicar seus conhecimentos acadêmicos em situações de prática profissional, e pode ser realizado a partir do momento que o aluno é aprovado na disciplina de Algoritmos e Estrutura de Dados, disciplina esta considerada como um requisito mínimo para aplicação de algumas práticas da computação. Como complemento, o Estágio Não Obrigatório não pode ultrapassar 120 (cento e vinte) horas para aproveitamento como Atividades Curriculares Complementares.

\section{Atividades Complementares}

Objetivando atingir o perfil profissional definido e exigido pelo mercado e
também pela sociedade, a Matriz Curricular do Curso prevê a realização de
atividades complementares, que deverão ser realizadas ao longo do curso. 

A ampliação do horizonte da formação profissional. Atividades como iniciação
científica e tecnológica, programas acadêmicos amplos, programas de extensão
universitária, visitas técnicas, participação e apresentação de trabalhos em
eventos científicos, organização de eventos, estágios extra-curriculares,
participação em seminários e palestras, realização de cursos.

As atividades complementares,  nos termos explicitados pelas respectivas
Resoluções do CNE, são disciplinadas  e sistematizadas pela Resolução n.
313/2003 do CEPE/UFRPE. O aluno deverá cursar  \textbf{obrigatoriamente um
mínimo de 320 horas em Atividades Complementares Curriculares e um mínimo de 420 horas em
disciplinas optativas}. Tais restrições visam que o aluno se integre  às atividades da UFRPE e comunidade e
que também curse um mínimo de disciplinas  que compõem Áreas de Saber
constituídas por disciplinas optativas,  respectivamente. As atividades
complementares estão inseridas no eixo de  formação livre, cujos créditos podem
ser obtidos em quaisquer atividades acadêmicas  curriculares. Além disso, a
solicitação da creditação das atividades curriculares  complementares deverá ser
feita pelo aluno por meio de requerimento documentado e  encaminhado ao
Colegiado para proceder conforme Art. 37 da referida  Resolução do CEPE:

Deferido o aproveitamento pelas instâncias competentes,  o Coordenado de Curso
remeterá ao DRCA, para creditar no histórico escolar do  Aluno a carga horária e
Créditos, correspondente ao aprovado, considerando a Tabela~\ref{tab:atividadescomplementares} reproduzida a seguir.

% Table generated by Excel2LaTeX from sheet 'Plan1'
\newpage
    \begin{longtable}{r|r|l|r|r}
      \caption{Atividades Complementares}
      \label{tab:atividadescomplementares}%
      \\
\hline
    \textbf{No.} & \textbf{Código} & \textbf{Descrição} & \textbf{Horas/Aula} & \textbf{Créditos} \\
    \hline
%    \midrule
    1     & 14001 & Monitoria I & 60    & 4 \\     \hline
    2     & 14002 & Monitoria II & 60    & 4 \\     \hline
    3     & 14003 & Programa Especial de Treinamento I & 60    & 4 \\     \hline
    4     & 14004 & Programa Especial de Treinamento II & 60    & 4 \\     \hline
    5     & 14005 & Projeto de Pesquisa I & 60    & 4 \\     \hline
    6     & 14006 & Projeto de Pesquisa II & 60    & 4 \\     \hline
    7     & 14007 & Vivência Profissional Complementar I & 60    & 4 \\    \hline
    8     & 14008 & Vivência Profissional Complementar II & 60    & 4 \\    \hline
    9     & 14009 & Programa de Extensão I & 60    & 4 \\    \hline
    10    & 14010 & Programa de Extensão II & 60    & 4 \\    \hline
    11    & 14011 & Programa de Alfabetização I & 60    & 4 \\    \hline
    12    & 14012 & Programa de Alfabetização II & 60    & 4 \\    \hline
    13    & 14013 & Projeto de Extensão I & 60    & 4 \\    \hline
    14    & 14014 & Projeto de Extensão II & 60    & 4 \\    \hline
    15    & 14015 & Discussões Temáticas I & 15    & 1 \\    \hline
    16    & 14016 & Discussões Temáticas II & 15    & 1 \\    \hline
    17    & 14017 & Tópicos Especiais I & 15    & 1 \\    \hline
    18    & 14018 & Tópicos Especiais II & 30    & 2 \\    \hline
    19    & 14019 & Prática Integrada I & 15    & 1 \\    \hline
    20    & 14020 & Prática Integrada II & 30    & 2 \\    \hline
    21    & 14021 & Cursos de Extensão I & 30    & 2 \\    \hline
    22    & 14022 & Cursos de Extensão II & 60    & 4 \\    \hline
    23    & 14023 & Evento de Extensão I & 30    & 2 \\    \hline
    24    & 14024 & Evento de Extensão II & 60    & 4 \\    \hline
    25    & 14025 & Publicação Técnico-Científica I & 30    & 2 \\    \hline
    26    & 14026 & Publicação Técnico-Científica I & 60    & 4 \\    \hline
    27    & 14027 & Produto de Extensão I & 30    & 2 \\    \hline
    28    & 14028 & Produto de Extensão II & 60    & 4 \\    \hline
    29    & 14029 & Prestação de Serviço I & 30    & 2 \\    \hline
    30    & 14030 & Prestação de Serviço II & 60    & 4 \\    \hline
    \end{longtable}%



%%%%%%%%%%%%%%%%%%%%%%%%%%%%%%%%%%%%%%%%%%%%%%%%%%%%%%%%%%%%%%%%%%%%%%%%%%%%%%%%%%%%%
%
% section Curricularização da Extensão
%
%%%%%%%%%%%%%%%%%%%%%%%%%%%%%%%%%%%%%%%%%%%%%%%%%%%%%%%%%%%%%%%%%%%%%%%%%%%%%%%%%%%%%%

\section{Curricularização de Atividades de Extensão}
\label{chap:curricularizacao-da-extensao}

As atividades de extensão do curso de BCC, seguem a Resolução UFAPE/CONSEPE nº 7/2022, de 22/08/2022, que dispõe sobre a Integralização das Atividades de Extensão como componente curricular dos Cursos de Graduação da UFAPE.

As atividades de extensão devem estar presentes no processo formativo do estudante do curso de BCC e devem compor, no mínimo, 10\% (dez por cento) da carga horária total do curso (320 horas), que devem ser desenvolvidas em atividades de extensão, ao longo do vínculo com o curso.

Conforme preconiza o Art.~3º, da supracitada Resolução, as Atividades Curriculares de Extensão (ACEX) são as atividades institucionalizadas sob as modalidades de Extensão, Projeto ou Programa, as quais objetivam o protagonismo do estudante na práxis voltada ao desenvolvimento de suas habilidades, como parte da sua formação humana, ética e cidadã e à solução de problemas da sociedade. O Art.~4º da supracitada Resolução, estabelece que as ACEX, em suas diferentes modalidades, devem:

\begin{enumerate}[label=\Roman*.]
   \item Atender à Resolução que dispõe da Política de Extensão vigente da UFAPE;
   \item Estar previstas nos Projetos Pedagógicos dos cursos (PPC) de graduação da UFAPE;
   \item Ser diferenciada de ACC, para não haver duplicidade no cômputo de carga horária.
\end{enumerate}
  
As ACEX poderão ser integralizadas por meio de:

\begin{itemize}
   \item Projetos a serem desenvolvidos de forma independente ou vinculado a um Programa.
   \item Atividades de Extensão institucionais certificadas e registradas na Pró-Reitoria de Extensão e Cultura da UFAPE, não sendo cumulativas e não vinculadas como ACC, devendo ser diferenciada, para não ocorrer o cômputo da carga horária em duplicidade.
   \item Projetos ou Programas prevêem a realização de Cursos, Eventos e Prestação de Serviços.
   \item Projetos ou Programas em andamento, nos quais poderão ser vinculados Cursos, Eventos ou Prestação de Serviço posteriormente à institucionalização na UFAPE, desde que seja informada no relatório parcial ou final.
\end{itemize}

O pedido de aproveitamento, juntamente com a documentação comprobatória (comprovação de cadastro junto à Pró-Reitoria de Extensão e Cultura, da Universidade Federal do Agreste de Pernambuco, e declaração/certificado de participação) que se refere o Art. 4° deverá ser encaminhado à Coordenação de Curso e avaliado pelo Colegiado de Coordenação Didática do Curso de Bacharelado em Ciência da Computação, da Universidade Federal do Agreste de Pernambuco.

Os casos omissos deverão ser apresentados e discutidos no Núcleo Docente Estruturante para posterior encaminhamento e deliberação do Colegiado de Coordenação Didática do Curso de Bacharelado em Ciência da Computação, da Universidade Federal do Agreste de Pernambuco.



%%%%%%%%%%%%%%%%%%%%%%%%%%%%%%%%%%%%%%%%%%%%%%%%%%%%%%%%%%%%%%%%%%%%%%%%%%%%%%%%%%%%%
%
% chapter Perfil do egresso
%
%%%%%%%%%%%%%%%%%%%%%%%%%%%%%%%%%%%%%%%%%%%%%%%%%%%%%%%%%%%%%%%%%%%%%%%%%%%%%%%%%%%%%%



\chapter{Perfis}

\section{Estudante ingressante}

Deve se interessar pela computação e, em particular, pela ciência. Deve possuir
entusiasmo para conhecer e dominar novos assuntos, além de disposição  para
construir sua própria reputação por meio dos produtos do seu esforço  próprio ou
resultantes de trabalho em equipe do qual participa.  Deve possuir atitude e a
necessidade de realizar, mesmo sem supervisão. De engajar-se em representações
locais, regionais, nacionais e internacionais, através de representações de
classe, visando a atualização e fortalecimento da sociedade. Estes atributos são
esperados na conduta do estudante ingressante e utilizados ao longo do curso.

\subsection{Descrição dos requisitos psicofísicos}
 
Para atender ao perfil profissional definido, as atividades do curso priorizam o
exercício dos requisitos inerentes ao desempenho da profissão, a citar:

\begin{itemize}
  \item método e disciplina de trabalho;
  \item raciocínio lógico e abstrato; 
  \item capacidade de trabalho em equipe; 
  \item criatividade, produtividade e iniciativa; 
  \item disposição para efetuar trabalho complexo e minucioso; 
  \item compromisso com o desenvolvimento tecnológico; 
  \item compromisso com o ser humano; 
  \item senso crítico, seriedade e responsabilidade. 
\end{itemize}


\section{Egresso}

Do egresso de um curso de Bacharelado em Ciência da Computação é exigida uma
predisposição e aptidões para a área, além de um conjunto de competências, 
habilidades e atitudes a serem adquiridas durante a realização do curso.

\subsection{Definição do perfil profissional}

Por definição, o Bacharel em Ciência da Computação deve ser um profissional 
qualificado para a pesquisa e desenvolvimento na área de computação, para  o
projeto e construção de hardware e software básico e também para o uso de 
sistemas computadorizados em outras áreas da atividade humana, a fim de
viabilizar  ou aumentar a produtividade e a qualidade de todos os tipos de
procedimentos.  Na UFRPE todo egresso deve ser um profissional(1)  com domínio e
capacidade para trabalhar na área da Computação, desenvolvendo projetos de
computadores  e sistemas de computação, programas e sistemas de informação; (2)
atento  ao caráter ecológico, social e ético; e (3) que exerça suas atividades
na  sociedade com responsabilidade.

\subsection{Problemas que os egressos deverão estar aptos a resolver}

Dado o perfil profissional desejado (seção 1), o egresso deverá estar apto a resolver as seguintes classes de problemas: 

\begin{itemize}
  \item concepção, especificação, projeto, construção, avaliação e adaptação de sistemas digitais; 
  \item análise e projeto de estrutura lógica e funcional de computadores e sua
implementação;
  \item desenvolvimento e implementação de software básico e de apoio para
  sistemas computacionais;
  \item projeto e desenvolvimento de sistemas e programas usando linguagens de
  programação;
  \item projeto e desenvolvimento de sistemas de estruturação de informação; 
  \item projeto e desenvolvimento de redes de processamento local e remota, em
  matéria de hardware e de software;
\end{itemize}

O egresso do curso de Bacharelado em Ciência da Computação deve estar
preparado para propor soluções inovadoras e  adequadas para problemas propostos,
capacitado a acompanhar e avaliar avanços  tecnológicos em computação, bem como
aplicar e implementar as evoluções, reposições e adaptações que se façam
necessárias, tanto de forma reativa com pró-ativa, logo deve estar apto a
desenvolver as seguintes funções no mercado de trabalho:

\begin{itemize}
  \item \textbf{empreendedor} --- descobrimento e empreendimento de novas
  oportunidades para aplicações usando sistemas computacionais e avaliando a 
  conveniência de se investir no desenvolvimento da aplicação;
  \item \textbf{consultor} --- consultoria e assessoria a empresas de diversas
  áreas no que tange ao uso adequado de sistemas computacionais;
  \item \textbf{coordenador de equipe} --- coordenação de equipes envolvidas em
  projetos na área de computação e informática;
  \item \textbf{membro de equipe} --- participação de forma colaborativa e 
  integrada de equipes que desenvolvem projetos na área de informática;
  \item \textbf{pesquisador} --- participação em projetos de pesquisa científica
  e tecnológica.
\end{itemize}


%%%%%%%%%%%%%%%%%%%%%%%%%%%%%%%%%%%%%%%%%%%%%%%%%%%%%%%%%%%%%%%%%%%%%%%%%%%%%%%%%%%%%
%
% chapter Caracterização e Dados do Curso
%
%%%%%%%%%%%%%%%%%%%%%%%%%%%%%%%%%%%%%%%%%%%%%%%%%%%%%%%%%%%%%%%%%%%%%%%%%%%%%%%%%%%%%%


\chapter{Caracterização e Dados do Curso}

A Tabela~\ref{ltab:DadosCaracteristicas1} apresenta os parâmetros de caracterização de oferta do curso de
Bacharelado em Ciência da Computação para a Unidade Acadêmica de Garanhuns da UFRPE.

\begin{small}
\begin{longtable}{ll}
 \caption{Dados e características do curso}
 \label{ltab:DadosCaracteristicas1}
\\
\hline
Curso & Ciência da Computação \\ 
\hline
Modalidade & Bacharelado \\ 
Título do egresso & Bacharel em Ciência da Computação \\ 
Área do conhecimento & Ciências Exatas e da Terra \\ 
Modalidade de Educação & Presencial (ensino a distância restrito ao percentual imposto \\ 
 & por legislação em vigor). \\ 
Habilitação & Ciência da Computação \\ 
Local de Oferta & Universidade Federal Rural de Pernambuco \\ 
 & Unidade Acadêmica de Garanhuns \\ 
Número de vagas & 80 (oitenta) vagas anuais, sendo 40 (quarenta) vagas semestrais \\ 
Carga horária & 3.200 (três mil e duzentas) horas \\ 
Horário de funcionamento & Predominantemente noturno \\ 
Forma de acesso & Processo seletivo (ENEM). Em caso de existência de vagas\\
 & é possível o ingresso por: (a) transferência de outras IES;\\
 & (b) portadores de diploma; (c) reingresso; ou, (d) reopção,\\
 & de acordo com  processo seletivo específico e resoluções da UFRPE. \\
\hline
\end{longtable}
\end{small}



%%%%%%%%%%%%%%%%%%%%%%%%%%%%%%%%%%%%%%%%%%%%%%%%%%%%%%%%%%%%%%%%%%%%%%%%%%%%%%%%%%%%%
%
% chapter Organização Curricular
%
%%%%%%%%%%%%%%%%%%%%%%%%%%%%%%%%%%%%%%%%%%%%%%%%%%%%%%%%%%%%%%%%%%%%%%%%%%%%%%%%%%%%%%

\chapter{Organização do Curso (e Recursos Humanos)}

A estrutura curricular proposta, (visualizar a Matriz Curricular no apêndice 1) possui uma carga-horária mínima de 3.200 horas para a integralização dos créditos exigidos para a obtenção do título de Bacharel em Ciência da Computação.

Como 20\% da carga horária total do curso de Graduação  (máximo de 640 horas para este curso)  pode ocorrer à distância (MEC, Portaria No. 4.059, de 10/12/2004), através de  disciplinas que permitam desenvolver suas atividades didáticas, módulos ou unidades de ensino-aprendizagem centrados na auto-aprendizagem e com mediação de recursos didáticos organizados em ambientes de informação que utilizem tecnologia de comunicação remota, avaliações presenciais, incluídos métodos e práticas de ensino-aprendizagem que incorporem o uso de tecnologias de informação e comunicação para a realização dos objetivos pedagógicos, bem como prever encontros presenciais e atividades de tutoria, através de docentes qualificados, e que poderão ser encorajadas pela Instituição a serem ofertadas à distância. Assim, pode-se compartilhar a oferta das disciplinas nos diversos \textit{Campi} onde o curso será ofertado oferecendo ao aluno uma maior variedade na opção de disciplinas e maior integração. 

Além disso, estar-se-á auxiliando os recém-criados cursos a desenvolverem seus grupos de pesquisa e extensão.

De acordo com as diretrizes do MEC, os currículos dos cursos da área de Computação e Informática podem ser compostos por quatro grandes áreas de formação:

\begin{itemize}
  \item Formação básica;
  \item Formação tecnológica;
  \item Formação complementar;
  \item Formação humanística.
\end{itemize}

A formação básica compreende os princípios básicos da Ciência da Computação, além da Matemática necessária para defini-los formalmente, a Física e Eletricidade necessárias para permitir o entendimento e o projeto de computadores viáveis e a formação pedagógica, que introduz os conhecimentos básicos da construção do conhecimento. 

A formação tecnológica, também chamada de aplicada ou profissional, aplica os conhecimentos básicos adquiridos no desenvolvimento tecnológico da Computação. Os objetivos são criar instrumentos (ferramentas) de interesse da sociedade ou robustecer tecnologicamente os sistemas de computação, para permitir a construção de ferramentas antes inviáveis ou ineficientes. Entre temas de disciplinas de formação tecnológica, pode-se citar: Sistemas Operacionais, Redes de Computadores, Banco de Dados, Sistemas Multimídia, Interface Humano-Computador, Realidade Virtual, Inteligência Artificial, Computação Gráfica e Processamento de Imagens.

A formação complementar permite uma interação dos egressos dos cursos com outras profissões, tendo como objetivo dotar o aluno do conhecimento necessário para resolver problemas de outras áreas. Para que possa interagir com profissionais de outras áreas na busca de soluções computacionais complexas para seus problemas, profissional de Computação deve conhecer de forma geral e abrangente essas áreas. Como exemplos de áreas de formação complementar têm-se: Educação, Economia, Administração, Biologia, Música, entre outras. A formação humanística dá ao egresso uma dimensão social e humana, por meio do tratamento de temas como o empreendedorismo, a ética em computação, a sociologia e a filosofia. 

O Curso de BCC da UAG foi dividido de acordo com a Tabela~\ref{tab:CHNucleoFormacao}. Os percentuais das áreas de formação são calculados tendo como referência o total parcial de carga horária que não leva em conta nem as atividades complementares, excluindo as optativas, nem o Trabalho de Conclusão de Curso (TCC). As disciplinas optativas por fazerem parte de eixos de conhecimento do curso estão contabilizadas dentro da formação tecnológica.

% Table generated by Excel2LaTeX from sheet 'Plan3'
\begin{table}[htbp]
  \centering
  \caption{Distribuição de carga horária por núcleos de formação}
    \begin{tabular}{l|r|r|r}
\hline \textbf{Área de Formação} & \textbf{Créditos} & \textbf{Carga Horária} & \textbf{\%} \\ \hline
    CH Básico & 80    & 1200  & 41,38\% \\ \hline
    CH Profissional & 70    & 1050  & 36,21\% \\ \hline
    CH Optativas & 28    & 420   & 17,50\% \\ \hline
    CH Complementares &       & 320   & 13,33\% \\ \hline
    CH Estágio Supervisionado &      & 300   & 12,50\% \\ \hline
    CH Trabalho de Conclusão de Curso &      & 180   & 6,67\% \\ \hline
          &       &       &  \\ \hline
    Matemática & 28    & 420   & 17,50\% \\ \hline
    Ciências Básicas (da Natureza) & 4     & 60    & 2,50\% \\ \hline
    Contexto Social e Profissional & 10    & 150   & 6,25\% \\ \hline
    Fundamentos da Computação & 48    & 720   & 30,00\% \\ \hline
    Tecnologia da Computação & 42    & 630   & 26,25\% \\ \hline
    Optativas & 28    & 420   & 17,50\% \\ \hline
          &       &       &  \\ \hline
    Total Parcial & 160   & 2400  & 100,00\% \\ \hline
    Atividades Complementares &       & 320   &  \\ \hline
    Estágio Supervisionado &       & 300   &  \\ \hline
    Trabalho de Conclusão de Curso &      & 180   &  \\ \hline
    Total Geral &       & 3200  &  \\ \hline


    \end{tabular}%
  \label{tab:CHNucleoFormacao}%
\end{table}%


Para integralização das 3.200 horas da carga-horária mínima exigida, o aluno deverá cumprir esta carga-horária distribuída da seguinte forma: 

\begin{itemize}
  \item 1.980 horas em disciplinas obrigatórias;
  \item 420 horas em disciplinas optativas;
  \item 320 horas em atividades complementares, ver Tabela~\ref{tab:atividadescomplementares}, pág.~\pageref{tab:atividadescomplementares}.
  \item 300 horas de estágio supervisionado.
  \item 180 horas de trabalho de conclusão de curso supervisionado por um Professor da área de Computação;

\end{itemize}

A critério do aluno, outras disciplinas complementares, optativas e eletivas (todas as disciplinas ofertadas na UFRPE são consideradas eletivas para o aluno do Bacharelado em Ciência da Computação, desde que aprovadas pelo Colegiado do Curso) poderão ser incorporadas ao seu histórico escolar. Contudo, as eletivas não são contabilizadas para integralização do currículo de créditos. Os prazos de conclusão de curso estão sumarizados na Tabela~\ref{tab:prazoCurriculo}, a seguir. 

% Table generated by Excel2LaTeX from sheet 'Plan3'
\begin{table}[htbp]
  \centering
  \caption{Prazos para integralização do currículo}
    \begin{tabular}{l|r}
\hline
    \textbf{Prazo} & \textbf{Tempo} \\ \hline 
    Mínimo & 8 semestres \\ \hline
    Pleno & 9 semestres \\ \hline
    Máximo & 15 semestres \\ \hline

    \end{tabular}%
  \label{tab:prazoCurriculo}%
\end{table}%


\section{Infraestrutura para Manutenção do Curso}

Para a adequada manutenção e evolução de um curso de Bacharelado em Ciência da
Computação são necessários 3 (três) itens fundamentais:

\begin{enumerate}
  \item recursos humanos, com docentes qualificados e técnicos operacionais
  (para os serviços de apoio à secretaria, apoio didático e coordenação);
  \item infraestrutura didática, (acesso à \textit{Internet}, laboratórios
  temáticos de informática);
  \item biblioteca, contendo os recursos bibliográficos, revistas impressas e
  digitais;
  \item manutenção do núcleo docente estruturante; 
  \item colegiado de coordenação didática; e,
  \item programa de tutoria institucional.
\end{enumerate} 

Tais recursos são sumarizados nas subseções a seguir.



\subsection{Recursos Humanos}

Para a completa execução do curso de Bacharelado em  Ciência da Computação será
necessária a contratação e docentes para curso, 1 (uma)  secretária, em nível
superior; 2 (dois) técnicos de apoio educacional, em nível médio; e, 2 (dois)
assistentes em administração, em nível médio.

\subsection{Laboratórios e ambientes}

Para a infraestrutura de informática, no que diz respeito aos laboratórios, os
cursos de computação devem oferecer um bom laboratório de software.  É
imprescindível que haja conexão com a \textit{Internet} e que os alunos tenham
acesso no mínimo aos dois ambientes computacionais e de redes mais comuns
atualmente: Windows 7 e UNIX, como uma distribuição Linux, rodando em
\textit{dual-boot}. 

Os laboratórios de hardware devem ser completos, com
instrumental  necessário para  matérias como arquitetura de computadores, 
circuitos digitais e automação: osciloscópios e analisadores digitais,  kits de
programação e simulação de sistemas de automação e de circuitos digitais, além
dos laboratórios temáticos/específicos, abaixo descritos.

\begin{enumerate}
  \item 4 (quatro) salas de aula com quadros com capacidade para 40 (quarenta)
  estudantes;
  \item 10 (dez) laboratórios de ensino com computadores e quadro com capacidade
  para 40 (quarenta) computadores;
  \item 1 (um) anfiteatro para 150 (cento e cinquenta) pessoas;
  \item 1 (uma) sala para servidores (computadores) da rede de informática da
  UAG;
  \item 1 (uma) sala para servidores (computadores) para uso do curso de BCC;
  \item 1 (uma) sala para técnicos de informática da UAG;
  \item 1 (uma) sala para técnicos de informática do curso de BCC;
  \item 1 (um) laboratório de manutenção de computadores;
  \item 1 (uma) sala para almoxarifado/depósito de material;
  \item 2 (duas) salas para monitoria com capacidade para 20 (vinte) estudantes;
  \item 30 (trinta) salas individuais para professores;
  \item 1 (uma) sala para professores visitantes;
  \item 1 (uma) sala para coordenação do curso de BCC;
  \item 1 (um) laboratórios de pesquisa em Computação;
  \item 1 (um) laboratório de Matemática Computacional;
  \item 1 (um) laboratório de Estatística Computacional;
  \item 1 (uma) sala para biblioteca setorial;
  \item 2 (duas) salas sendo uma para sede do Garanhuns Digital e outra para o
  grupo PET computação;
  \item 1 (uma) sala para sede do Garanhuns Digital;
  \item 1 (uma) sala para o centro acadêmico do curso de BCC;
  \item 1 (uma) copa;
  \item 1 (uma) área para empreendedorismo, com 6(seis) salas (empresa júnior e
  incubadora);
  \item 2 (dois) banheiros com chuveiros, feminino e masculino; 
  \item 6 (seis) salas para estudantes de pós-graduação;
  \item 1 (um) sala para secretaria da pós-graduação;
  \item 8 (oito) salas de aula para pós-graduação com capacidade para 20
  (vinte); estudantes;
  \item 2 (duas) salas de computação para pós-graduação com capacidade para 10
  (dez) computadores;
  \item 1 (uma) sede do grupo PET Computação.
\end{enumerate}



\subsection{Biblioteca}
 
Para a biblioteca, convém adquirir, pelo menos, 13 (treze) exemplares de
cada livro-texto (bibliografia básica) de cada  ementa de disciplina e pelo menos 2 (dois) exemplares de
(bibliografia complementar) por ementa de disciplina.

Há a necessidade de assinatura de revistas especializadas, em âmbito nacional e internacional, como por exemplo, 
\textit{Communications of th}e ACM, IEEE \textit{Software} e MIS \textit{Quarterly}, e
tecnológicas ou de microinformática, como PC World, Exame Informática, entre outras.

\section{Núcleo Docente Estruturante}


O Núcleo Docente Estruturante (NDE) será estabelecido pelo Colegiado do Curso
segundo a Resolução 065/2011-CEPE, sendo constituído pelo(a) Coordenador(a)
de Curso, como seu presidente nato, e por pelo menos 30\% (trinta por cento) dos
docentes efetivos e atuantes no curso, relacionados pelo Colegiado de Curso que
satisfizerem os requisitos previstos na resolução.

A coordenação do curso possui um coordenador e outros docentes que formam o Núcleo Docente Estruturante, ou NDE. 
Ao NDE cabe a manutenção do presente Projeto Pedagógico de Curso (PPC) e a correspondente implementação. O NDE é 
um órgão consultivo, cujas sugestões e decorrentes ações devem ser avaliadas e aprovadas pelo Conselho Diretor (CD) do 
Instituto de Informática. A definição precisa das atribuições e da constituição do NDE, dentre outras, deverão ser fornecidas em regimento próprio devidamente aprovado pelo Colegiado de Coordenação Didática.


Este grupo deve avaliar constantemente o andamento do Curso, especialmente nos 
primeiros anos, propondo melhorias e ajustes ao PPC e ao funcionamento do Curso,
de forma a possibilitar a realização dos objetivos propostos para o Curso.


\section{Colegiado de Coordenação Didática}

A coordenação didática de cada curso de graduação e de
pós-graduação é exercida por um Colegiado de Coordenação Didática, conforme preconiza o Art. 16, do Estatuto da Unidade Acadêmica de Garanhuns, constituído
pelo Coordenador do Curso, como presidente, pelo seu substituto eventual,  como
vicepresidente, por docentes dos primeiros quatro períodos do curso (quatro
representantes)  e do quinto ao último período do curso (cinco representantes), 
que ministrem disciplinas no curso, por representante(s) do corpo discente de
graduação e de pós-graduação,  escolhidos na forma da legislação vigente, com
mandato de um (1) ano,  permitida uma recondução.

São atribuições dos Colegiados de Coordenação
Didática de Curso, conforme o Art. 6., do Regimento Geral da Unidade Acadêmica de Garanhuns: 
a) elaborar modificações ao currículo do curso,  propondo-as
ao Colegiado Geral de Coordenação Didática; b) propor ao Colegiado Geral de 
Coordenação Didática o elenco de disciplinas optativas do curso; c) Promover, 
através de propostas devidamente, justificadas, ao Colegiado Geral de
Coordenação Didática,  a melhoria contínua do curso; d) propor ao Colegiado
Geral de  Coordenação Didática modificações nos planos dos respectivos cursos; 
e) propor, em cada período letivo, os planos de ensino das disciplinas do
Currículo do Curso;  f) apreciar e deliberar sobre as solicitações acerca do 
aproveitamento de estudos e adaptações, ouvidos os docentes da Unidade com 
competência para julgar e emitir parecer sobre o conteúdo de tais solicitações; 
g) aprovar o Regimento do Centro Acadêmico do Curso, submetendo-o depois à 
homologação do Conselho Universitário; i) exercer as demais funções que lhe são,
explícita ou implicitamente, deferidas em lei, no Estatuto e neste Regimento
Geral;  j) deliberar sobre os casos omissos na esfera de sua competência.

\section{Programa de Tutoria Institucional}

O programa de tutoria foi iniciado na UFRPE-UAG com estudantes do primeiro período, estratificados em quatro grupos distintos orientados por quatro professores do curso. Uma das  características deste programa é manter uma postura de  acompanhamento e apoio direto ao aluno, entretanto, a participação do discente não é obrigatória, bem como seu ingresso em outros programas existentes na faculdade, como iniciação científica ou programas de extensão.

Como resultado do programa espera-se que seja incrementada a qualidade de formação acadêmica dos graduandos, propiciando ao aluno oportunidades de participação em atividades científicas que transcendam a formação básica, além do auxílio em suas dificuldades a partir de estratégias de recuperação para alunos com problemas de aprendizagem. Estas ações visam combater as causas da evasão e da formação deficiente, contribuem para o desenvolvimento e modernização da própria IES.

\subsection{Organização}

Cabe a coordenação do curso o papel de organizar o programa, nomeando os tutores dentre seu corpo docente, para o período de um ano e estipulando a quantidade e relação nominal de alunos que cada tutor orientará.

A coordenação do curso é responsável ainda, por estabelecer e divulgar os critérios básicos do programa de tutoria, comunicar aos docentes e discentes da graduação as ações de recuperação propostas pelos docentes tutores, e finalmente ao final de cada período avaliar o aproveitamento dos alunos.

\subsection{Responsabilidades do Tutor}

O tutor desempenha um papel importante no programa de tutoria. É sua atribuição auxiliar o tutorado em suas dificuldades durante o decorrer do curso, propor uma discussão sobre o plano de atividades dos tutorados sob sua responsabilidade direta, sem, entretanto ultrapassar os limites de autoridade do docente da disciplina e também ao final de cada período relatar as atividades do tutorado.

\subsection{Responsabilidades do Tutorando}

Como tutorandos entenda-se todos os alunos matriculados no curso e que por algum motivo necessitem de auxílio. Este deve participar do plano de atividades proposto pelo tutor, e, acompanhamento do seu desempenho escolar, através do histórico escolar e atividades complementares, incentivados a participara de atividades de extensão, monitoria, iniciação científica.

O tutorando será incentivado, com anuência do titular da disciplina e auxílio de seu tutor, apresentar um seminário aberto à comunidade científica sobre os assuntos específicos que apresenta dificuldade.


\section{Acompanhamento e avaliação do processo de ensino e aprendizagem}

O resultado da avaliação do processo de ensino e aprendizagem do curso é um dos insumos para o planejamento pedagógico. Basicamente consiste nas lições aprendidas no semestre anterior, e que serão acrescidas à base de conhecimento do curso.

Esta base de conhecimento é um dos principais ativos do curso e será mantida semestralmente ou sempre na ocorrência de evento que justifique alteração. É atribuição do Núcleo Docente Estruturante do curso a  avaliação do processo de ensino e aprendizagem.

As lições aprendidas serão obtidas através da avaliação e análise do semestre anterior (fechamento do semestre) e da avaliação 
discente das estratégias adotadas pelo curso. A avaliação discente deve ser feita formalmente (inclusive como forma de 
registro) e por disciplina, de acordo com as normas institucionais. Estas avaliações servirão de insumos para a análise do semestre anterior, da qual participam o corpo docente e discente, sob a responsabilidade da coordenação do curso.

Esta avaliação seguirá a resolução vigente que rege o sistema de avaliação das componentes curriculares na
UFRPE, tendo como referência o perfil do egresso, os objetivos do curso e as competências profissionais orientadoras para a formação do Bacharel em Ciência da Computação. Esta avaliação será complementada pelas seguintes ações:

\begin{itemize}
\item Reuniões semestrais do Coordenador e Vice-Coordenador com os alunos, tentando identificar pontos positivos e negativos no processo ensino e de aprendizagem empregado pelos diversos professores, possivelmente utilizando questionários preenchidos pelos alunos e professores das disciplinas.
\item Utilização dos resultados das avaliações docentes feitas pela UFRPE para identificar oportunidades de melhorias e soluções.
\end{itemize}

Esta análise é o produto do sistema de avaliação do processo de ensino e aprendizagem, que servirá de insumo para o 
planejamento pedagógico do semestre posterior. Tal análise deverá ocorrer ao final de cada semestre.


\section{Acompanhamento e avaliação do Projeto Pedagógico de Curso}


O presente Projeto Pedagógico do Curso de Bacharelado em Ciência da Computação surge de uma demanda da comunidade e busca à formação e qualificação profissional, consequentemente estará em permanente avaliação e sujeito à revisão e aperfeiçoamento.

Serão implantados mecanismos de avaliação docente com periodicidade semestral e de direcionamento do Colegiado do Curso, com periodicidade mínima mensal. Reuniões periódicas do Colegiado de Curso utilizarão os resultados das avaliações docentes para melhorar a alocação docente e treinamento dos profissionais envolvidos com o curso.

Avaliações de periodicidade semestral serão realizadas e reuniões de direcionamento estratatégico anuais para adequação a futuras mudanças e anseios conforme Parecer CNE/CP n. 009/2001.

Pretende-se realizar uma política de avaliação permanente em conjunto com a comunidade e, em especial com os grupos de trabalhos e fóruns de discussão permanentes, capazes de tornar públicas as interpretações do processo avaliativo, construindo parâmetros, critérios e padrões com o coletivo do Colegiado do Curso, região onde o curso se insere e da Universidade para que corresponda às políticas da SBC e às demandas da comunidade.

O Curso de Bacharelado em Ciência da Computação  prevê um sistema de avaliação contínuo de modo que o colegiado tenha subsídios para  efetuar melhorias na qualidade do curso. Este sistema de avaliação será realizado de quatro formas distintas:

\begin{enumerate}
\item Reunião periódica dos professores do colegiado do curso pelo menos uma vez ao semestre para uma avaliação do curso e do semestre.
\item Avaliação das disciplinas do curso, professores, coordenação e direção, por parte dos alunos, de acordo com normas institucionais próprias. As disciplinas obrigatórias poderão ser alteradas, e novas disciplinas optativas poderão ser propostas, a fim de realizar rápidas adequações ao mercado de trabalho, às orientações da Sociedade Brasileira de Computação, Ministério da Educação, ou ainda, às normas internas e institucionais, mediante discussão e aprovação pelo Núcleo Docente Estruturante, e, submetido à Coordenação do Colegiado Didático do curso, para validação.
\item Preenchimento pelo aluno de questionário específico, após a realização de estágio supervisionado obrigatório, indicando como foi a sua inserção e adaptação na empresa e apresentando sugestões de melhorias na atividade de estágio e no curso.
\item Avaliação do Curso e da Instituição pelos egressos com mais de um ano de graduação, através de formulário específico.
\end{enumerate}







%%%%%%%%%%%%%%%%%%%%%%%%%%%%%%%%%%%%%%%%%%%%%%%%%%%%%%%%%%%%%%%%%%%%%%%%%%%%%%%%%%%%%
%
% chapter Atividades complementares
%
%%%%%%%%%%%%%%%%%%%%%%%%%%%%%%%%%%%%%%%%%%%%%%%%%%%%%%%%%%%%%%%%%%%%%%%%%%%%%%%%%%%%%%

\chapter{Atividades Complementares}
As atividades complementares,  nos termos explicitados pelas respectivas
Resoluções do CNE, são disciplinadas  e sistematizadas pela Resolução n.
313/2003 do CEPE/UFRPE. O aluno deverá cursar  \textbf{obrigatoriamente um
mínimo de 200 horas em Atividades Complementares Curriculares e um mínimo de  420 horas em
disciplinas optativas}. Tais restrições visam que o aluno se integre  às atividades da UFRPE e comunidade e
que também curse um mínimo de disciplinas  que compõem Áreas de Saber
constituídas por disciplinas optativas,  respectivamente. As atividades
complementares estão inseridas no eixo de  formação livre, cujos créditos podem
ser obtidos em quaisquer atividades acadêmicas  curriculares. Além disso, a
solicitação da creditação das atividades curriculares  complementares deverá ser
feita pelo aluno por meio de requerimento documentado e  encaminhado ao
Colegiado para proceder conforme Art. 37 da referida  Resolução do CEPE:

Deferido o aproveitamento pelas instâncias competentes,  o Coordenado de Curso
remeterá ao DRCA, para creditar no histórico escolar do  Aluno a carga horária e
Créditos, correspondente ao aprovado, considerando a Tabela~\ref{tab:atividadescomplementares} reproduzida aqui.

% Table generated by Excel2LaTeX from sheet 'Plan1'

    \begin{longtable}{r|r|l|r|r}
      \caption{Atividades Complementares}
      \label{tab:atividadescomplementares}%
      \\
\hline
    \textbf{No.} & \textbf{Código} & \textbf{Descrição} & \textbf{Horas/Aula} & \textbf{Créditos} \\
    \hline
%    \midrule
    1     & 14001 & Monitoria I & 60    & 4 \\
    \hline
    2     & 14002 & Monitoria II & 60    & 4 \\
    \hline
    3     & 14003 & Programa Especial de Treinamento I & 60    & 4 \\
    \hline
    4     & 14004 & Programa Especial de Treinamento II & 60    & 4 \\
    \hline
    5     & 14005 & Projeto de Pesquisa I & 60    & 4 \\
    \hline
    6     & 14006 & Projeto de Pesquisa II & 60    & 4 \\
    \hline
    7     & 14007 & Vivência Profissional Complementar I & 60    & 4 \\
    \hline
    8     & 14008 & Vivência Profissional Complementar II & 60    & 4 \\
    \hline
    9     & 14009 & Programa de Extensão I & 60    & 4 \\
    \hline
    10    & 14010 & Programa de Extensão II & 60    & 4 \\
    \hline
    11    & 14011 & Programa de Alfabetização I & 60    & 4 \\
    \hline
    12    & 14012 & Programa de Alfabetização II & 60    & 4 \\
    \hline
    13    & 14013 & Projeto de Extensão I & 60    & 4 \\
    \hline
    14    & 14014 & Projeto de Extensão II & 60    & 4 \\
    \hline
    15    & 14015 & Discussões Temáticas I & 15    & 1 \\
    \hline
    16    & 14016 & Discussões Temáticas II & 15    & 1 \\
    \hline
    17    & 14017 & Tópicos Especiais I & 15    & 1 \\
    \hline
    18    & 14018 & Tópicos Especiais II & 30    & 2 \\
    \hline
    19    & 14019 & Prática Integrada I & 15    & 1 \\
    \hline
    20    & 14020 & Prática Integrada II & 30    & 2 \\
    \hline
    21    & 14021 & Cursos de Extensão I & 30    & 2 \\
    \hline
    22    & 14022 & Cursos de Extensão II & 60    & 4 \\
    \hline
    23    & 14023 & Evento de Extensão I & 30    & 2 \\
    \hline
    24    & 14024 & Evento de Extensão II & 60    & 4 \\
    \hline
    25    & 14025 & Publicação Técnico-Científica I & 30    & 2 \\
    \hline
    26    & 14026 & Publicação Técnico-Científica I & 60    & 4 \\
    \hline
    27    & 14027 & Produto de Extensão I & 30    & 2 \\
    \hline
    28    & 14028 & Produto de Extensão II & 60    & 4 \\
    \hline
    29    & 14029 & Prestação de Serviço I & 30    & 2 \\
    \hline
    30    & 14030 & Prestação de Serviço II & 60    & 4 \\
    \hline
%    \bottomrule
    \end{longtable}%


% ----------------------------------------------------------
% Finaliza a parte no bookmark do PDF
% para que se inicie o bookmark na raiz
% e adiciona espaço de parte no Sumário
% ----------------------------------------------------------
\phantompart

% ----------------------------------------------------------
% ELEMENTOS PÓS-TEXTUAIS
% ----------------------------------------------------------
\postextual
% ----------------------------------------------------------

% ----------------------------------------------------------
% Referências bibliográficas
% ----------------------------------------------------------
%\bibliography{abntex2-modelo-references}
\chapter*[Referências]{Referências}
\addcontentsline{toc}{chapter}{Referências}

\noindent Sociedade Brasileira de Computação - Currículos de Referência para Cursos de Graduação em Ciência da Computação. Disponível em: http://www.sbc.org.br/index.php?language=1 \&subject=28\&content=downloads\&id=82

\noindent Bacharelado em Sistemas de Informação da Universidade Federal de Minas Gerais. Disponível em: http://www.dcc.ufmg.br/ensino/infbsi 

\noindent Bacharelado em Sistemas de Informação da Universidade Federal de Santa Catarina. Disponível em: http://www.inf.ufsc.br/sin/index.html


\noindent Editora Universitária Champagnat, editor. Anais do II Curso de Qualidade de Cursos de Graduação da área de Computação e Informática, Curitiba - PR, julho 2000. SBC.

\noindent Editora Universitária Champagnat, editor. Anais do III Curso de Qualidade de Cursos de Graduação da área de Computação e Informática, Curitiba - PR, julho 2001. SBC.

\noindent Ministério da Educação.  Oferta de disciplinas não presenciais em cursos presenciais reconhecidos.  MEC, Portaria No. 4.059, de 10/12/2004.

 
\noindent Lei de Diretrizes e Base da Educação Nacional, 1996. Presidente da República.

\noindent Diretoria de Educação da Sociedade Brasileira de Computação. Currículo de Referência da SBC. Disponível em: http://www.sbc.org.br/educacao/

\noindent Conselho de Ensino Pesquisa e Extensão CEPE/UFRPE. Resolução 313/2003, outubro 2003.

\noindent Comissão de Especialistas de Ensino de Computação e Informática do Ministério da Educação (CEEInf/MEC). Diretrizes Curriculares da Área de Computação. 2002.  Disponível em: http://www.inf.ufrgs.br/mec/

\noindent Henry Etzkowitz and Loet Leydesdorff. The dynamics of innovation: from national systems and ``mode 2'' to a triple helix of university-industry-government relations. Research Policy, 29:109-123, 2000. Elsevier Science B.V.

\noindent Henry Etzkowitz, Andrew Webster, Christiane Gebhardt, and Branca Regina Cantisano Terra. The future of the university and the university of the future: evolution of ivory tower to entrepreneurial paradigm. Research Policy, 29:313-330, 2000. Elsevier Science B.V.

\noindent Torsten Leidig. L3- towards an open learning environment. ACM Journal of Educational Resources in Computing, 1(1), Spring 2001. Article 45.

\noindent Darius Mahdjoubi. The linear model of technological innovation, 1997. Disponível em: http://www.gslis.utexas.edu/~darius/lnr\_mdl/lnr\_mdl.html

\noindent Darius Mahdjoubi. Non-linear models of technological innovation, 1997. Disponível em: http://www.gslis.utexas.edu/~darius/non\_mdl/non\_mdl.html.

\noindent IEEE Computer Society and Association for Computing Machinery. Computing curricula 2001: The joint task force on computing curricula.  Disponível em: http://www.computer.org/ education/cc2001/report/ march 2001.
% ----------------------------------------------------------
% Glossário
% ----------------------------------------------------------
%
% Consulte o manual da classe abntex2 para orientações sobre o glossário.
%
%\glossary

% ----------------------------------------------------------
% Apêndices
% ----------------------------------------------------------

% ---
% Inicia os apêndices
% ---
\begin{apendicesenv}

% Imprime uma página indicando o início dos apêndices
\partapendices

% ----------------------------------------------------------
\chapter{Matriz Curricular}
% ----------------------------------------------------------

\begin{landscape}
  \begin{figure}[hbtp]
    \centering
      \caption{\label{fig:matrizcurricular}Matriz Curricular do curso de Bacharelado em Ciência da Computação, UFRPE/UAG}
      \includegraphics[width=1.55\textwidth]{images/matriz6.eps}
  \end{figure}
\end{landscape}

% ----------------------------------------------------------
\chapter{Disciplinas obrigatórias}
% ----------------------------------------------------------
% Table generated by Excel2LaTeX from sheet 'disciplinasObrigatórias'
\begin{longtable}{r|r|l|r|r|r}
     \caption{Perfil Curricular}
     \label{ltab:teste}
     \\
%    \toprule
\hline
%  \rowcolor[rgb]{0.91,0.89,0.88}
    N.    & Período & Disciplinas                       & h/semana & Créditos & CH \\ \hline
\multicolumn{6}{l}{Núcleo: \textbf{Matemática}}\\                                     \hline 
    1     & I       & Cálculo p/ Computação I           & 4     & 4        & 60 \\ \hline
    2     & I       & Geometria analítica               & 4     & 4        & 60 \\ \hline
    3     & I       & Lógica Matemática                 & 4     & 4        & 60 \\ \hline
    4     & II      & Cálculo p/ Computação II          & 4     & 4     & 60 \\     \hline
    5     & II      & Álgebra Linear                    & 4     & 4     & 60 \\     \hline
    6     & III     & Probabilidade e Estatística       & 4     & 4     & 60 \\     \hline
    7     & III     & Matemática Discreta               & 4     & 4     & 60 \\     \hline
          &         & Subtotal (\textbf{Matemática})    & \textbf{28}    & \textbf{28}    & \textbf{420} \\ \hline
\multicolumn{6}{l}{Núcleo: \textbf{Fundamentos da Computação}} \\ \hline
    8     & I     & Introdução à Computação & 2     & 2     & 30 \\ \hline
    9     & I     & Introdução à Programação & 6     & 6     & 90 \\ \hline
    10    & II    & Algoritmos e Estruturas de Dados I & 4     & 4     & 60 \\ \hline
    11    & II    & Programação Orientada à Objetos & 4     & 4     & 60 \\ \hline
    12    & III   & Sistemas Digitais & 4     & 4     & 60 \\ \hline
    13    & III   & Algoritmos e Estruturas de Dados II & 4     & 4     & 60 \\ \hline
    14    & IV    & Arquitetura de Computadores & 4     & 4     & 60 \\ \hline
    15    & IV    & Projeto e Análise de Algoritmos & 4     & 4     & 60 \\ \hline
    16    & IV    & Paradigmas de Linguagens de Programação & 4     & 4     & 60 \\ \hline
    17    & V     & Sistemas Operacionais & 4     & 4     & 60 \\ \hline
    18    & V     & Sistemas de Informação e Tecnologias & 4     & 4     & 60 \\ \hline
    19    & V     & Teoria da Computação & 4     & 4     & 60 \\ \hline
          &       & SubTotal (\textbf{Fundamentos da Computação}) & \textbf{48}    & \textbf{48}    & \textbf{720} \\ \hline
\multicolumn{6}{l}{Núcleo: \textbf{Ciências Básicas}} \\ \hline
    20    & II    & Física p/ Computação & 4     & 4     & 60 \\ \hline
          &       & SubTotal (\textbf{Ciências Básicas}) & \textbf{4}     & \textbf{4}     & \textbf{60} \\      \hline
\multicolumn{6}{l}{Núcleo: \textbf{Tecnologia da Computação}} \\         \hline 
    21    & IV    & Bancos de Dados & 4     & 4     & 60 \\     \hline
    22    & IV    & Engenharia de Software & 4     & 4     & 60 \\     \hline
    23    & V     & Redes de Computadores & 4     & 4     & 60 \\     \hline
    24    & V     & Inteligência Artificial & 4     & 4     & 60 \\ \hline 
    25    & VI    & Reconhecimento de Padrões & 4     & 4     & 60 \\     \hline
    26    & VI    & Compiladores & 4     & 4     & 60 \\     \hline
    27    & VI    & Computação Gráfica & 4     & 4     & 60 \\     \hline
    28    & VI    & Sistemas Distribuídos & 4     & 4     & 60 \\     \hline
    29    & VII   & Projeto de Desenvolvimento & 4     & 4     & 60 \\     \hline
    30    & VII   & Interação Humano-Computador & 4     & 4     & 60 \\    \hline
          &       & SubTotal (\textbf{Tecnologia da Computação}) & \textbf{40}    & \textbf{40}    & \textbf{600} \\          \hline
\multicolumn{6}{l}{Núcleo: \textbf{Contexto Social e Profissional}}\\     \hline
    31    & III   & Metodologia Científica & 2     & 2     & 30 \\     \hline
    32    & III   & Inglês & 2     & 2     & 30 \\    \hline
    33    & VI    & Empreendedorismo & 4     & 4     & 60 \\    \hline
    34    & VII  & Computadores e Sociedade & 2     & 2     & 30 \\ \hline
    35    & IX  & Estágio & 8     & 8     & 120 \\           \hline
          &       & SubTotal (\textbf{Contexto Social e Profissional}) & \textbf{18}    & \textbf{18}    & \textbf{270} \\    \hline
\multicolumn{6}{l}{\textbf{Optativas}} \\ \hline
    36    & VII   & Optativa 1 & 4     & 4     & 60 \\ \hline
    37    & VII   & Optativa 2 & 4     & 4     & 60 \\ \hline
    38    & VIII  & Optativa 3 & 4     & 4     & 60 \\ \hline
    39    & VIII  & Optativa 4 & 4     & 4     & 60 \\ \hline
    40    & VIII  & Optativa 5 & 4     & 4     & 60 \\ \hline
    41    & VIII  & Optativa 6 & 4     & 4     & 60 \\ \hline
    42    & VIII  & Optativa 7 & 4     & 4     & 60 \\ \hline
    \textbf{} & \textbf{} & SubTotal (\textbf{Optativas}) & \textbf{28} & \textbf{28} & \textbf{420} \\ \hline
\end{longtable}%


% ----------------------------------------------------------
\chapter{Disciplinas Optativas}
% ----------------------------------------------------------

\begin{longtable}{r|r|l|r|r|r}
     \caption{Disciplinas Optativas}
     \label{ltab:teste}
     \\
%    \toprule
\hline
%  \rowcolor[rgb]{0.91,0.89,0.88}
    N.    & Período & Disciplinas                       & h/semana & Créditos & CH \\ \hline

\multicolumn{6}{l}{\textbf{Área Temática: Banco de Dados} } \\ \hline
    1     &       & Banco de Dados Avançados & 4     & 4     & 60 \\ \hline
    2     &       & Projeto de Banco de Dados & 4     & 4     & 60 \\ \hline
    3     &       & Integração de Dados e Data-Warehouse & 4     & 4     & 60 \\ \hline
    4     &       & Modelagem Conceitual de Dados & 4     & 4     & 60 \\ \hline
    5     &       & Tópicos Especiais em Banco de Dados & 4     & 4     & 60 \\ \hline

\multicolumn{6}{l}{\textbf{Área Temática: Engenharia da Computação}}\\ \hline
    1     &       & Sistemas de Tempo Real & 4     & 4     & 60 \\ \hline
    2     &       & Projeto de Sistemas Embarcados & 4     & 4     & 60 \\ \hline
    3     &       & Prototipação de Circuitos Digitais & 4     & 4     & 60 \\ \hline
    4     &       & \small{Tóp. Avançados em Engenharia da Computação} & 4     & 4     & 60 \\ \hline

\multicolumn{6}{l}{\textbf{Área Temática: Engenharia de Software}} \\ \hline
    1     &       & Engenharia de Software Experimental & 4     & 4     & 60 \\ \hline
    2     &       & Programação Orientada a Aspectos & 4     & 4     & 60 \\ \hline
    3     &       & Programação Orientada a Objetos II & 4     & 4     & 60 \\ \hline
    4     &       & Programação Paralela e Distribuída (Clusters) & 4     & 4     & 60 \\ \hline
    5     &       & Programação Web & 4   & 4  & 60 \\ \hline
    6     &       & Desenvolvimento Distribuído de Software & 4     & 4     & 60 \\ \hline
    7     &       & Análise e Projeto de Sistemas & 4     & 4     & 60 \\ \hline
    8     &       & Gerenciamento de Projetos & 4     & 4     & 60 \\ \hline
    9     &       & Engenharia de Requisitos & 4     & 4     & 60 \\ \hline
    10    &       & Qualidade de Software & 2     & 2     & 30 \\ \hline
    11    &       & Teste de Software & 4     & 4     & 60 \\ \hline
    12    &       & Estimativas e Medição de Software & 2     & 2     & 30 \\ \hline
    13    &       & Desenvolvimento de Aplicações Móveis & 4     & 4     & 60 \\ \hline
    14    &       & Metodologias Ágeis & 4     & 4     & 60 \\ \hline
    15    &       & \small{Tópicos Especiais em Engenharia de Software} & 4     & 4     & 60 \\ \hline

\multicolumn{6}{l}{\textbf{Área Temática: Inteligência Computacional} }  \\ \hline
    1     &       & Redes Neurais & 4     & 4     & 60 \\ \hline
    2     &       & Computação Evolucionária & 4     & 4     & 60 \\ \hline
    3     &       & Mineração de Dados & 4     & 4     & 60 \\ \hline
    4     &       & Visão Computacional & 4     & 4     & 60 \\ \hline
    5     &       & Agentes Inteligentes & 4     & 4     & 60 \\ \hline
    6     &       & Aprendizagem de Máquina & 4     & 4     & 60 \\ \hline
    7     &       & Biometria & 4     & 4     & 60 \\ \hline
    8     &       & Sistemas de Recomendação & 4     & 4     & 60 \\ \hline
    9     &       & Reconhecimento de Padrões II & 4     & 4     & 60 \\ \hline
    10    &       & Sistemas Hibridos & 4     & 4     & 60 \\ \hline
    11    &       & Computacao Forense & 4     & 4     & 60 \\ \hline
    12    &       & \small{Métodos Paramétricos Aprendizagem de Máquina} & 4     & 4     & 60 \\ \hline
    13    &       & Bioinformática & 4     & 4     & 60 \\ \hline
    14    &       & Introdução à Computação Quântica & 4    & 4 & 60 \\ \hline 
    15    &       & \small{Tópicos Especiais em Inteligência Computacional} & 4     & 4     & 60 \\ \hline
    16    &       & \small{Tópicos Especiais em Aprendizagem de Máquina} & 4     & 4     & 60 \\ \hline
    17    &       & \small{Tópicos Especiais em Redes Neurais} & 4     & 4     & 60 \\ \hline
    18    &       & \small{Tópicos Especiais em Computação Evolucionária} & 4     & 4     & 60 \\ \hline
    19    &       & \small{Tópicos Especiais em Inteligência Artificial} & 4     & 4     & 60 \\ \hline

\multicolumn{6}{l}{\textbf{Área Temática: Matemática e Simulação Computacional}} \\ \hline
    1     &       & Análise Numérica & 4     & 4     & 60 \\ \hline
    2     &       & Cálculo Numérico e Computacional & 4     & 4     & 60 \\ \hline
    3     &       & Métodos de Otimização & 4     & 4     & 60 \\ \hline
    4     &       & Modelagem Matemática & 4     & 4     & 60 \\ \hline
    5     &       & Programação Matemática & 4   & 4  & 60 \\ \hline
    6     &       & Simulação de Sistemas & 4     & 4     & 60 \\ \hline
    7     &       & Cálculo para Computação III & 4     & 4     & 60 \\ \hline
    8     &       & \small{Tópicos Especiais em Matemática Computacional} & 4     & 4     & 60 \\ \hline
    9     &       & \small{Tópicos Especiais em Pesquisa Operacional} & 4     & 4     & 60 \\ \hline

\multicolumn{6}{l}{\textbf{Área Temática: Metodologia e Técnicas da Computação}} \\ \hline
    1     &       & \small{Análise Estatística para Métodos Quantitativos} & 4     & 4     & 60 \\ \hline
    2     &       & Métodos de Pesquisa em Computação & 4     & 4     & 60 \\ \hline
    3     &       & Análise de Dados Simbólicos & 4     & 4     & 60 \\ \hline
    4     &       & Métodos de Pesquisa Qualitativa & 4     & 4     & 60 \\ \hline
    5     &       & Revisão Sistemática de Literatura & 4     & 4     & 60 \\ \hline
    6     &       & \small{Tópicos Especiais em Metodologia de Pesquisa} & 4     & 4     & 60 \\ \hline

\multicolumn{6}{l}{\textbf{Área Temática: Mídia e Interação}} \\ \hline
    1     &       & Processamento de Imagens I & 4     & 4     & 60 \\ \hline
    2     &       & Processamento de Imagens II & 4     & 4     & 60 \\ \hline
    3     &       & Modelagem Geométrica & 4     & 4     & 60 \\ \hline 
    4     &       & Realidade Virtual e Aumentada & 4     & 4     & 60 \\ \hline 
    5     &       & Processamento Digital de Sinais & 4     & 4     & 60 \\ \hline
    6     &       & Processamento de Audio & 4     & 4     & 60 \\ \hline
    7     &       & Desenvolvimento de Jogos & 4     & 4     & 60 \\ \hline
    8     &       & Princípios da Animação Gráfica 3D & 4     & 4     & 60 \\ \hline
    9     &       & Wavelets & 4     & 4     & 60 \\ \hline
   10     &       & \small{Tópicos Especiais em Processamento de Imagens} & 4     & 4     & 60 \\ \hline
   11     &       & \small{Tópicos Especiais em Processamento de Sinais} & 4     & 4     & 60 \\ \hline
   12     &       & \small{Tópicos Especiais em Mídia e Interação} & 4     & 4     & 60 \\ \hline
   13     &       & \small{Tópicos Especiais em Computação Gráfica} & 4     & 4     & 60 \\ \hline
   14     &       & \small{Tóp. Especiais em Realidade Virtual e Aumentada} & 4     & 4     & 60 \\ \hline

\multicolumn{6}{l}{\textbf{Área Temática: Redes e Sistemas Distribuídos}} \\ \hline
    1     &       & Segurança de Redes de Computadores & 4     & 4     & 60 \\ \hline
    2     &       & Gerenciamento de Redes de Computadores & 4     & 4     & 60 \\ \hline
    3     &       & \small{Tópicos Especiais Redes e Sistemas Distribuídos} & 4     & 4     & 60 \\ \hline

\multicolumn{6}{l}{\textbf{Área Temática: Tecnologia Educacional}} \\ \hline
    1     &       & Informática na Educação & 4     & 4     & 60 \\ \hline
    2     &       & Desenvolvimento de Software Educacional & 4     & 4     & 60 \\ \hline
    3     &       & Tecnologia Educacional e Cognição & 4     & 4     & 60 \\ \hline
    4     &       & Projeto de Sistemas Educacionais & 4     & 4     & 60 \\ \hline
    5     &       & Linguagem Brasileira de Sinais & 4     & 4     & 60 \\ \hline
    6     &       & \small{Tópicos Especiais em Tecnologia Educacional} & 4     & 4     & 60 \\ \hline

\multicolumn{6}{l}{\textbf{Área Temática: Tecnologias da Informação}} \\ \hline
    1     &       & Gestão da TI & 4     & 4     & 60 \\ \hline
    2     &       & Gestão de Processos de Negócio & 4     & 4     & 60 \\ \hline
    3     &       & Gestão de Serviços em TI & 4     & 4     & 60 \\ \hline
    4     &       & Governança em TI & 4     & 4     & 60 \\ \hline
    5     &       & Contabilidade e Administração de Empresas & 4     & 4     & 60 \\ \hline
    6     &       & Contabilidade e Administração Financeira & 4     & 4     & 60 \\ \hline
    7     &       & \small{Tópicos Especiais em Gestão de Projetos} & 4     & 4     & 60 \\ \hline


    \end{longtable}%



% %-----------------------------------------------------------
% \chapter{Pré e co-requisitos}
% %-----------------------------------------------------------

\begin{landscape}

    \begin{longtable}{r|r|l|l|l}
     \caption{Disciplinas com os pré e co-requisitos}
     \label{ltab:PreReq}
     \\
     \hline
    \textbf{N.} & \textbf{Período} & \textbf{Disciplinas} & \textbf{Pré-requisitos} & \textbf{Co-requisitos} \\
\hline
\multicolumn{5}{l}{\textbf{Núcleo: Matemática}}\\ \hline
    1     & I     & Cálculo p/ Computação I &       &       \\ \hline
    2     & I     & Geometria Analítica &       &        \\ \hline
    3     & I     & Lógica Matemática &       &        \\ \hline
    4     & II    & Cálculo p/ Computação II & Cálculo p/ Computação I &       \\ \hline
    5     & II    & Algebra Linear & Geometria Analítica &        \\ \hline
    6     & III   & Probabilidade e Estatística & Cálculo p/ Computação II &       \\ \hline
    7     & III   & Matemática Discreta & Lógica Matemática &        \\ \hline

\multicolumn{5}{l}{\textbf{Núcleo: Fundamentos da Computação}}\\ \hline
    8     & I     & Introdução à Computação &       &        \\ \hline
    9     & I     & Introdução à Programação &       &      \\ \hline
    10    & II    & Algoritmos e Estruturas de Dados I & Introdução à Programação &       \\ \hline
    11    & II    & Programação Orientada a Objetos & Introdução à Programação &       \\ \hline
    12    & III   & Sistemas Digitais & Física p/ Computação &       \\ \hline
    13    & III   & Algoritmos e Estruturas de Dados II & Algoritmos e Estruturas de Dados I &       \\ \hline
    14    & IV    & Arquitetura de Computadores & Sistemas Digitais &       \\ \hline
    15    & IV    & Projeto e Análise de Algoritmos & Algoritmos e Estruturas de Dados II &       \\ \hline
    16    & IV    & Paradigmas de Linguagens de Programação & Algoritmos de Estruturas de Dados II &       \\ \hline
    17    & V     & Sistemas Operacionais & Arquitetura de Computadores &       \\ \hline
    18    & V     & Sistemas de Informação e Tecnologias &       &       \\ \hline
    19    & V     & Teoria da Computação & Matemática Discreta &        \\ \hline

\multicolumn{5}{l}{\textbf{Núcleo: Ciências Básicas}} \\ \hline
    20    & II    & Física p/ Computação & Cálculo p/ Computação I &        \\ \hline

\multicolumn{5}{l}{\textbf{Núcleo: Tecnologia da Computação}} \\ \hline
    21    & IV    & Bancos de Dados & Algoritmos e Estruturas de Dados II &        \\ \hline
    22    & IV    & Engenharia de Software & Programação Orientada a Objetos &       \\ \hline
    23    & V     & Redes de Computadores & Introdução à Computação &       \\ \hline
    24    & V     & Inteligência Artificial & Projeto e Análise de Algoritmos &       \\ 
          &       &       & Paradigmas de Ling. de Programação &          \\ \hline
    25    & VI    & Reconhecimento de Padrões &  Algoritmos e Estruturas de Dados II &  \\
          &       &       & Matemática Discreta &       \\ \hline
    26    & VI    & Compiladores & Teoria da Computação &       \\ 
              &       &       & Algoritmos e Estruturas de Dados II &         \\ \hline
    27    & VI    & Computação Gráfica & Álgebra Linear &        \\ 
          &       &       & Introdução a Computação &         \\ \hline
    28    & VI    & Sistemas Distribuídos & Redes de Computadores &        \\ \hline
    29    & VII   & Projeto de Desenvolvimento & Engenharia de Software &   \\ 
          &       &       & Banco de Dados &        \\ 
          &       &       & Empreendimentos em TIC &        \\ \hline
    30    & VII   & Interação Humano-Computador & Engenharia de Software &        \\ \hline
    31    & IX    & Trabalho de Conclusão de Curso & Disciplinas Cursadas &       \\ \hline

\multicolumn{5}{l}{\textbf{Núcleo: Contexto Social e Profissional}} \\ \hline
    32    & III   & Metodologia Científica &       &       \\ \hline
    33    & III   & Inglês &       &        \\ \hline
    34    & VI    & Empreendimentos em TIC &       &        \\ \hline
    35    & VIII  & Computadores e Sociedade &       &        \\ \hline


    \end{longtable}%
\end{landscape}

%---------------------------------------------------------

%\chapter{Ementário}
%---------------------------------------------------------

%\chapter{Ementário}
%---------------------------------------------------------

%\singlespacing
\setlength{\parindent}{0in}
\chapter{Ementas das Disciplinas Obrigatórias}
\section*{1o. Período}

IDENTIFICAÇÃO

DISCIPLINA: Cálculo para Computação I CÓDIGO:\\
DEPARTAMENTO: Unidade Acadêmica de Garanhuns\\
ÁREA:\\
CARGA HORÁRIA TOTAL: 60 NÚMERO DE CRÉDITOS: 04\\
CARGA HORÁRIA SEMANAL: 04 TEÓRICAS: 04 PRÁTICAS: 00 TOTAL:04\\
PRÉ-REQUISITOS: não há\\
CO-REQUISITOS:não há

EMENTA

Conjuntos numéricos. Funções elementares: linear, afim, quadrática,
modular. Funções diretas e inversas. Funções exponenciais e
logarítmicas. Introdução à trigonometria. Funções trigonométricas.
Limite e continuidade. Derivadas e aplicações.

BIBLIOGRAFIA

Básica:

STEWART, JAMES. Cálculo, vol 1, Editora Pioneira Thomson Learning,
2001.\\
ANTON, HOWARD. Cálculo, um novo horizonte, vol 1, 6. edição, Editora
Bookman; reimpressão 2004.\\
GUIDORIZZI, HAMILTON L.. Um Curso de Cálculo, vol.1, 5. edição, Ao Livro
Técnico S.A., 2002.\\
Complementar:\\
LEITHOLD, LOUIS. Cálculo com Geometria Analítica, vol. 1, 3. edição,
Editora Harbra, 2002.\\
SWOKOSWKI, E. William. Cálculo com Geometria Analítica, vol. 1, 2.
edição, Editora Makron Books do Brasil, 1995.\\

\newpage IDENTIFICAÇÃO

DISCIPLINA: Geometria Analítica CÓDIGO:\\
DEPARTAMENTO: Unidade Acadêmica de Garanhuns\\
ÁREA:\\
CARGA HORÁRIA TOTAL: 60 HORAS NÚMERO DE CRÉDITOS: 04\\
CARGA HORÁRIA SEMANAL: 04 TEÓRICAS: 04 PRÁTICAS: 00 TOTAL: 04\\
PRÉ-REQUISITOS: não há\\
CO-REQUISITOS: não há

EMENTA

Vetores no $R^2$; Produto interno no $R^2$; Estudo da reta no $R^2$; Lugares geométricos no $R^2$ (circunferência, elipse, parábola e hipérbole); Vetores no R3; Produto interno no $R^3$, vetorial e misto; Aplicações: áreas e volumes; Equação da reta no $R^3$ e equação do plano; Equação da superfície esférica.

BIBLIOGRAFIA

Básica:\\
BOULOS, P. Geometria Analítica: um tratamento vetorial, Mc-Graw Hill, 1987.\\
STEINBUCH, Alfredo e WINTERLE, Paulo - Geometria Analítica, Mc Graw-Hill, 1987.\\
LEITHOLD, L. O Cálculo com geometria analítica. v. 2. 3 ed. São Paulo :
Harbra, 1994.\\
Complementar:\\
VENTURI, J.   Álgebra Vetorial e Geometria Analítica, Editora UFPR, Curitiba, 1989.\\

\newpage IDENTIFICAÇÃO

DISCIPLINA: Lógica Matemática CÓDIGO:\\
DEPARTAMENTO: Unidade Acadêmica de Garanhuns\\
ÁREA: \\
CARGA HORÁRIA TOTAL: 60 NÚMERO DE CRÉDITOS: 04\\
CARGA HORÁRIA SEMANAL: 04 TEÓRICAS: 04 PRÁTICAS: 00 TOTAL: 04\\
PRÉ-REQUISITOS: não há\\
CO-REQUISITOS: não há

EMENTA

Proposições e conectivos; Operações lógicas sobre proposições;
Construção de tabelas-verdade; Tautologias, contradições e
contingencias; Implicação lógica; Equivalência lógica; Álgebra das
proposições; Método dedutivo; Argumentos, regras de inferência;
Validade mediante tabela verdade; Validade mediante regras de
inferência; Métodos de demonstrações; Sentenças abertas; Operações
lógicas sobre sentenças abertas; Quantificadores.

BIBLIOGRAFIA

Básica: \\
GERSTING, JUDITH L. Fundamentos matemáticos para a ciência da
computação: um tratamento moderno de matemática discreta. LTC, 2008.\\
ALENCAR FILHO, Edgard de, Iniciação à Lógica Matemática. 18. ed. 203 p,
São Paulo : Nobel, 2000.\\
DAGHLIAN, Jacob ; Lógica e Álgebra de Boole, 4. ed. 167 p., São Paulo :
Atlas, 1995.\\
Complementar: \\
DEL PICCHIA, Walter; Métodos Numéricos para Resolução de Problemas
Lógicos. São Paulo : Edgard Blücher, 1993.\\
SALMON, Wesley C. Lógica, 3a Edição. Rio de Janeiro: LTC Editora, 2002.\\


\newpage IDENTIFICAÇÃO

DISCIPLINA: Introdução à Ciência da Computação CÓDIGO:\\
DEPARTAMENTO: Unidade Acadêmica de Garanhuns\\
ÁREA: \\
CARGA HORÁRIA TOTAL: 30 NÚMERO DE CRÉDITOS: 02\\
CARGA HORÁRIA SEMANAL: 02 TEÓRICAS: 02 PRÁTICAS: 00 TOTAL: 02\\
PRÉ-REQUISITOS: não há\\
CO-REQUISITOS: não há
 
EMENTA 
 
Introdução à Ciência da Computação: a ciência, o curso e a profissão.
História e evolução da Ciência da Computação. Conceitos básicos.
Classificação de sistemas computacionais. Noções de sistemas
operacionais, redes, tipos de linguagens, compiladores e
interpretadores. Tópicos complementares. Seminário.

BIBLIOGRAFIA 

Básica:\\
Introdução à Ciência da Computação; Ricardo Daniel Fedeli, Enrico
Giulio, Franco Polloni, Fernando Eduardo Peres. Thomson Pioneira, 2010.
2. Edição. ISBN 9788522108459.\\
Introdução à informática; H. L. Capron, J. A. Johnson ; tradução José
Carlos\\
Barbosa dos Santos. São Paulo, SP : Prentice Hall, 2006. 8. Edição. ISBN: 8587918885.\\
BROOKSHEAR, J. Glenn. Ciência da Computação: uma visão abrangente; Porto
Alegre: Bookmam, 2005. 7. Edição. ISBN: 8573075376.\\
Complementar:\\
Introdução À Ciência da Computação; Fábio Carneiro Mokarzel; Nei
Yoshihiro Soma. Editora Campus, 2008. 1. Edição. ISBN 8535218793.\\
Introdução à informática, Peter Norton. Makron Books, 2005. ISBN
8534605157.\\
Informática Conceitos Básicos, Fernando de Castro Velloso. Campus, 2004.
ISBN 8535215360.\\


\newpage IDENTIFICAÇÃO

DISCIPLINA: Introdução à Programação CÓDIGO:\\
DEPARTAMENTO: Unidade Acadêmica de Garanhuns\\
ÁREA: \\
CARGA HORÁRIA TOTAL: 90 NÚMERO DE CRÉDITOS: 06\\
CARGA HORÁRIA SEMANAL: 06 TEÓRICAS: 03 PRÁTICAS: 03 TOTAL: 06\\
PRÉ-REQUISITOS: não há\\
CO-REQUISITOS: não há

EMENTA 

Introdução a algoritmos e pseudocódigos. Introdução à programação
imperativa: variáveis, constantes e expressões. Controle de fluxo de
execução e repetição. Estruturas triviais de dados: vetores, matrizes e
registros. Noções de funções. Comandos de atribuição e declaração de
constantes, variáveis e tipos de dados. Expressões. Ponteiros.
Instruções condicionais de controle de fluxo. Bibliotecas definidas
pelo usuário. Recursividade. Alocação dinâmica de memória.

BIBLIOGRAFIA 

Básica:\\\
Fundamentos da programação de computadores. Ana Fernanda Gomes Ascencio,
Edilene Aparecida Veneruchi de Campo. 2. Edição. Prentice Hall. ISBN
8587918362. (2002).\\
Lógica de Programação. A construção de Algoritmos e estruturas de dados.
André Luiz e Henri Frederico. 3. Edição. Pearson. Prentice Hall. ISBN:
8576050242.\\
C Completo e Total. Herbert Schildt. Pearson, 3. Edição, 2009. ISBN.
8534605955.\\
Treinamento em Linguagem C. Victorine Viviane Mizrahi. 2. Edição.
Prentice Hall.ISBN 9788576051916.\\
Complementar:\\
Introdução à programação: 500 algorítmos resolvidos. Anita Lopes e Guto
Garcia. Campus. ISBN 8535210199 (2002).\\
Algoritmos e Programação, Teoria e Prática. Marco Medina e Cristina
Fertig. Novatec. ISBN 857522073X. (2005).\\
Projetos de Algorítmos com implementações em C e pascal. Nivio Ziviani.
Thomson Learning (Pioneira). ISBN 8522103909. (2004).\\
\newpage

\section*{2o. Período}

IDENTIFICAÇÃO

DISCIPLINA: Cálculo para Computação II CÓDIGO:\\ 
DEPARTAMENTO: Unidade Acadêmica de Garanhuns\\
ÁREA: \\
CARGA HORÁRIA TOTAL: 60 NÚMERO DE CRÉDITOS: 04\\
CARGA HORÁRIA SEMANAL: 04 TEÓRICAS: 04 PRÁTICAS: 00 TOTAL: 04\\
PRÉ-REQUISITOS: Cálculo para Computação I\\
CO-REQUISITOS: não há

EMENTA 

Integrais indefinidas, definidas e Integrais impróprias. Seqüências e
séries numéricas. Série de potência. Curvas planas e coordenadas
polares. Funções reais de várias variáveis, Limites e continuidade de
funções de várias variáveis, derivadas parciais, diferenciabilidade,
máximos e mínimos de funções. Integrais múltiplas.

BIBLIOGRAFIA 

Básica: \\
STEWART, James. Cálculo, vol 2, Editora Pioneira Thomson Learning, 2001.\\
ANTON, Howard. Cálculo, um novo horizonte, vol 2; 6. edição, Editora
Bookman; reimpressão 2004.\\
GUIDORIZZI, Hamilton L. Um Curso de Cálculo, vol.2, 5. edição, Ao Livro
Técnico S.A., 2002.\\
Complementar:\\
LEITHOLD, Louis. Cálculo com Geometria Analítica, vol. 2, 3a edição,
Editora Harbra, 2002.\\
SWOKOSWKI, Earl William. Cálculo com Geometria Analítica, vol. 1, 2.
edição, Editora Makron Books do Brasil, 1995.\\


\newpage IDENTIFICAÇÃO

DISCIPLINA: Álgebra Linear CÓDIGO:\\ 
DEPARTAMENTO: Unidade Acadêmica de Garanhuns\\
ÁREA: \\
CARGA HORÁRIA TOTAL: 60 HORAS NÚMERO DE CRÉDITOS: 04\\
CARGA HORÁRIA SEMANAL: 04 TEÓRICAS: 04 PRÁTICAS: 00 TOTAL: 04\\
PRÉ-REQUISITOS: Geometria Analítica\\
CO-REQUISITOS: não há

EMENTA 

Vetores; Matrizes; Operações elementares e sistemas de equações; Espaços vetorias (subespaços, dependência e independência linear, base e
dimensão, espaço linha, espaço coluna e posto, dimensão do conjunto solução de um sistema linear); Determinantes; Transformações lineares
(núcleo e imagem, representação por matrizes, mudança de base, auto valor, auto vetor e diagonalização).

BIBLIOGRAFIA 

Básica:\\
STEINSBRUCH, Alfredo  Álgebra Linear. Editora: Makron Books.\\
LIPSCHUTZ, Seymour  Álgebra Linear  Editora: Mc Graw Hill.\\
BOLDRINI, J. L. Álgebra Linear. Harbra. 2006.\\
Complementar:\\
LAWSON, Terry. Álgebra Linear. Edgar Blucher. 2008.\\


\newpage IDENTIFICAÇÃO

DISCIPLINA: Algoritmos e Estruturas de Dados I CÓDIGO:\\ 
DEPARTAMENTO: Unidade Acadêmica de Garanhuns\\
ÁREA: \\
CARGA HORÁRIA TOTAL: 60 HORAS NÚMERO DE CRÉDITOS: 04\\
CARGA HORÁRIA SEMANAL: 04 TEÓRICAS: 02 PRÁTICAS: 02 TOTAL: 04\\
PRÉ-REQUISITOS: Introdução à Programação\\
CO-REQUISITOS: não há

EMENTA 

Resolução de problemas e desenvolvimento de algoritmosConceitos de C:
Modularização de programas e tratamento de arquivo. Tipo de dado
abstrato. Algoritmos de Ordenação e Busca. Estruturas de dados
elaboradas (lista, fila e pilha) e seus algoritmos. Análise do
problema, estratégias de solução, representação e documentação.

BIBLIOGRAFIA 

Básica:\\
TENENBAUM, A.; LANGSAM, Y. \& AUGENSTEIN, M.J.-Data Structures Using C
And C++, 2nd Edition, Prentice-Hall, 1996.\\
CELES ,W.; CERQUEIRA, R. e RANGEL, J.L.- Introdução a Estrutura de Dados
Uma Introdução com Tecnicas de Programação em C, Coleção: Campus/SBC,
Rio de Janeiro: Campus, 2004.\\
Algoritmos e Estruturas de Dados. A. de M. Guimarães e N. A. de C.
Lages. Editora LTC, 1994.\\
Complementar:\\
Estruturas de Dados e seus Algoritmos. J. L. Szwarcfiter e L.
Markenzion. Segunda Edição. LTC, 1994.\\
Desenvolvimento de Algoritmos e Estruturas de Dados. R. Terada, Makron
Books, 1991.\\
ZIVIANI, N., Projeto de Algoritmos: com implementação em Pascal e C, 2a
Ed., Editora Thomson Pioneira, 2004.\\


\newpage IDENTIFICAÇÃO

DISCIPLINA: Programação Orientada a Objetos CÓDIGO:\\ 
DEPARTAMENTO: Unidade Acadêmica de Garanhuns\\
ÁREA: \\
CARGA HORÁRIA TOTAL: 60 HORAS NÚMERO DE CRÉDITOS: 04\\
CARGA HORÁRIA SEMANAL: 04 TEÓRICAS: 02 PRÁTICAS: 02 TOTAL: 04\\
PRÉ-REQUISITOS: Introdução à Programação\\
CO-REQUISITOS: não há

EMENTA 

Contextualização de Programação Orientada a Objetos. Conceitos de
orientação a Objetos: Objetos, Operações, Mensagens, Métodos e Estados.
Ambientes de Desenvolvimento de Software Orientado a Objetos. Classes,
Métodos e Objetos. Tipo de Entrada de Dados. Tipo de Saída de Dados.
Estruturas de Controle e Repetição. Modificadores de Classe e de
Métodos. Construtores e Finalizadores. Herança Simples e Múltipla:
Super e sub classes. Polimorfismo, Abstrações e Generalizações.
Diagramas UML. Interface Gráfica com Usuário.

BIBLIOGRAFIA 

Básica:\\
Harvey M. Deitel, Paul J. Deitel. Java How to Program. Prentice Hall. 7.
Edição, 2006. ISBN : 0132222205.\\
Lynn Andrea Stein. Interactive Programming in Java, 2003.\\
Bruce Eckel. Thinking in Java. Segunda edição, 2000.\\
Complementar:\\
Barnes, D. J., Kölling, M. Programação Orientada a Objetos com Java, Ed.
Pearson/Prentice-Hall, 2004.\\
Santos, R. Introdução à Programação Orientada à Objetos Usando Java, Ed.
Campus, 2003.\\
Cay S. Horstmann e Gary Cornell. Core Java 2, Volume I  Fundamentos.
Sun Microsystems Press, Makron Books do Brasil, 2001.\\


\newpage IDENTIFICAÇÃO

DISCIPLINA: Física para Computação CÓDIGO:\\
DEPARTAMENTO: Unidade Acadêmica de Garanhuns\\ ÁREA: \\
CARGA HORÁRIA TOTAL: 60 HORAS NÚMERO DE CRÉDITOS: 04\\
CARGA HORÁRIA SEMANAL: 04 TEÓRICAS: 02 PRÁTICAS: 02 TOTAL: 04\\
PRÉ-REQUISITOS: Cálculo I para Computação\\
CO-REQUISITOS: não há

EMENTA 

Conceito de: (a) Carga elétrica, (b) Campo elétrico, (c) Potencial
elétrico, (d) corrente elétrica, (e) potência elétrica; Resistência
elétrica e lei de Ohm; associação de resistores: associação em série e
em paralelo, transformação estrela-triângulo; bateria elétrica;
circuitos resistivos e leis de Kirchhoff; Capacitor e circuitos RC;
Fontes do campo magnético, solenóide e imãs; Indutor, auto-indução,
indutância mútua, circuitos RL; Corrente alternada, circuitos RLC,
transformadores, motores e geradores elétricos; Espectro
eletromagnético, propagação de ondas eletromagnéticas, lasers; Metais,
isolantes e semicondutores; Diodo e circuitos com diodos; Transistor e
circuitos com transistor; Circuitos eletrônicos básicos.

BIBLIOGRAFIA 

Básica:\\\
Fundamentos de Física, volume 3, Halliday, Resnick e Walker. ISBN.:
8521614869. LTC. (2006).\\
Eletrônica, volume 1 e 2, A. P. Malvino. 7. edição McGraw-Hill,
Laboratório de Eletricidade e Eletrônica, Francisco Gabriel Capuano e
Maria Aparecida Mendes Marinho. 24. edição Érica.\\
Complementar:\\
Física, volume 2, Paul A. Tipler e Gene Mosca. 5. edição LTC\\
Física Conceitual, Paul G. Hewitt. 9. edição. Bookman\\
Dispositivos eletrônicos e teoria de circuitos. 8. edição. Robert L.
Boylestad e Louis Nashelsky. Pearson.
\newpage

\section*{3o. Período}

IDENTIFICAÇÃO 

DISCIPLINA: Probabilidade e Estatística CÓDIGO:\\
DEPARTAMENTO: Unidade Acadêmica de Garanhuns\\
ÁREA: \\
CARGA HORÁRIA TOTAL: 60 HORAS NÚMERO DE CRÉDITOS: 04\\
CARGA HORÁRIA SEMANAL: 04 TEÓRICAS: 04 PRÁTICAS: 00 TOTAL: 04\\
PRÉ-REQUISITOS: Cálculo para Computação II
CO-REQUISITOS: não há

EMENTA 

Análise combinatória. Planejamento de uma pesquisa. Análise exploratória
de dados. Probabilidade. Variáveis aleatórias discretas e contínuas.
Principais modelos teóricos. Estimação de parâmetros. Testes de
hipóteses.

BIBLIOGRAFIA 

Básica:\\\
BARBETTA, P. A.; REIS, M. M., BORNIA, A. C. Estatística para Cursos de Engenharia e Informática . São Paulo: Editora Atlas, 2004;\\
BUSSAB, W. O., MORETTIN, P. A. Estatística básica . 5 ed. São Paulo: Editora Saraiva, 2002.\\
Complementar:\\
MAGALHÃES, A. N., LIMA, A.C.P.  Noções de probabilidade e estatística. 6 ed. São Paulo: EDUSP, 2005.\\
MONTGOMERY, D.C., RUNGER, G. C. Estatística aplicada e probabilidade para engenheiros . Rio de Janeiro: LTC, 2003;\\
TRIOLA, M. F.  Introdução à Estatística . 9. ed. Rio de Janeiro: LTC, 2005.\\

\newpage IDENTIFICAÇÃO

DISCIPLINA: Matemática Discreta para Computação CÓDIGO:\\
DEPARTAMENTO: Unidade Acadêmica de Garanhuns\\
ÁREA: \\
CARGA HORÁRIA TOTAL: 60 NÚMERO DE CRÉDITOS: 04\\
CARGA HORÁRIA SEMANAL: 04 TEÓRICAS:04 PRÁTICAS: 00 TOTAL: 04\\
PRÉ-REQUISITOS: Lógica Matemática\\
CO-REQUISITOS: não há

EMENTA 

Grafos e árvores. Álgebra de boole. Inteiros, divisores e primos.

BIBLIOGRAFIA 

Básica:\\
EDWARD, Scheinerman, Matemática Discreta: Uma Introdução. Thomson
Learning, (2003).\\
JUDITH, L. Gersting, Fundamentos Matemáticos para Ciência da Computação
 5a edição. LTC (2004).\\
LOVÀSZ, L. PELIKÁN, J. Matemática Discreta. SBM. 2003.\\
Complementar:\\
SEYMOUR, Lipschutz, MARC, Lipson, Teoria e Problemas da Matemática
Discreta  coleção Shaum. Bookaman (2004).\\


\newpage IDENTIFICAÇÃO

DISCIPLINA: Metodologia Científica CÓDIGO:\\
DEPARTAMENTO: Unidade Acadêmica de Garanhuns\\
ÁREA: \\
CARGA HORÁRIA TOTAL: 30 NÚMERO DE CRÉDITOS: 02\\
CARGA HORÁRIA SEMANAL: 02 TEÓRICAS: 02 PRÁTICAS: 00 TOTAL: 02\\
PRÉ-REQUISITOS: não há\\
CO-REQUISITOS: não há

EMENTA 

O método científico e a prática da pesquisa; função social da pesquisa e
da ciência. Tipos e características da pesquisa. Instrumentalização
metodológica. Projeto de pesquisa. Relatório de pesquisa.

BIBLIOGRAFIA 

Básica:\\
LAKATOS \& MARCONE. Fundamentos de Metodologia Cientifica. 7. ed. 320 p.
Atlas, 2010\\
LAKATOS \& MARCONE. Metodologia do Trabalho Científico. 7. ed. 228 p.
São Paulo : Atlas, 2007\\
WAZLAWICK, Raul Sidney. Metodologia De Pesquisa Para Ciência Da
Computaçao. 184p. Campus. 2009.\\
Complementar:\\
CHALMERS, A. O quê é ciência afinal. 2. Ed. 230p. Brasiliense. 2009.\\
SIQUEIRA et. al. Como Elaborar Projetos De Pesquisa. 140p, Editora FGV
2007.\\
MATIAS-PEREIRA. Manual de Metodologia da Pesquisa Cientifica. 2. Ed.
240p. Atlas. 2010.\\

\newpage IDENTIFICAÇÃO

DISCIPLINA: Sistemas Digitais CÓDIGO:\\ 
DEPARTAMENTO: Unidade Acadêmica de Garanhuns\\
ÁREA: \\
CARGA HORÁRIA TOTAL: 60 NÚMERO DE CRÉDITOS: 04\\
CARGA HORÁRIA SEMANAL: 04 TEÓRICAS: 03 PRÁTICAS: 01 TOTAL: 04\\
PRÉ-REQUISITOS: Física para computação\\
CO-REQUISITOS: não há

EMENTA 

Sistemas de Numeração e Códigos; Aritmética Binária; Porta Lógicas;
Análise e Projeto de Circuitos Combinacionais; Minimização por Mapa de
Karnaugh; Somadores; Decodificadores; Codificadores; Multiplexadores;
Demultiplexadores; Análise e Síntese de Circuitos Sequenciais; Latches
e Flip-Flops; Minimização de Estado; Registradores; Registradores de
Deslocamento; Dispositivos Lógicos Programáveis; Memória).

BIBLIOGRAFIA 

Básica:\\
TOCCI, Ronald J. Sistemas Digitais: Princípios e Aplicações. 10. Ed.
Pearson. São Paulo, 2007.\\
WAKERLY, John F. Digital design: principles and practices.
Prentice-Hall. 2000. D'AMORE, Roberto. VHDL: Descrição
e Síntese de Circuitos Digitais. 1. Ed. LTC. 2005.\\
Complementar:\\
1300 Esquemas e Circuitos Eletrônicos. R. Bourgeron. 8528901165. Hermus
(2002).\\
IDOETA, I., CAPUANO, F. G. Elementos de Eletrônica Digital. Érica, 1984\\
MALVINO, Eletrônica. Volume 1, 4. Edição, Makron Books.\\


\newpage IDENTIFICAÇÃO

DISCIPLINA: Inglês CÓDIGO:\\ 
DEPARTAMENTO: Unidade Acadêmica de Garanhuns\\
ÁREA: \\
CARGA HORÁRIA TOTAL: 60 NÚMERO DE CRÉDITOS: 04\\
CARGA HORÁRIA SEMANAL: 04 TEÓRICAS: 03 PRÁTICAS: 01 TOTAL: 04\\
PRÉ-REQUISITOS: não há\\
CO-REQUISITOS: não há

EMENTA 

Leitura de textos em inglês, visando o desenvolvimento de estratégias
globais de leitura e de análise lingüística.

BIBLIOGRAFIA 

Básica:\\
MURPHY, Raymond. Essential grammar in use: a self-study reference and
practice book for elementary students of English. Cambridge: CUP, 1990.\\
GRELLET, François. Developing reading skills. Cambridge: CUP, 1998.\\
PINTO, Dilce et al. Compreensão inteligente de textos. Grasping the
meaning. Vol. 1 e 2, Ao livro técnico, Rio de Janeiro, 1991.\\
Complementar:\\
ESTERAS, Santiago Remacha. Infotech: English for computer users. 4th ed.
Cambridge: Cambridge University Press, 2008.\\
GLEDINNING, Eric \& McEWAN, John. Oxford English for Information
technology. 2nd ed. Oxford: Oxford University Press, 2008.\\


\newpage IDENTIFICAÇÃO

DISCIPLINA: Algoritmos e Estruturas de Dados II CÓDIGO:\\ 
DEPARTAMENTO: Unidade Acadêmica de Garanhuns\\
ÁREA: \\
CARGA HORÁRIA TOTAL: 60 NÚMERO DE CRÉDITOS: 04\\
CARGA HORÁRIA SEMANAL: 04 TEÓRICAS: 02 PRÁTICAS: 02 TOTAL:04\\
PRÉ-REQUISITOS: Algoritmos e Estruturas de Dados I\\
CO-REQUISITOS: não há

EMENTA 

Formas de representação e abstração de dados em memória, Dispositivos e técnicas para armazenamento de dados, Estruturas Abstratas de Dados:  Lista, Pilha, Fila, Métodos de busca e classificação de dados em memória secundária de computadores. Parâmetros físicos e lógicos dos dispositivos para armazenagem de dados. Métodos de representação e abstração de dados Métodos de acesso a arquivos, sequenciais, sequenciais indexadas,  indexadas e diretas. Estruturas árvores-B e  hashing, Árvore binária, Árvore Balanceada, Heap, Grafo.

BIBLIOGRAFIA 

Básica:\\
Estruturas de Dados e Algoritmos em Java. Goodrich, M.T.; Tamassia, R.;
IV Edição, Bookman, São Paulo, 2007. ISBN: 978-85-60031-50-4.\\
KOFFMAN, Elliot B., WOLFGANG, Paul A. T. Data Structures: Abstraction
and Design Using Java. 2010.\\
DALE, Nell., JOYCE, Daniel T., WEEMS, Chip. Object-Oriented Data
Structures Using Java. Third Edition.\\
Complementar:\\
Estruturas de Dados e seus Algoritmos. J. L. Szwarcfiter e L.
Markenzion. Segunda Edição. LTC, 1994.\\
Algoritmos e Estruturas de Dados. A. de M. Guimarães e N. A. de C.
Lages. Editora LTC, 1994.\\
Estruturas de Dados e Algoritmos em C++. A. Drozdek. Pioneira, 2002.

\newpage 

\section*{4o. Período}

IDENTIFICAÇÃO

DISCIPLINA: Arquitetura de Computadores CÓDIGO:\\ 
DEPARTAMENTO: Unidade Acadêmica de Garanhuns\\
ÁREA:\\
CARGA HORÁRIA TOTAL: 60 NÚMERO DE CRÉDITOS: 04\\
CARGA HORÁRIA SEMANAL: 04 TEÓRICAS: 03 PRÁTICAS: 01 TOTAL: 04 \\
PRÉ-REQUISITOS: Sistemas Digitais\\
CO-REQUISITOS: não há

EMENTA 

Organização de Computadores; Conjunto de Instruções, Mecanismos de
Interrupção e de Exceção; Barramento, Comunicações; Interfaces e
Periféricos, Hierarquia de Memória; Multiprocessadores;
Multicomputadores; Arquiteturas Paralelas.

BIBLIOGRAFIA 

Básica:\\
PATTERSON, D; HENNESSY, J. L. Organização e Projeto de Computadores. 3.
Ed. Campus. 2005.\\
STALLINGS, W. Arquitetura e Organização de Computadores. 8. Ed. Pearson.
2010.\\
MONTEIRO, M. A. Introdução à Organização de Computadores. 5. Ed. LTC.
2007.\\
Complementar:\\
HENNESSY, J. L; PATTERSON, D. Arquiteturas de Computadores: Uma
abordagem quantitativa. 3. Ed. Campus. 2003.\\
TOCCI, Ronald J. Sistemas Digitais: Princípios e Aplicações. 10. Ed.
Pearson. São Paulo, 2007.\\
TANENBAUM, A. S. Organização Estruturada de Computadores. 5. Ed. Pearson Ed. 2009.\\


\newpage IDENTIFICAÇÃO

DISCIPLINA: Projeto e Análise de Algoritmos CÓDIGO:\\
DEPARTAMENTO: Unidade Acadêmica de Garanhuns\\
ÁREA: \\
CARGA HORÁRIA TOTAL: 60 NÚMERO DE CRÉDITOS: 04\\
CARGA HORÁRIA SEMANAL: 04 TEÓRICAS: 02 PRÁTICAS: 02 TOTAL:04\\
PRÉ-REQUISITOS: Algoritmos e Estruturas de Dados II\\
CO-REQUISITOS: não há

EMENTA 

Corretude de Algoritmos Recursivos e Iterativos. Notação e Análise
Assintótica. Eficiência de Algoritmos Recursivos e Não-Recursivos.
Análise Empírica. Grafos e Sub-grafos; Isomorfismo, Matrizes de
Adjacência e Incidência, Caminhos e Ciclos. Árvores: Caracterização de
Árvores, Cortes de Arestas, Cortes de Vértices. Conectividade:
Conectividade de Vértices e Arestas; Ciclos Eulerianos e Hamiltonianos.
Emparelhamentos. Coloração de Vértices e de Arestas. Planaridade. Redes
de Fluxo. Algoritmos Gulosos. Programação Dinâmica. Backtracking.
Branch-and-Bound. NP-Completude.

BIBLIOGRAFIA 

Básica:\\
CORMEN, T. H., LEISERSON, C. E., RIVEST, R. L., STEIN, C. Algoritmos:
Teoria e prática. Editora Campus, tradução da 2. edição Americana,
2002.\\
GERSTING, Judith L. Fundamentos Matemáticos para a Ciência da
Computação. LTC  Livros Técnicos e Científicos, 1982.\\
PAPADIMITRIOU, C. H., VAZIRANI, U. V., DASGUPTA, S. Algoritmos.
McGraw-Hill, 2006.\\
Complementar:\\
SIPSER, Michael. Introdução à Teoria da Computação. Thomson Pioneira, 2.
edição, 2007.\\
MEDINA, M., FERTIG, C. Algoritmos e Programação: Teoria e Prática.
Novatec, 2005.\\
HOPCROFT, J. E.; MOTWANI, R.; ULLMAN, J.D.: Introdução a Teoria de
Autômatos, Linguagens e Computação. Rio de Janeiro: Campus, 2. edição, 2002.\\


\newpage IDENTIFICAÇÃO

DISCIPLINA: Paradigmas de Linguagens de Programação CÓDIGO:\\ 
DEPARTAMENTO: Unidade Acadêmica de Garanhuns\\
ÁREA: \\
CARGA HORÁRIA TOTAL: 60 NÚMERO DE CRÉDITOS: 04\\
CARGA HORÁRIA SEMANAL: 04 TEÓRICAS: 02 PRÁTICAS: 02 TOTAL: 04\\
PRÉ-REQUISITOS: não há\\
CO-REQUISITOS: não há

EMENTA 

Valores e Tipos; Variáveis e Armazenamento; Associações e Escopo;
Abstração de Procedimentos; Abstração de Dados; Abstração Genérica;
Sistema de tipos; Controle de Fluxo; Concorrência; Paradigma
Imperativo; Paradigma Orientado a Objetos; Paradigma Concorrente;
Paradigma Funcional; Paradigma Lógico; Paradigma Scripting.

BIBLIOGRAFIA 

Básica:\\
David A. Watt. Programming Language Concepts and Paradigms. 2004.\\
R. Sebesta. Conceitos de Linguagens de Programação  4a Ed., Bookman,
2000.\\
Varejão, Flávio. Linguagens de Programação  Conceitos e Técnicas.
Editora Campus, 2004.\\
Complementar:\\
C. Ghezzi \& M. Jazayeri Programming Language Concepts-3rd Ed., John
Wiley\&Sons, 1997.\\
R. Sethi Programming Languages: Concepts and Languages-2nd Ed., Addison
Wesley. 1996.\\
Pratt, T.W. Programming Languages  Design and Implementation, Third
Edition, Prentice-Hall, 1996.\\


\newpage IDENTIFICAÇÃO

DISCIPLINA: Banco de Dados CÓDIGO:\\ 
DEPARTAMENTO: Unidade Acadêmica de Garanhuns\\
ÁREA: \\
CARGA HORÁRIA TOTAL: 60 NÚMERO DE CRÉDITOS: 04\\
CARGA HORÁRIA SEMANAL: 04 TEÓRICAS: 02 PRÁTICAS: 02 TOTAL: 04\\
PRÉ-REQUISITOS: Algoritmos e Estruturas de Dados II\\
CO-REQUISITOS: não há

EMENTA 

Conceituação. Arquitetura de SGDB. Modelagem de dados: modelo E-R e suas
variações, abstrações por agregação e generalização. Modelos de
representação (relacional, hierárquico e redes). Normalização e
manutenção da integridade. Arquitetura de Sistemas de Bancos de Dados,
SQL.

BIBLIOGRAFIA 

Básica: \\
SILBERSCHATZ, Abraham et al. ; Sistemas de bancos de dados. 3. ed. 778
p. São Paulo : Makron Books, 1997\\
DATE, C. J. ; Banco de dados. Rio de Janeiro : Campus, 1990\\
Complementar:\\
MOLINA Hector Garcia, Ullman, Jefrey  IMPLEMENTAÇÃO DE SISTEMAS DE
BANCO DE DADOS  ED. CAMPUS LONEY, Kevin et al.; Oracle  Sistema
de computador; Banco de dados ralacional -- Oracle Rio de Janeiro:
Campus, 2000\\


\newpage IDENTIFICAÇÃO

DISCIPLINA: Engenharia de Software CÓDIGO:\\ 
DEPARTAMENTO: Unidade Acadêmica de Garanhuns\\
ÁREA: \\
CARGA HORÁRIA TOTAL: 60 NÚMERO DE CRÉDITOS: 04\\
CARGA HORÁRIA SEMANAL: 04 TEÓRICAS: 02 PRÁTICAS: 02 TOTAL: 04\\
PRÉ-REQUISITOS: Programação Orientada a Objetos\\
CO-REQUISITOS: não há

EMENTA 

Contextualização da Engenharia de Software; Fundamentação dos Princípios
da Engenharia de Software; Conceituação de Produto e Processo de
Software; Tipos de Processos de Software; Comparação entre os
Paradigmas de Desenvolvimento Software; Caracterização do Projeto de
Software; UML; Gerenciamento de Projetos; Gerenciamento Ágil; Processo
de Engenharia de Requisito; Requisitos; Testes de Software; Estilos
Arquiteturais; Evolução e Refatoração; Definição de Qualidade de
Software.

BIBLIOGRAFIA 

Básica:\\
Sommerville, Ian. Engenharia de Software. ISBN.: 8588639076. Addison 
Wesley.\\
Roger S. Pressman. Engenharia de Software. ISBN.: 8586804576.
McGraw-Hill.\\
FLEEGER, Shari Lawrence. Software Engineering : Theory and Practice,
Englewood Cliffs, N. J.:Prentice Hall, 2001, 659p. ISBN 0130290491.\\
Complementar:\\
Steve McConnell. Code Complete: A Practical Handbook of Software
Construction. ISBN.: 0735619670. Ed. Microsoft Pres.\\
Freeman E. , Freeman E., Bates B., Sierra K.; Head First Design
Patterns. ISBN.: 0596007124. Ed. O'Reilly Media.\
\newpage 

\section*{5o. Período}

IDENTIFICAÇÃO

DISCIPLINA: Sistemas Operacionais CÓDIGO:\\ 
DEPARTAMENTO: Unidade Acadêmica de Garanhuns\\
ÁREA: \\
CARGA HORÁRIA TOTAL: 60 NÚMERO DE CRÉDITOS: 04\\
CARGA HORÁRIA SEMANAL: 04 TEÓRICAS: 03 PRÁTICAS: 01 TOTAL: 04\\
PRÉ-REQUISITOS: Arquitetura de Computadores 
CO-REQUISITOS: não há

EMENTA 

Estrutura de um Sistema Operacional. Processos Concorrentes.
Escalonamento. Gerenciamento de Memória. Memória Virtual. Gerenciamento
de Disco. Sistemas de Arquivos. Proteção e Segurança. Sistemas
Distribuídos. Estudos de Caso.

BIBLIOGRAFIA 

Básico:\\
TANENBAUM, A. Sistemas Operacionais Modernos. Prentice-Hall. 2009.\\
DEITEL; CHOFFNES. Sistemas Operacionais. Pearson Education. 3a Ed.
2005.\\
STALLINGS, W. Operating Systems: Internals and Design Principles.7a Ed.
2011.\\
Complementar:\\
WOODHULL, A. S.; TANEMBAUM, A. Sistemas Operacionais: Projeto e
Implementação. Artmed. 2008.\\
SILBERSCHATZ, A; GALVIN, P.; GAGNE, G. Operating Systems Concepts. 6a
Ed. John Willey. 2003.\\
OLIVEIRA, R. S., CARISSIMI, A. S., TOSCANI, S. Sistemas Operacionais. 2.
Ed. Editora Sagra Luzzato: Porto Alegre. 2001.\\


\newpage IDENTIFICAÇÃO

DISCIPLINA: Sistemas de Informação e Tecnologias CÓDIGO:\\ 
DEPARTAMENTO: Unidade Acadêmica de Garanhuns\\
ÁREA:\\
CARGA HORÁRIA TOTAL: 60 NÚMERO DE CRÉDITOS: 04\\
CARGA HORÁRIA SEMANAL: 04 TEÓRICAS: 04 PRÁTICAS: 00 TOTAL: 04\\
PRÉ-REQUISITOS: não há\\
CO-REQUISITOS: não há

EMENTA 

A organização. Configuração estrutural. A TI na empresa e a Revolução da
Web.Visão sistêmica de estratégias integradoras de áreas e informação
como apoio ao processo decisório. Aplicações organizacionais;
Planejamento. Elementos da Tomada de decisão numa organização. Decisão
e controle. Sistemas de Informação Transacionais, Gerenciais e de Apoio
às Operações e à Decisão. ERPs. CRMs. SCMs. Business intellligence.
Gestão do conhecimento. A importância do planejamento em TI. Tendências
em TI nas organizações.

BIBLIOGRAFIA 

Básica:\\
TURBAN, E.; MCLEAN, E; WETHERBE, J. Tecnologia da Informação para
Gestão. Transformando os Negócios na Economia Digital. Tradução de
Renate Schinke. Revisão técnica de Ângela F. Brodbeck. Porto Alegre:
Bookman, 2004.\\
O'BRIEN, James. Sistemas de Informação e as decisões
gerenciais na era da Internet. São Paulo: Editora Saraiva, 2001.\\
Rezende, D. A E. Abreu, A F de . Tecnologia Da Informação Aplicada A
Sistemas de Informação Empresariais: O Papel Estratégico Da Informação
E Dos Sistemas de Informação Nas Empresas. São Paulo: Editora Atlas, 2010.\\
Complementar:\\
GORDON, S.R.; GORDON, J.R. Sistema de informação. 3. ed. Rio de Janeiro:
LTC. Livros Técnicos e Científicos Editora S.A., 2006.\\


\newpage IDENTIFICAÇÃO

DISCIPLINA: Teoria da Computação CÓDIGO:\\
DEPARTAMENTO: Unidade Acadêmica de Garanhuns\\
ÁREA: \\
CARGA HORÁRIA TOTAL: 60 NÚMERO DE CRÉDITOS: 04\\
CARGA HORÁRIA SEMANAL: 04 TEÓRICAS: 04 PRÁTICAS: 00 TOTAL: 04\\
PRÉ-REQUISITOS: Matemática Discreta\\
CO-REQUISITOS: não há

EMENTA 

Conceitos Básicos. Alfabetos e Linguagens. Gramáticas. Autômatos
Finitos; Linguagens Regulares. Linguagens Livres de Contexto; Autômato
de Pilhas; Gramáticas Livre de Contexto Ambígua. Formas Normais.
Linguagens Recursivamente Enumeráveis e Sensíveis ao Contexto.
Hierarquia de Chomsky. Indecidibilidade.

BIBLIOGRAFIA 

Básica:\\
MENEZES, Paulo F B: Linguagens Formais e Autômatos. P. Alegre: Sagra
Luzzatto, 2005 (5a. edição).\\
HOPCROFT, J. E.; MOTWANI, R.; ULLMAN, J.D.: Introdução à Teoria de
Autômatos, Linguagens e Computação. Rio de Janeiro: Elsevier, 2002.\\
Complementar:\\
GERSTING, J. L: Fundamentos Matemáticos para a Ciência da Computação.
Rio de Janeiro: LTC, 2008 (5.. edição).\\

\newpage IDENTIFICAÇÃO

DISCIPLINA: Redes de Computadores CÓDIGO:\\ 
DEPARTAMENTO: Unidade Acadêmica de Garanhuns\\
ÁREA: \\
CARGA HORÁRIA TOTAL: 60 NÚMERO DE CRÉDITOS: 04\\
CARGA HORÁRIA SEMANAL: 04 TEÓRICAS: 04 PRÁTICAS: 00 TOTAL: 04\\
PRÉ-REQUISITOS: Introdução a Computação\\
CO-REQUISITOS: não há

EMENTA 

Conceitos básicos de redes; Modelo OSI; Modelo TCP/IP; Conceitos básicos
de roteamento; Conceitos básicos de switching e roteamento
intermediário; e Tecnologias WAN.

BIBLIOGRAFIA 

Básica:\\
Kurose, James F.; Ross, Keith W. Redes de computadores e a internet: uma
abordagem top-down. 3. Edição. São Paulo: Pearson Addison Wesley, 2008.\\
Tanenbaum, Andrew S. Redes de computadores. 4. Edição. Rio de Janeiro:
Editora Campus, 2002.\\
Peterson,Larry. Redes de computadores: Uma abordagem de Sistemas. 3.
Edição. Editora Campos, 2004.\\
Complementar:\\
COMER, Douglas E. Interligação de Redes com TCP/IP, Volumes I 5. Edição.
Prentice Hall, 2006.\\
COMER, Douglas E. Interligação de Redes com TCP/IP, Volumes II 5.
Edição. Prentice Hall, 2006.\\
SOARES, LEMOS e COLCHER -- Redes Locais -- Das LANs, MANs e WANs às
redes ATM. 2. Edição. Ed. Campus, 1995.\\

\newpage IDENTIFICAÇÃO

DISCIPLINA: Inteligência Artificial CÓDIGO:\\ 
DEPARTAMENTO: Unidade Acadêmica de Garanhuns\\
ÁREA: \\
CARGA HORÁRIA TOTAL: 60 HORAS NÚMERO DE CRÉDITOS: 04\\
CARGA HORÁRIA SEMANAL: 04 TEÓRICAS: 02 PRÁTICAS: 02 TOTAL: 04\\
PRÉ-REQUISITOS: Projeto e Análise de Algoritmos; Paradigmas de Linguagens de Programação\\
CO-REQUISITOS: não há

EMENTA 

Introdução. Sistemas especialistas. Agentes Inteligentes. Resolução de
problemas por meio de busca. Problema de satisfação de restrição.
Linguagens Simbólicas. Esquemas para representação do conhecimento:
lógicos, em rede, estruturados, procedurais. Formalismos para a
representação de conhecimento incerto. Redes Bayesianas. Conjuntos e
lógica Fuzzy. Introdução à Computação Evolucionária. Algoritmos
Genéticos. Ajuste de parâmetros em algoritmos genéticos. Projeto.

BIBLIOGRAFIA 

Básica:\\
RUSSELL, Stuart e NORVIG, Peter. Inteligência Artificial. Campus, Rio de
Janeiro, 2004.\\
REZENDE, Solange O. (2003). Sistemas Inteligentes -- Fundamentos e
aplicações. Barueri, SP. Editora Manole. 2003.\\
EIBEN, Agoston E. \& Smith, J.E. Introduction to Evolutionary Computing.
Springer. 2008.\\
Complementar:\\
LUGER, George. Inteligência Artificial: Estruturas e Estratégias para a
Solução. Bookman, Porto Alegre, 2004.\\
LIDEN, Ricardo. Algoritmos Genéticos. Brasport. 2006\\
RICH, Elaine; KNIGHT, Kevin. Inteligência artificial. 2. ed. Makron
Books, São Paulo, 1994.\\
BITTENCOURT, G: Inteligência Artificial -- Ferramentas e Teorias.
Editora da UFSC. 2.. Edição. Florianópolis, 2001.
\newpage 

\section*{6o. Período}

IDENTIFICAÇÃO

DISCIPLINA: Empreendimentos em TIC CÓDIGO:\\ 
DEPARTAMENTO: Unidade Acadêmica de Garanhuns\\
ÁREA:\\
CARGA HORÁRIA TOTAL: 60 NÚMERO DE CRÉDITOS: 04\\
CARGA HORÁRIA SEMANAL: 04 TEÓRICAS: 04 PRÁTICAS: 00 TOTAL: 04\\
PRÉ-REQUISITOS: não há\\
CO-REQUISITOS: não há

EMENTA 

Importância e Contribuição Econômica e Social do Empreendedorismo;
Empreendedores: Características e comportamentos empreendedores.
Motivações; Inovação tecnológica e empreendedorismo; O sucesso e o
fracasso de novos empreendimentos; Cases de sucesso em computação.
Plano de negócio: Importância, estruturação e apresentação. Pesquisa de
Mercado; Ferramentas Gerenciais para o empreendedor. Incubadoras: o que
são, objetivos; Caminhos a seguir e recursos disponíveis para o
empreendedor: Incubação, MCT, FINEP, Venture Capital, Start-ups.

BIBLIOGRAFIA 

Básica:\\
DRUCKER, P. F. Inovação e Espírito Empreendedor. São Paulo: Pioneira,
1986.\\
DOLABELA, Fernando. Oficina do empreendedor. São Paulo: Cultura, 1999.\\
Leite, Emanuel (2002) -- O Fenómeno do Empreendedorismo, Recife, Edições
Bagaço.\\
Complementar:\\
DEGEN, Ronald Jean. O Empreendedor: empreender como opção de carreira. São Paulo: Pearson Prentice Hall, 2009.\\
Sites Específicos:\\
http://www.planodenegocios.com.br - portal brasileiro de plano de negócios, com artigos, cursos e eventos, links e bibligrafia sobre o tema;\\
http://www.empreendedor.com.br - revista para empreendedores;\\
http://www.anprotec.org.br - ANPROTEC;\\
http://emanueleite.blogspot.com/\\


\newpage IDENTIFICAÇÃO

DISCIPLINA: Reconhecimento de Padrões CÓDIGO:\\ 
DEPARTAMENTO: Unidade Acadêmica de Garanhuns\\
ÁREA: \\
CARGA HORÁRIA TOTAL: 60 NÚMERO DE CRÉDITOS: 04\\
CARGA HORÁRIA SEMANAL: 04 TEÓRICAS: 04 PRÁTICAS: 00 TOTAL: 04\\
PRÉ-REQUISITOS: Algoritmos e Estruturas de Dados II, Matemática
Discreta.\\
CO-REQUISITOS: Computação Gráfica.

EMENTA 

Introdução a Aprendizagem de Máquina. Paradigmas de Aprendizagem de
Máquina. Classificação e avaliação de classificadores. Regressão.
Agrupamento. Redes Neurais. Introdução a Processamento de Imagens.
Realce, filtragem e restauração de imagens. Segmentação de imagens.
Compressão e comunicação de imagens. Noções de visão computacional e
reconhecimento de padrões. Projeto.

BIBLIOGRAFIA 

Básica:\\
TAN, Pang-Ning; STEINBACH; Michael; KUMAR, Vipin. Introdução ao Data
Mining. Ciência Moderna. 2009.\\
HAYKIN, Simon. Redes Neurais: Princípios e Prática, 2. edição. Bookman.
2001.\\
GONZALEZ, Rafael C.; WOODS, Richard E.. Processamento Digital De Imagens
(3 Ed). Pearson Education. 2010.\\
Complementar:\\
WITTEN, I. H., Frank, E. (2005). Data mining : practical machine
learning tools and techniques. Second Edition. Elsevier.\\
BEALE, R.; Jackson, T. Neural Computing  An Introduction. Institute
of Physics Publishing. 1990.\\
MITCHELL, T. (1997). Machine Learning. McGraw-Hill.\\
THEODORIDIS. S., Koutroumbas, K. (2009). Pattern Recognition. Fourth
Edition, Academic Press.\\
REZENDE, Solange O. (2003). Sistemas Inteligentes -- Fundamentos e
aplicações. Barueri, SP. Editora Manole. 2003.\\
BRAGA, A. P., Ludermir, T.B. e de Carvalho A. P. L. F. (2007). A Redes
Neurais Artificiais -- Teoria e Aplicações. LTC.\\
MARQUES, Jorge Salvador. (2005). Reconhecimento de Padrões  Métodos
Estatísticos e Neurais. IST Press.\\
DUDA, R.O., Hart, P.E., Stork, D.G. (2001) Pattern Classification.
Second Edition. Wiley.\\
BISHOP, C. M. (2006). Pattern Recognition and Machine Learning.
Springer.\\
A. K. Jain. Fundamentals of digital image processing. Prentice-Hall
International Editions, Englewood Cliffs, NJ, 1989.\\
Bernd Jähne. Digital Image Processing 6th edition. Springer. 2005.\\


\newpage IDENTIFICAÇÃO

DISCIPLINA: Compiladores CÓDIGO:\\ 
DEPARTAMENTO: Unidade Acadêmica de Garanhuns\\
ÁREA: \\
CARGA HORÁRIA TOTAL: 60 NÚMERO DE CRÉDITOS: 04\\
CARGA HORÁRIA SEMANAL: 04 TEÓRICAS: 02 PRÁTICAS: 02 TOTAL: 04\\
PRÉ-REQUISITOS: Teoria da Computação; Algoritmos e Estruturas de Dados II\\
CO-REQUISITOS: não há

EMENTA 

Processadores de linguagem: compilador e interpretador. Introdução à
compilação. Fases da compilação. Relações sobre gramáticas. Análise
léxica. Análise sintática ascendente e descendente. Tabelas de
símbolos. Esquemas de tradução. Análise semântica. Geração de código
intermediário. Ambientes de execução. Introdução à otimização de
código.

BIBLIOGRAFIA 

Básica:\\
AHO, A. V., SETHI, R., ULLMAN, J. D. Compiladores: Princípios, Técnicas
e Ferramentas. Rio de Janeiro: Guanabara Koogan, 1986.\\
PRICE, A. M. Implementação de Linguagens de Programação: Compiladores.
Sagra Luzzato, 2001.\\
Complementar:\\
DELAMARO, M. E. Como Construir um Compilador -- Utilizando Ferramentas
Java. São Paulo: Novatec, 1. edição, 2004.\\
SEBESTA, R. W. Conceitos de Linguagens de Programação. Porto Alegre:
Bookman, 5. edição, 2003.\\
MENEZES, P. B. Linguagens Formais e Autômatos. Sagra Luzzato, 5. edição,
2005.\\
SETZER, Valdemar. A Construção de um Compilador. Campus, Rio de Janeiro, 1986.

\newpage IDENTIFICAÇÃO

DISCIPLINA: Computação Gráfica CÓDIGO:\\ 
DEPARTAMENTO: Unidade Acadêmica de Garanhuns\\
ÁREA: \\
CARGA HORÁRIA TOTAL: 60 NÚMERO DE CRÉDITOS: 04\\
CARGA HORÁRIA SEMANAL: 04 TEÓRICAS: 04 PRÁTICAS: 00 TOTAL: 04\\
PRÉ-REQUISITOS: Introdução à Programação; Álgebra Linear\\
CO-REQUISITOS: não há

EMENTA 

Introdução à computação gráfica. Biblioteca gráfica OpenGL.
Representação de objetos. Dispositivos periféricos gráficos. Processo
de visualização. Curvas e superfícies paramétricas. Eliminação de
superfícies ocultas. Geração de imagens com realismo. Tópicos
complementares em computação gráfica.

BIBLIOGRAFIA 

Básica:\\
Computação Gráfica: Teoria e Prática, AZEVEDO, E. e Conci, A. Editora
Campus, Elsevier, 2003. Rio de Janeiro.ISBN 8535212533.\\
Fundamentos de Computação Gráfica, Gomes, J. e Velho, L. IMPA, 2003.\\
Geometric Algebra for Computer Graphics, John A. Vince, Springer, 2008.
ISBN 1846289963.\\
Complementar:\\
Fundamentals of Computer Graphics, Second Ed. Peter Shirley, et al. A K
Peters Ltd, 2005. ISBN 1568812698.\\
OpenGL(R) Programming Guide, Shreiner D., et al. Addison-Wesley, 5th
Edition, 2005. ISBN 0321335732.\\
Foley, J.D. van Dam, A. Feiner K.S., Jughes, J.F.,
 Computer Graphics: Principles And
Practice , Addison Wesley, 1993.\\


\newpage IDENTIFICAÇÃO

DISCIPLINA: Sistemas Distribuídos CÓDIGO:\\ 
DEPARTAMENTO: Unidade Acadêmica de Garanhuns\\
ÁREA: \\
CARGA HORÁRIA TOTAL: 60 NÚMERO DE CRÉDITOS: 04\\
CARGA HORÁRIA SEMANAL: 04 TEÓRICAS: 04 PRÁTICAS: 00 TOTAL: 04\\
PRÉ-REQUISITOS: Redes de Computadores\\
CO-REQUISITOS: não há

EMENTA 

Introdução aos sistemas distribuídos; Definições de processos e threads;
Comunicação em sistemas distribuídos; Sincronização em sistemas
distribuídos; Conceitos de middleware;Redes P2P: conceitos básicos,
arquiteturas, aplicações;Introdução a grades computacionais;Tecnologias
de middleware tradicionais.

BIBLIOGRAFIA 

Básica:\\
G. Coulouris, J. Dollimore e T. Kindberg,  Sistemas
Distribuídos: Conceitos e Projetos . 4. Edição.
Bookman Companhia, 2007.\\
A.S. Tanenbaum and M.V. Steen,  Sistemas Distribuídos:
Princípios e Paradigmas . 2. Edição. Prentice Hall
Brasil, 2007.\\
M. DANTAS. Computação Distribuída: Redes, Grids e Clusters
Computacionais. Rio de janeiro: Axcel Books, 2005.
Complementar:\\
U. RIBEIRO. Sitemas Distribuídos: Desenvolvendo Aplicações de Alta
Perfarmance no Linux. Rio de Janeiro: Axcel Books, 2005.\\
Marques, José Alves. Tecnologia de Sistemas Distribuídos. 1. Edição.
FCA, 1998.\\
Tanenbaum, Andrew S. Redes de computadores. 4. Edição. Rio de Janeiro:
Editora Campus, 2002.
\newpage 

\section*{7o. Período}

IDENTIFICAÇÃO

DISCIPLINA: Projeto de Desenvolvimento de Software CÓDIGO:\\ 
DEPARTAMENTO: Unidade Acadêmica de Garanhuns\\ 
ÁREA: \\
CARGA HORÁRIA TOTAL: 90 NÚMERO DE CRÉDITOS: 06\\
CARGA HORÁRIA SEMANAL: 06 TEÓRICAS: 02 PRÁTICAS: 05 TOTAL: 06\\
PRÉ-REQUISITOS: Engenharia de Software; Fundamentos de Banco de Dados\\
CO-REQUISITOS: não há

EMENTA 

Desenvolvimento de um sistema de computação usando conceitos aprendidos
anteriormente. Sistemas multidisciplinares devem ser estimulados bem
como o trabalho em equipe.

BIBLIOGRAFIA 

Básica:\\
PRESMANN. Engenharia de Software 7. ed. 776 p. McGraw-Hill, 2011\\
SOMMERVILLE. Engenharia de Software. 8. ed. 568 p. Addison Wesley, 2007\\
BRAUDE. Projeto de Software. Bookman. 2005.\\
Complementar:\\
WAZLAWICK. Análise e Projetos de Sistemas de Informação. 2. Ed. 344p.
Campus. 2010.\\
BEZERRA. Princípios de Análise e Projetos de Sistemas UML. 380p, Campus
2006.\\
HELM et. al. Padrões de Projeto 1. Ed. 366p. Bookman. 2005.\\


\newpage IDENTIFICAÇÃO

DISCIPLINA: Interface Humano-Computador CÓDIGO:\\
DEPARTAMENTO: Unidade Acadêmica de Garanhuns\\
ÁREA: \\
CARGA HORÁRIA TOTAL: 60 NÚMERO DE CRÉDITOS: 04\\
CARGA HORÁRIA SEMANAL: 04 TEÓRICAS: 02 PRÁTICAS: 02 TOTAL: 04\\
PRÉ-REQUISITOS: Engenharia de Software
CO-REQUISITOS: não há

EMENTA 

Conhecer os fundamentos de fatores humanos em IHC, os modelos mentais
(metáforas), os paradigmas de IHC (engenharia semiótica e cognitiva) e
os métodos, técnicas, suporte e avaliação de design de interação.

BIBLIOGRAFIA 

Básica:\\
Barbosa, S. D. J.; Silva, B. S. Interação Humano-Computador. Rio de
Janeiro: Elservier, 2010.\\
Preece, Rogers e Sharp; Design de Interação  Além da Interação
Homem-computador. São Paulo: Bookman, 2005.\\
Complementar:\\
de Souza, C.S.; Leite, J.C.; Prates, R.O.; Barbosa, S.D.J. Projeto de
Interfaces de Usuário: perspectivas cognitivas e semióticas. Jornada de
Atualização em Informática (JAI), Congresso da SBC, 1999.\\
Prates, R.O.; Barbosa, S.D.J. Avaliação de Interfaces de Usuário 
Conceitos e Métodos. Jornada de Atualização em Informática (JAI), XXIII
Congresso da SBC, 2003.

\newpage IDENTIFICAÇÃO

DISCIPLINA: Computadores e Sociedade CÓDIGO:\\ 
DEPARTAMENTO: Unidade Acadêmica de Garanhuns\\
ÁREA: \\
CARGA HORÁRIA TOTAL: 30 NÚMERO DE CRÉDITOS: 02\\
CARGA HORÁRIA SEMANAL: 02 TEÓRICAS: 02 PRÁTICAS: 00 TOTAL: 02\\
PRÉ-REQUISITOS: não há\\
CO-REQUISITOS: não há

EMENTA 

CONSEQÜÊNCIAS DA INFORMATIZAÇÃO DA SOCIEDADE: A informatização e o
aspecto educacional; Efeitos políticos e econômicos; Impactos sociais;
Informatização e privacidade; POLÍTICA NACIONAL DE INFORMÁTICA:
Indústria nacional de informática; O papel do analista e sistemas na
sociedade; AUTOMAÇÃO DE ATIVIDADES: Comerciais; Industriais;
Escritórios; APLICAÇÕES DA INFORMÁTICA: Científica; Administrativa;
Jurídica; Humanística; Educação; ERGONOMIA E DOENÇAS PROFISSIONAIS:
Tipos; Características.

BIBLIOGRAFIA 

Básica:\\
MINISTÉRIO da Ciência e Tecnologia. Sociedade da Informação 
no Brasil  Livro Verde. Brasília: Imprensa Nacional,
2000.Bibliografia
Complementar:\\
MASIEIRO, Paulo C. Ética em Computação. São Paulo : Editora da
Universidade de São Paulo. 2000.\\
YOUSSEF, Antonio N.; FERNANDES, Vicente P. Informática e Sociedade. 2
ed. São Paulo: Ática, 1998.\\
NEGROPONTE, Nicholas. A Vida Digital. São Paulo: Companhia das Letras,
1998.\\
SCHAFF, Adam. A Sociedade Informática. São Paulo: Brasiliense, 1995.\\
NAVEIRA \& SILVA, L. A quarta onda. Rio de Janeiro : Record, 1989.\\
TOFFLER, Alvin. A Terceira Onda. Rio de Janeiro, Record, 1995.\\
SOUSA, M. S. L. Introdução aos aspectos humanos da 
interação homem: computador. UFRGS, 1995. \\
REZENDE, D. A. Engenharia de software e sistemas de informação. Rio de Janeiro: Brasport, 2002,
358p. ISBN 8574520942. ISBN 8522103909. (2004).


%--------------------------------------------------------
\chapter{Ementas das Disciplinas Optativas}
%--------------------------------------------------------

\section*{Área Temática: Banco de Dados}

IDENTIFICAÇÃO

DISCIPLINA: Banco de Dados Avançado CÓDIGO:\\
DEPARTAMENTO: Unidade Acadêmica de Garanhuns\\
ÁREA: \\
CARGA HORÁRIA TOTAL: 60 NÚMERO DE CRÉDITOS: 04\\
CARGA HORÁRIA SEMANAL: 04 TEÓRICAS: 02 PRÁTICAS: 02 TOTAL: 04\\
PRÉ-REQUISITOS: Banco de Dados\\
CO-REQUISITOS: não há

EMENTA 

Conceitos avançados de Banco de Dados. SQL Avançado. Aspectos
Operacionais de SGBD (Controle de Concorrência, Restrições de
Integridade, Segurança e Recuperação após falhas). Modelo EER (ER
Estendido). Data Warehouse. Data Mining. BDs Móveis.BDs Multimídia. BDs Geográficos. BDs Distribuídos. BD Orientado a
Objetos. BD Objeto-Relacional. BDs Ativos. BDs Temporais. BDs
Biológicos. BDs P2P. BDs Autônomos. Cloud Data Bases (Banco de Dados em
Nuvens).

BIBLIOGRAFIA 

Básica:\\
SILBERSCHATZ, Abraham et al. Sistemas de bancos de dados. 5. ed. 778 p.
São Paulo : Makron Books, 2006\\
ÖZSU and VALDURIEZ. Princípios de Sistemas de Bancos de dados
Distribuídos, Editora Campus, 2001\\
ELSMARI \& NAVATHE. Sistemas de Banco de Dados. 6. Ed. 744p. Pearson.
2011.\\
Complementar:\\
REUSER, Carlos Alberto. Projeto de Banco de Dados. 6. Ed. 282p. Ed.
Bookman 2009.\\
RAMAKRISHNAN, GEHRKE. Sistema de Gerenciamento de Banco de Dados. 3. Ed.
884 páginas, Mcgraw Hill 2008.\\
RIGAUX et. al. Spatial Databases: with application to GIS, Morgan
Kaufmann, 2002.\

\newpage IDENTIFICAÇÃO

DISCIPLINA: Integração de Dados e Data Warehouse CÓDIGO:\\
DEPARTAMENTO: Unidade Acadêmica de Garanhuns\\
ÁREA: \\
CARGA HORÁRIA TOTAL: 60 NÚMERO DE CRÉDITOS: 04\\
CARGA HORÁRIA SEMANAL: 04 TEÓRICAS: 02 PRÁTICAS: 02 TOTAL: 04\\
PRÉ-REQUISITOS: Banco de Dados\\
CO-REQUISITOS: não há

EMENTA 

Definir o processo de construção de um Data Warehouse (DW). Explicar as
diferentes fases de construção de um DW. Listar os principais fatores
que definem um projeto com sucesso. Analisar e transformar exigências
empresariais em um modelo de negócios (conceitual). Utilizar diagramas
de relacionamentos de entidades para transformar o modelo de negócios
em um modelo dimensional (lógico). Transformar o modelo dimensional em
um projeto de dados físico. Apresentar as principais estruturas que
cooperam no desempenho e criação de uma base DW.

BIBLIOGRAFIA 

Básica:\\
SILBERSCHATZ et al. Sistemas de bancos de dados. 5. ed. 778 p. São Paulo
: Makron Books, 2006\\
ÖZSU \& VALDURIEZ. Princípios de Sistemas de Bancos de dados
Distribuídos, Editora Campus, 2001\\
JARKE et, al. Fundamentals of Data Warehouses, Springer-Verlag. 2000.\\
Complementar:\\
HAN \& KAMBER. Data mining: concepts and techniques. Morgan-Kauffman.
2000.\\
ABITEBOUL et. al. Data on the Web, from relations to semistructured data
and XML. Morgan-Kauffman. 2000.\\
ELSMARI \& NAVATHE. Sistemas de Banco de Dados. 6. Ed. 744p. Pearson.
2011.

\newpage IDENTIFICAÇÃO

DISCIPLINA: Modelagem Conceitual de Dados CÓDIGO:\\
DEPARTAMENTO: Unidade Acadêmica de Garanhuns\\
ÁREA: \\
CARGA HORÁRIA TOTAL: 60 NÚMERO DE CRÉDITOS: 04\\
CARGA HORÁRIA SEMANAL: 04 TEÓRICAS: 02 PRÁTICAS: 02 TOTAL: 04\\
PRÉ-REQUISITOS: Banco de Dados
CO-REQUISITOS: não há

EMENTA 

Conceitos Básicos. Análise de Requisitos para Projeto Conceitual do
Banco de Dados. Verificação do Projeto Conceitual do Banco de Dados.
Estratégias para Especificação do Projeto Conceitual de Banco de Dados.
Aspectos Avançados de Projeto Conceitual de Banco de Dados com o Modelo
Entidade-Relacionamento (MER) e com a Linguagem de Modelagem Unificada
(UML). Metamodelos e UML Profile. Ferramentas CASE. Ontologias e
Modelos Conceituais de Banco de dados. Projeto Conceitual de Data
Warehouse. Tópicos Especiais. Projeto Prático.

BIBLIOGRAFIA 

Básica:\\
SILBERSCHATZ, Abraham et al. Sistemas de bancos de dados. 5. ed. 778 p.
São Paulo : Makron Books, 2006\\
TEOREY et. al.  Projeto e Modelagem de Banco De Dados, Elsevier.
2007.\\
ELSMARI \& NAVATHE. Sistemas de Banco de Dados. 6. Ed. 744p. Pearson.
2011.\\
Complementar:\\
REUSER, Carlos Alberto. Projeto de Banco de Dados. 6. Ed. 282p. Ed.
Bookman 2009.\\
RAMAKRISHNAN, GEHRKE. Sistema de Gerenciamento de Banco de Dados. 3. Ed.
884 páginas, Mcgraw Hill 2008.\\
DATE, C. J. Introdução a Sistemas de Banco de dados. 8. Ed. 900 p. Rio de Janeiro : Campus, 2004\\

\newpage IDENTIFICAÇÃO

DISCIPLINA: Projeto de Banco de Dados CÓDIGO:\\
DEPARTAMENTO: Unidade Acadêmica de Garanhuns\\
ÁREA: \\
CARGA HORÁRIA TOTAL: 60 NÚMERO DE CRÉDITOS: 04\\
CARGA HORÁRIA SEMANAL: 04 TEÓRICAS: 02 PRÁTICAS: 02 TOTAL: 04\\
PRÉ-REQUISITOS: Banco de Dados\\
CO-REQUISITOS: não há

EMENTA 

Realizar um projeto real com aplicação de mercado utilizando os
conceitos de Banco de Dados já adquirido e um conjunto de novos
conceitos apresentados nesta disciplina.

BIBLIOGRAFIA 

Básica:\\
SILBERSCHATZ, Abraham et al. Sistemas de bancos de dados. 5. ed.
778 p. São Paulo : Makron Books, 2006\\
REUSER, Carlos Alberto. Projeto de Banco de Dados. 6. Ed. 282p.
Ed. Bookman 2009.\\
ELSMARI \& NAVATHE. Sistemas de Banco de Dados. 6. Ed. 744p.
Pearson. 2011.\\
Complementar:\\
RAMAKRISHNAN, GEHRKE. Sistema de Gerenciamento de Banco de Dados. 3. Ed.
884 páginas, Mcgraw Hill 2008.\\
DATE, C. J. Introdução a Sistemas de Banco de dados. 8. Ed. 900 p. Rio
de Janeiro : Campus, 2004\\

\newpage IDENTIFICAÇÃO

DISCIPLINA: Tópicos Especiais em Banco de Dados CÓDIGO:\\
DEPARTAMENTO: Unidade Acadêmica de Garanhuns\\
ÁREA: \\
CARGA HORÁRIA TOTAL: 60 NÚMERO DE CRÉDITOS: 04\\
CARGA HORÁRIA SEMANAL: 04 TEÓRICAS: 02 PRÁTICAS: 02 TOTAL: 04\\
PRÉ-REQUISITOS: Banco de Dados\\
CO-REQUISITOS: não há

EMENTA 

Estudo de técnicas avançadas em na área multidisciplinar de Banco de
Dados permitindo ao aluno conhecer o estado da arte nesta área de
pesquisa.

BIBLIOGRAFIA 

Artigos e Períodicos Recentes da Área de Pesquisa
\newpage 

\section*{Área Temática: Engenharia da Computação}

IDENTIFICAÇÃO

DISCIPLINA: Sistemas de Tempo Real CÓDIGO:\\
DEPARTAMENTO: Unidade Acadêmica de Garanhuns\\
ÁREA: \\
CARGA HORÁRIA TOTAL: 60 NÚMERO DE CRÉDITOS: 04\\
CARGA HORÁRIA SEMANAL: 04 TEÓRICAS: 02 PRÁTICAS: 02 TOTAL: 04\\
PRÉ-REQUISITOS: Arquitetura e Organização de Computadores\\
CO-REQUISITOS: não há

EMENTA 

Sistemas de tempo real. Processos de tempo real. Interações entre
processos. Tempo de execução. Escalonamento de processos. Garantia de
escalonamento.Kernels e sistemas operacionais de tempo real. Introdução a tolerância a
falhas. Dispositivos para aumentar robustez de sistemas embarcados

BIBLIOGRAFIA 

Básica:\\
BURNS, A.; WELLINGS, A.; Real-Time Systems and Programming Languages. 3a
Ed. Addison-Wesley. 2001.\\
KOPETZ, H. Real-Time Systems: Design Principles for Distributed Embedded
Applications. Springer. 1997.\\
FARINES, J. M.; FRAGA, J. S.; OLIVEIRA, R. S. 12. Escola de Computação.
IME-USP. 2000.\\
Complementar:\\
GOLDSMITH, S. A practical guide to Real-Time Systems. Prentice Hall.
1993.\\
SHAW, A.C. Sistemas e Software de Tempo Real. Bookman. 2003.\\
LI, Q.; YAO, C. Real-Time Concepts for Embedded Systems. CMP. 2003.

\newpage IDENTIFICAÇÃO

DISCIPLINA: Projeto de Sistemas Embarcados CÓDIGO:\\ 
DEPARTAMENTO: Unidade Acadêmica de Garanhuns\\
ÁREA: \\
CARGA HORÁRIA TOTAL: 60 NÚMERO DE CRÉDITOS: 04\\
CARGA HORÁRIA SEMANAL: 04 TEÓRICAS: 02 PRÁTICAS: 02 TOTAL: 04\\
PRÉ-REQUISITOS: Arquitetura e Organização de Computadores\\
CO-REQUISITOS: não há

EMENTA 

Modelos de Especificação de Sistemas Embutidos. Técnicas de
Particionamento de Sistemas. Técnicas de Estimativas. Técnicas para
Geração de Interfaces. Técnicas para Síntese de Software. Técnicas de
Co-simulação. Prototipação de Sistemas.

BIBLIOGRAFIA 

Básica:\\
VAHID, F.; GIVARGIS, T. Embedded system design: a unified
hardware/software introduction. Willey. 2002.\\
GAJSKI, D. et al. Embedded System Design  Modeling, Synthesis and
Verification. Springer. 2009.\\
CATSOULIS, J. Designing Embedded Hardware. 2. Ed.
O'Reilly Media. 2005.\
Complementar:\
HENNESSY, J. L; PATTERSON, D. Arquiteturas de Computadores -- Uma
abordagem quantitativa. 3. Ed. Campus. 2003.\\
SASS, R.; SCHMIDT, A.G. Embedded Systems Design with Platform FPGAs --
Principles and Practices. Morgan Kaufmann. 2010.\\
SIMON, E. S. An Embedded Software Primer. Addison-Wesley Professional.
1999.

\newpage IDENTIFICAÇÃO

DISCIPLINA: Prototipação de Circuitos Digitais CÓDIGO:\\ 
DEPARTAMENTO: Unidade Acadêmica de Garanhuns\\
ÁREA: \\
CARGA HORÁRIA TOTAL: 60 NÚMERO DE CRÉDITOS: 04\\
CARGA HORÁRIA SEMANAL: 04 TEÓRICAS: 02 PRÁTICAS: 02 TOTAL: 04\\
PRÉ-REQUISITOS: Arquitetura e Organização de Computadores
CO-REQUISITOS: não há

EMENTA 

Ferramentas de CAD (Computer Aided-Design). Metodologias de projeto.
Tecnologia para implementação de circuitos de alta integração. Estilos
de projetos para implementação de circuitos integrados. Projeto e
implementação de circuitos integrados usando ferramentas de CAD.
Laboratório/projeto de um estudo de caso.

BIBLIOGRAFIA 

Básica:\\
WESTE, N., ESHRAGHIAN, K., Principles of CMOS VLSI Design- A Systems
Perspective, Addison-Wesley Publishing Company, 1988.\\
REIS, Ricardo A. L., Concepção de Circuitos integrados. Série Livros
didáticos. Editora Sagra Luzzatto, 2000.\\
Complementar:\\
Notas de aula; Artigos científicos

\newpage IDENTIFICAÇÃO

DISCIPLINA: Tópicos Avançados em Engenharia da Computação CÓDIGO:\\ 
DEPARTAMENTO: Unidade Acadêmica de Garanhuns\\ 
ÁREA:\\
CARGA HORÁRIA TOTAL: 60 NÚMERO DE CRÉDITOS: 04\\
CARGA HORÁRIA SEMANAL: 04 TEÓRICAS: 02 PRÁTICAS: 02 TOTAL: 04\\
PRÉ-REQUISITOS: não há\\
CO-REQUISITOS: não há

EMENTA 

Estudo de técnicas avançadas em Engenharia da Computação permitindo ao
aluno conhecer o estado da arte nesta área de pesquisa.

BIBLIOGRAFIA 

Artigos recentes da área de pesquisa.
\newpage 

\section*{Área Temática: Engenharia de Software}

IDENTIFICAÇÃO

DISCIPLINA: Engenharia de Software Experimental CÓDIGO:\\ 
DEPARTAMENTO: Unidade Acadêmica de Garanhuns\\
ÁREA: \\
CARGA HORÁRIA TOTAL: 60 HORAS NÚMERO DE CRÉDITOS: 04\\
CARGA HORÁRIA SEMANAL: 04 TEÓRICAS: 03 PRÁTICAS: TOTAL: 01\\
PRÉ-REQUISITOS: não há\\
CO-REQUISITOS: não há

EMENTA 

Introdução a Princípios de Experimentação; Métodos e Técnicas de
Experimentação: Estudo de Caso; Survey; Experimento. Projeto, Execução
e Avaliação de Experimentos em Engenharia de Software; Definição de
Métricas; Empacotamento e Replicação de Experimentos.

BIBLIOGRAFIA 

Básica:\\
TRAVASSOS, G.; GUROV, D.; AMARAL, E. Introdução a Engenharia de Software
Experimental. Technical Report ES-590/02, COPPE/UFRJ, Abril, 2002.\\
SHULL, Forrest., SINGER, Janice., SJøBERG, Dag I. K.. Guide to Advanced
Empirical Software Engineering; Springer, 2010.\\
BASILI, V.; SALBY, R.; HUTCHENS, D. Experimentation in Software
Engineering. IEEE Transactions on Software Engineering,
SE-12(7):733-743, Julho, 1986.\\
Complementar:\\
JAIN, R. The art of computer systems performance analysis: techniques
for experimental design, measurement, simulation and modeling.
Wiley,1991.\\
CONRADI, Reidar. Empirical Methods and Studies in Software Engineering:
Experiences from ESERNET; Springer, 2003.\\
PFLEEGER, S. Design and Analysis in Software Engineering, Part 1: The
Language of Case Studies and Formal Experiments. Software.

\newpage IDENTIFICAÇÃO

DISCIPLINA: Programação Orientada a Aspectos CÓDIGO:\\ 
DEPARTAMENTO: Unidade Acadêmica de Garanhuns\\ ÁREA: \\
CARGA HORÁRIA TOTAL: 60 HORAS NÚMERO DE CRÉDITOS: 04\\
CARGA HORÁRIA SEMANAL: 04 TEÓRICAS: 02 PRÁTICAS: 02 TOTAL: 04\\
PRÉ-REQUISITOS: não há
CO-REQUISITOS: não há

EMENTA 

Introdução ao paradigma orientado a aspectos: Problemas no paradigma OO; Modularidade; Extensibilidade; Reusabilidade; AspectJ; Aspectos; Pointcuts; Join points; Advice; Static and dynamic crosscutting; Discutir detalhes da semântica estática e dinâmica de AspectJ; Apresentação de exemplos de aplicações do paradigma orientado a aspectos; Padrões de projetos orientados a aspectos; Padrões GOF (implementações em Java e AspectJ ); PDC (Persistent Data Collection); Refactorings orientados a aspectos; AspectJ dioms para construção de software orientado a aspectos; Catálogo de refactorings e code smells para AspectJ.

BIBLIOGRAFIA 

Básica:\\
Ladda, Ramnivas. AspectJ in Action: Practical Aspect-Oriented
programming. Greenwich, CT, USA, 1 ed., 2003.\\
Kiselev, I. Aspect-Oriented Programming with AspectJ. ISBN:~0672324105.
Sams Publisher. 1 ed., 2002.\\
Miles, R.; AspectJ Cookbook. ISBN: 0596006543. Ed. O'REILLY \& ASSOC.
2009.\\
Complementar:\\
Sommerville, Ian. Engenharia de Software. ISBN.: 8588639076. Addison --
Wesley.\\
Roger S. Pressman. Engenharia de Software. ISBN.: 8586804576.
McGraw-Hill.\

\newpage IDENTIFICAÇÃO

DISCIPLINA: Programação Orientada a Objetos II CÓDIGO:\\ 
DEPARTAMENTO: Unidade Acadêmica de Garanhuns\\ 
ÁREA: \\
CARGA HORÁRIA TOTAL: 60 HORAS NÚMERO DE CRÉDITOS: 04\\
CARGA HORÁRIA SEMANAL: TEÓRICAS: PRÁTICAS: TOTAL: \\
PRÉ-REQUISITOS: não há\\
CO-REQUISITOS: não há

EMENTA 

Desenvolvimento de Interface Gráfica com Usuário (Aplicações Desktop); Desenvolvimento em equipe; Utilização de SVN e CVS; Merge e Branch de códigos fonte; Instalações de Pluggins para Eclipse; Utilização de APIs; XML; Utilização de JUnit; Construção e utilização de Build.xml; Eclipse Debugging.

BIBLIOGRAFIA 

Básica:\\
Harvey M. Deitel, Paul J. Deitel. Java How to Program. Prentice Hall. 7.
Edição, 2006. ISBN : 0132222205.\\
Lynn Andrea Stein.~Interactive Programming in Java, 2003.\\
Bruce Eckel.~Thinking in Java. Segunda edição, 2000.\\
Complementar:\\
Barnes, D. J., Kölling, M. Programação Orientada a Objetos com Java, Ed.
Pearson/Prentice-Hall, 2004.\\
Santos, R. Introdução à Programação Orientada à Objetos Usando Java, Ed.
Campus, 2003.\\
Cay S. Horstmann e Gary Cornell.~Core Java 2, Volume I  Fundamentos.
Sun Microsystems Press, Makron Books do Brasil, 2001.

\newpage IDENTIFICAÇÃO

DISCIPLINA: Programação Paralela e Distribuída CÓDIGO:\\ 
DEPARTAMENTO: Unidade Acadêmica de Garanhuns\\
ÁREA: \\
CARGA HORÁRIA TOTAL: 60 HORAS NÚMERO DE CRÉDITOS: 04\\
CARGA HORÁRIA SEMANAL: TEÓRICAS: 3 PRÁTICAS: TOTAL: 1\\
PRÉ-REQUISITOS: não há\\
CO-REQUISITOS: não há

EMENTA 

Arquiteturas paralelas, programação paralela e aspectos de desempenho;
Processos, comunicação e sincronização (IPC); Threads, comunicação e
sincronização em memória compartilhada; Paralelismo com threads;
Algoritmos de escalonamento; Processadores paralelos e distribuídos;
Comunicação em Rede (sockets); Computação com Passagem de Mensagem
(MPI)

BIBLIOGRAFIA 

Básica:\\
Wilkinson, B. and Allen, M. Parallel Programming: Techniques and
Applications Using Networked Workdstations and Parallel Computers.
Pearson Prentice Hall, 2005.\\
Dongarra,J.;Foster,I.;Fox,G.;Gropp,W.;White,A.;Torczon,L.;Kennedy,K.
Sourcebook of Parallel Computing. Morgan Kaufmann Pub, 2002.\\
DEITEL, Harvey M.; DEITEL, Paul J. Java: Como Programar. 4a Edição.
Bookman, 2002.\\
Complementar:\\
ORFALI, Robert; HARVEY, Dan. Client/Server Programming with Java and
CORBA. 2nd Edition. John Wiley, 1998.\\
Quinn, M.J.Parallel Programming in C with MPI and OpenMP. McGrawHill,
2004.\\
Grama,A.;Gupta,A.;Karypis,G.;Kumar,V. Introduction to Parallel
Computing. Adisson-Wesley, 2003.

\newpage IDENTIFICAÇÃO

DISCIPLINA: Programação Web CÓDIGO:\\ 
DEPARTAMENTO: Unidade Acadêmica de Garanhuns\\ 
ÁREA: \\
CARGA HORÁRIA TOTAL: 60 HORAS NÚMERO DE CRÉDITOS: 04\\
CARGA HORÁRIA SEMANAL: TEÓRICAS: PRÁTICAS: TOTAL: \\
PRÉ-REQUISITOS: não há\\
CO-REQUISITOS: não há

EMENTA 

Introdução à Programação WEB (Elementos Básicos); Desenvolvimento de interfaces gráficas para a Internet; Desenvolvimento de serviços e sistemas de informação para a Internet; Servidores de Aplicação; Segurança e desempenho de sistemas para a Internet; Integração de sistemas baseados na Internet; JBoss; Desenvolvimento de interfaces gráficas para a Internet com HTML, Javascript, e Applets); Desenvolvimento de serviços e sistemas de informação para a Internet com Servlets, Java Server Pages; Performance de sistemas para a Internet; XML/CSS/XLST (Web Design)

BIBLIOGRAFIA 

Básica:\\
Nicholas Kassem ET. AL.;~Designing Enterprise Applications with the Java
2 Platform, Enterprise Edition. Addison-Wesley, 2000.\\
Jason Hunter, William Crawford. Java Servlet Programming.
O'Reilly \& Associates, 1998.\\
Baehr, C. M; Web Development A Visual-Spatial Approach. ISBN:
0131701223. Ed. Prentice Hall. 2006.\\
Complementar:\\
David Flanagan, Dan Shafer. Javascript: The Definitive Guide.
O'Reilly \& Associates, 3rd edition, 1998\\
Downet, T.; Web Development with Java. ISBN: 1846288622. Ed. Springer
Verlag NY. 2007

\newpage IDENTIFICAÇÃO

DISCIPLINA: Desenvolvimento Distribuído de Software CÓDIGO:\\ 
DEPARTAMENTO: Unidade Acadêmica de Garanhuns\\ 
ÁREA: \\
CARGA HORÁRIA TOTAL: 60 HORAS NÚMERO DE CRÉDITOS: 04\\
CARGA HORÁRIA SEMANAL: TEÓRICAS: 3 PRÁTICAS: TOTAL: 1\\
PRÉ-REQUISITOS: não há\\
CO-REQUISITOS: não há

EMENTA 

Gerência de Projeto; Engenharia de Requisitos; Processos de Software; Adaptação de Processos; Modelos de desenvolvimento; Métodos Ágeis; Scrum; XP.

BIBLIOGRAFIA 

Básico:\\
PRIKLADNICKI, Rafael., AUDY, Jorge. Desenvolvimento Distribuído de
Software; Elsevier, 2007.\\
CARMEL, E. Global Software Teams: Collaboration Across Borders and Time
Zones;Prentice-Hall, EUA, 1999.\\
ECKSTEIN, Jutta. Agile Software Development with Distributed Teams;
Dorset House, 2010.\\
Complementar:\\
COHN, Mike. Succeeding with Agile: Software Development Using Scrum;
Addison-Wesley, 2009.\\
LOELIGER, Jon. Version Control with Git: Powerful Tools and Techniques
for Collaborative Software Development; O'Reilly,
2009.\\
KAROLAK, Dale Walter. Global Software Development: Managing Virtual
Teams and Environments; Wiley, 1998.

\newpage IDENTIFICAÇÃO

Disciplina: Análise e Projeto de Sistemas CÓDIGO:\\ 
Departamento: Unidade Acadêmica de Garanhuns\\
ÁREA: \\
Carga Horária Total: 60 horas Número de créditos: 04\\
Carga Horária Semanal:04 horas Teóricas: 50h Práticas: 10h Total: 60h\\
Pré-requisitos: Engenharia de Software\\
CO-REQUISITOS: não há

EMENTA 

Apresentar aos alunos os princípios de linha de produção dentro da Eng. de Software. Este conceito define as etapas de como um software orientado a objeto deve ser construído desde seu levantamento de requisitos até a sua implantação.

BIBLIOGRAFIA 

Básica:\\
PRESMANN. Engenharia de Software 7. ed. 776 p. McGraw-Hill, 2011\\
HELM et. al. Padrões de Projeto 1. Ed. 366p. Bookman. 2005.\\
BEZERRA. Princípios de Análise e Projetos de Sistemas UML. 380p, Campus 2006.\\
Complementar:\\
BRAUDE. Projeto de Software. Bookman. 2005.\\
SOMMERVILLE. Engenharia de Software. 8. ed. 568 p. Addison Wesley, 2007\\
WEST et. al. Use A Cabeça Analise \& Projeto Orientado Ao Objeto. 472 p.
Starling Consult. 2007.

\newpage IDENTIFICAÇÃO

DISCIPLINA: Gerenciamento de Projetos CÓDIGO:\\ 
DEPARTAMENTO: Unidade Acadêmica de Garanhuns\\
ÁREA: \\
CARGA HORÁRIA TOTAL: 60 HORAS NÚMERO DE CRÉDITOS: 04\\
CARGA HORÁRIA SEMANAL: TEÓRICAS: PRÁTICAS: TOTAL: \\
PRÉ-REQUISITOS: não há\\
CO-REQUISITOS: não há

EMENTA 

Fundamentos de Gerenciamento de Projetos (GP); Descrição do Processo de GP; Atividades de planejamento e gerenciamento de um processo específico (o RUP, por exemplo); Desenvolvimento de Proposta e Plano de Projeto; Atividade de Iniciação de um Projeto; Atividade de Gerenciamento de Tempo; Atividade de Gerenciamento de Escopo; Atividade de Gerenciamento de Riscos; Fases de Implementação do Projeto; Fases de Implantação; Fase de Finalização do Projeto de Software; Ferramentas de planejamento e gerenciamento; Técnicas para estimativas e coleta de métricas

BIBLIOGRAFIA 

Básica:\\
Jack R. Meredith e Samuel J. Mantel, Jr.; Project Management: A
Managerial Approach, Third Edition, John Wiley \& Sons Inc., USA,
1995.\\
Robert K., Wysocki, Robert, Jr Beck, David B. Crane; Effective Project
Management, 2nd Edition. 384 pages. March 2, 2000. John Wiley \& Sons.
ISBN 0471360287.03 -- Bruce Eckel.\\
Barry W. Boehm; Software Engineering Economics. 767 pages (October
1981). Prentice Hall.~ISBN 01382212\\
Complementar:\\
Steve McConnell. Software Project Survival Guide. 250 pages. November
1997. Microsoft Press. ISBN 1572316217.\\
Tom Demarco, Timothy Lister. Peopleware: Productive Projects and Teams.
2nd~edition, February 1, 1999. Dorset House. ISBN 0932633439.\\
Frederick P., Jr. Brooks, Frederick P. Brooks Jr.; The Mythical
Man-Month : Essays on Software Engineering. 336 pages anniversary
edition (July 1995). Addison-Wesley Pub Co. ISBN 0201835959.\

\newpage IDENTIFICAÇÃO

DISCIPLINA: Especificação de Requisitos CÓDIGO:\\ 
DEPARTAMENTO: Unidade Acadêmica de Garanhuns\\ 
ÁREA: \\
CARGA HORÁRIA TOTAL: 60 HORAS NÚMERO DE CRÉDITOS: 04\\
CARGA HORÁRIA SEMANAL: 04 TEÓRICAS: 02 PRÁTICAS: 02 TOTAL: 04\\
PRÉ-REQUISITOS: Engenharia de Software\\
CO-REQUISITOS: não há

EMENTA 

Motivação; Técnicas de especificação de requisitos; Ferramentas de especificação de requisitos; Técnicas de validação de sistemas; Ferramentas de validação de sistemas; Estudo de caso; Projeto

BIBLIOGRAFIA 

Básica:\\
PRESMANN. Engenharia de Software 7. ed. 776 p. McGraw-Hill, 2011\\
KONTOYA \& SOMMERVILLE: Requirements Engineering: Processes and
Tecniques, John Wiley \& Sons, Ltd, 1998\\
COHN: User Stories Applied, Addison Wesley. 2004\\
Complementar:\\
SOMMERVILLE \& SAWYER: Requirements Engineering. John Wiley \& Sons,
Ltd, 1997 \\
SOMMERVILLE. Engenharia de Software. 8. ed. 568 p. Addison Wesley, 2007\\
LAMSWEERDE. Requirements Engineering. John Wiley \& Sons, Ltd, 2009\\

\newpage IDENTIFICAÇÃO

DISCIPLINA: Qualidade de Software CÓDIGO:\\ 
DEPARTAMENTO: Unidade Acadêmica de Garanhuns\\
ÁREA: \\
CARGA HORÁRIA TOTAL: 30 NÚMERO DE CRÉDITOS: 02\\
CARGA HORÁRIA SEMANAL: 02 TEÓRICAS: 01 PRÁTICAS: 01 TOTAL: 02\\
PRÉ-REQUISITOS: Engenharia de Software\\
CO-REQUISITOS: não há

EMENTA 

Apresentar aos alunos o tema Qualidade de Software, procurando discutir
aspectos relacionados a essa sub-área da Engenharia de Software.

BIBLIOGRAFIA 

Básica:\\
PRESMANN. Engenharia de Software 7. ed. 776 p. McGraw-Hill, 2011\\
KOSCIANSKI. Qualidade de Software, Editora Novatec, 2006.\\
BROOKS. O Mítico homem mês 1.Ed. 328p. Campus. 2009.\\
Complementar:\\
BARTIE. Garantia da Qualidade de Software, Editora Campus, 2002.\\
SOMMERVILLE. Engenharia de Software. 8. ed. 568 p. Addison Wesley, 2007\\
CHRISSIS et. al. CMMI -- Guidelines for Process Integration and Product
Improvement, Addison-Wesley, 2003.

\newpage IDENTIFICAÇÃO

DISCIPLINA: Paradigmas de Linguagens de Programação CÓDIGO:\\ 
DEPARTAMENTO: Unidade Acadêmica de Garanhuns\\
ÁREA: \\
CARGA HORÁRIA TOTAL: 60 HORAS NÚMERO DE CRÉDITOS: 04\\
CARGA HORÁRIA SEMANAL: 04 TEÓRICAS: 03 PRÁTICAS: TOTAL: 04 \\
PRÉ-REQUISITOS: não há\\
CO-REQUISITOS: não há

EMENTA 

Inspeção de software; Princípios e técnicas de testes de software; teste
de unidade; teste de integração; teste de regressão; Desenvolvimento
orientado a testes; Automação dos testes; Geração de casos de teste;
Teste de interfaces humanas; Teste de aplicações para a web; Testes
alfas, beta e de aceitação; Ferramentas de testes; Planos de testes.
Gerenciamento do processo de testes. Registro e acompanhamento de
problemas.

BIBLIOGRAFIA 

Básica:\\
DELAMARO. M.E.; MALDONADO, J. C.; JINO, M. Introdução ao Teste de
Software; Rio de Janeiro: Campus, 2007.\\
MOLINARI, Leonardo. Testes Funcionais de Software; Visual Books, 2008.\\
PRESSMAN, Roger S.. Engenharia de Software: uma Abordagem Profissional;
Bookman, 2011.\\
Complementar:\\
BECK, K. Test-driven development by example; EUA: Addison Wesley, 2002.\\
BLACK, Rex. Pragmatic Software Testing; Wiley, 2007.\\
MOLINARI, Leonardo. Testes de Software: Produzindo Sistemas Melhores e
mais Confiáveis; Érica, 2003.

\newpage IDENTIFICAÇÃO

DISCIPLINA: Medição de Software CÓDIGO:\\ 
DEPARTAMENTO: Unidade Acadêmica de Garanhuns\\
ÁREA: \\
CARGA HORÁRIA TOTAL: 30 HORAS NÚMERO DE CRÉDITOS: 02\\
CARGA HORÁRIA SEMANAL: 02 TEÓRICAS: 01 PRÁTICAS: 01 TOTAL: 02\\
PRÉ-REQUISITOS: Engenharia de Software\\
CO-REQUISITOS: não há

EMENTA 

Apresentar aos alunos os princípios subjacentes ao desenvolvimento de software que precisa para satisfazer objetivos específicos, onde estas metas devem ser expressas em termos mensuráveis Ela abrange os princípios de medição de software, e as maneiras pelas quais elas são usadas no planejamento de projetos de software, e no acompanhamento dos projetos bem como sua realização. Discute a aplicação de técnicas de medição de software para estes métodos de ensaio.

BIBLIOGRAFIA 

Básica:\\
PRESMANN. Engenharia de Software 7. ed. 776 p. McGraw-Hill, 2011\\
JONES. Applied Software Measurement 3. Ed. 700 p. McGraw-Hill, 2008\\
BROOKS. O Mítico homem mês 1.Ed. 328p. Campus. 2009.\\
Complementar:\\
VAZQUES ET. AL. Análise de Pontos de Função: Medição, Estimativas e
Gerenciamento de Projetos de Software  7. Edição. Ed. Érica. 2007\\
SOMMERVILLE. Engenharia de Software. 8. ed. 568 p. Addison Wesley, 2007\\
MCGARRY ET. AL. Practical Software Measurement: Objective Information
for Decision Makers. Addison-Wesley. 2001.

\newpage IDENTIFICAÇÃO

DISCIPLINA: Paradigmas de Linguagens de Programação CÓDIGO:\\ 
DEPARTAMENTO: Unidade Acadêmica de Garanhuns\\
ÁREA: \\
CARGA HORÁRIA TOTAL: 60 HORAS NÚMERO DE CRÉDITOS: 04\\
CARGA HORÁRIA SEMANAL: 04 TEÓRICAS: 03 PRÁTICAS: 01 TOTAL: 04 \\
PRÉ-REQUISITOS: não há\\
CO-REQUISITOS: 

EMENTA 

Inspeção de software; Princípios e técnicas de testes de software; teste
de unidade; teste de integração; teste de regressão; Desenvolvimento
orientado a testes; Automação dos testes
Geração de casos de teste; Teste de interfaces humanas; Teste de
aplicações para a web. 
Testes alfas, beta e de aceitação; Ferramentas de testes; Planos de
testes. Gerenciamento do processo de testes. Registro e acompanhamento
de problemas.

BIBLIOGRAFIA 

Básica:\\
DELAMARO. M.E.; MALDONADO, J. C.; JINO, M. Introdução ao Teste de
Software; Rio de Janeiro: Campus, 2007.\\
MOLINARI, Leonardo. Testes Funcionais de Software; Visual Books, 2008.\\
PRESSMAN, Roger S.. Engenharia de Software: uma Abordagem Profissional;
Bookman, 2011.\\
Complementar:\\
BECK, K. Test-driven development by example; EUA: Addison Wesley, 2002.\\
BLACK, Rex. Pragmatic Software Testing; Wiley, 2007.\\
MOLINARI, Leonardo. Testes de Software: Produzindo Sistemas Melhores e
mais Confiáveis; Érica, 2003.

\newpage IDENTIFICAÇÃO

DISCIPLINA: Tópicos Especiais em Engenharia de Software CÓDIGO: \\
DEPARTAMENTO: Unidade Acadêmica de Garanhuns\\
ÁREA: \\
CARGA HORÁRIA TOTAL: 60 NÚMERO DE CRÉDITOS: 04\\
CARGA HORÁRIA SEMANAL:04 horas TEÓRICAS: 60h PRÁTICAS: 0h TOTAL: 60h\\
PRÉ-REQUISITOS: Engenharia de Software\\
CO-REQUISITOS: não há

EMENTA 

Estudo de técnicas avançadas em na área multidisciplinar de Engenharia
de Software permitindo ao aluno conhecer o estado da arte nesta área de
pesquisa.

BIBLIOGRAFIA 

Artigos e Periódicos Recentes da Área de Pesquisa.

\newpage IDENTIFICAÇÃO

DISCIPLINA: Metodologias Ágeis CÓDIGO: \\
DEPARTAMENTO: Unidade Acadêmica de Garanhuns\\
ÁREA: \\
CARGA HORÁRIA TOTAL: 60 horas NÚMERO DE CRÉDITOS: 04\\
CARGA HORÁRIA SEMANAL:04 horas TEÓRICAS: 30h PRÁTICAS: 30h TOTAL: 60h\\
PRÉ-REQUISITOS: Gestão de Projetos e Análise e Projeto de Sistemas\\
CO-REQUISITOS: não há

EMENTA 

Apresentar aos alunos os problemas da Engenharia de Software
Tradicional. Como enxergar o desenvolvimento de software como uma
atividade de rápidas mudanças e como as metodologias ágeis trata esta
questão.

BIBLIOGRAFIA 

Básica:\\
Agile Planning and Estimation. Mike Cohn: Addison Wesley, 2004.\\
User Stories Applied. Por Mike Cohn  Addison Wesley (2003)\\
Agile Software Development Quality Assurance. Stamelos \& Sfetsos:
Information Science Centre, 2007.\\
Complementar:\\
BECK. K. Extreme Programming Explained. Addison Wesley. 2000.\\
SOMMERVILLE. Engenharia de Software. 8. ed. 568 p. Addison Wesley, 2007\\
BOEHM \& TURNER Balancing Agility and Discipline. Addison Wesley. 2003.
\newpage 

\section*{Área Temática: Inteligência Computacional}

IDENTIFICAÇÃO

DISCIPLINA: Redes Neurais CÓDIGO:\\ 
DEPARTAMENTO: Unidade Acadêmica de Garanhuns\\
ÁREA: \\
CARGA HORÁRIA TOTAL: 60 HORAS NÚMERO DE CRÉDITOS: 04\\
CARGA HORÁRIA SEMANAL: TEÓRICAS: PRÁTICAS: TOTAL: \\
PRÉ-REQUISITOS: Reconhecimento de Padrões\\
CO-REQUISITOS: não há

EMENTA 

Processos de aprendizagem de máquina. Perceptrons de camada única.
Perceptrons de múltiplas camadas. Redes de função de base radial. Mapas
auto-organizáveis. Tópicos avançados em Redes Neurais: máquinas de
vetor de suporte, análise de componentes principais, outros tópicos.

BIBLIOGRAFIA 

Básica:\\
BRAGA, A. P., LUDEMIR, T. B., CARVALHO, A. P. L. F. A Redes Neurais Artificiais  Teoria e Aplicações. LTC, 2007.\\
HAYKIN, Simon. Redes Neurais: Princípios e Prática, 3. edição. Bookman.
2008.\\
COPPIN, B. Inteligência Artificial. LTC, 2010.\\
Complementar:\\
MARQUES, J. S.. Reconhecimento de Padrões  Métodos Estatísticos e Neurais. IST Press, 2005.\\
WITTEN, I. H., FRANK, E. Data mining : practical machine learning tools and techniques. Second Edition. Elsevier, 2005.\\
THEODORIDIS, S., KOUTROUMBAS, K. Pattern Recognition. Fourth Edition, Academic Press, 2009.\\
DUDA, R.O., HART, P.E., STORK, D.G. Pattern Classification. Second Edition. Wiley, 2001.\\
BISHOP, C. M. Pattern Recognition and Machine Learning. Springer, 2006.\\
MITCHELL, T. Machine Learning. McGraw-Hill, 1997.

\newpage IDENTIFICAÇÃO

DISCIPLINA: Computação Evolucionária CÓDIGO:\\ 
DEPARTAMENTO: Unidade Acadêmica de Garanhuns\\
ÁREA: \\
CARGA HORÁRIA TOTAL: 60 HORAS NÚMERO DE CRÉDITOS: 04\\
CARGA HORÁRIA SEMANAL: TEÓRICAS: PRÁTICAS: TOTAL: \\
PRÉ-REQUISITOS: Inteligência Artificial\\
CO-REQUISITOS: não há

EMENTA 

A metáfora biológica. Elementos de um algoritmo de computação
evolucionária (CE). Algoritmos genéticos. Estratégia de evolução. Programação evolucionária. Programação genética. Controle de
parâmetros. Problema multi-modais. Satisfação de restrição. Tópicos avançados. Projeto.

BIBLIOGRAFIA 

Básica:\\
Eiben, A. E. \& Smith, J. E. . Introduction to Evolutionary Computing.
Springer Verlag, 2003.\\
Linden, R. Algoritmos Geneticos. BRASPORT, 2008.\\
Coello, C. A., Lamont, G. B., \& van Veldhuizen, D. A.. Evolutionary
Algorithms for Solving Multi-Objective Problems, Springer, 2nd Edition,
2007.\\
Complementar:\\
Fogel, D. B.. Evolutionary Computation, IEEE Press, 2003.\\
Ghosh, A. \& Tsutsui, S.. Advances in Evolutionary Computing: Theory and
Applications. Springer, 2003.\\
Goldberg, D. E.. of Genetic Algorithms in Search, Optimization, and
Machine Learning. Reading, Ma: Addison-Wesley, 1989.

\newpage IDENTIFICAÇÃO

DISCIPLINA: Mineração de Dados CÓDIGO:\\
DEPARTAMENTO: Unidade Acadêmica de Garanhuns\\
ÁREA:\\
CARGA HORÁRIA TOTAL: 60 HORAS NÚMERO DE CRÉDITOS: 04\\
CARGA HORÁRIA SEMANAL: TEÓRICAS: PRÁTICAS: TOTAL:\\
PRÉ-REQUISITOS: Inteligência Artificial e Reconhecimento de Padrões\\
CO-REQUISITOS: não há

EMENTA 

Introdução e Motivação ao Processo de Descoberta de Conhecimento em
Bases de Dados (KDD). Etapas do Processo de KDD. Conceitos e
Tecnologias de Suporte à Mineração de Dados. Pré-processamento dos
Dados. Extração de Padrões: Tarefas, Algoritmos e Paradigmas de
Mineração de Dados. Pós-processamento de Resultados. Métricas de
Avaliação: Complexidade, Eficiência e Escalabilidade. Tópicos
Avançados: Metaheurísticas, Paralelismo e Distribuição, Visualização,
Privacidade e Segurança, Representações e Estruturas de Dados
Nãoconvencionais, Mineração Multimodal (Textos e Multimídia), Mineração
de Dados Espaciais e Temporais. Técnicas, Ferramentas e Aplicações.

BIBLIOGRAFIA 

Básica:\\
I. H. Witten, E. Frank:  Data Mining: Practical
Machine Learning Tools and Techinques , Morgan
Kaufmann, 2005.\\
P. N. Tan, M. Steinbach, V. Kumar:  Introduction to
Data Mining , Addison-Wesley, 2005.\\
J. Han, M. Kamber:  Data Mining: Concepts and
Techniques , 2a. Ed., Morgan Kaufmann, 2006.\\
A. Abraham, C. Grosan, V. Ramos:  Swarm Intelligence
in Data Mining , Springer, 2006.\\
Complementar:\\
D. Hand, H. Mannila, P. Smith:  Principles of Data
Mining , MIT Press, 2001.\\
O. Maimon, L. Rokach:  Data Mining and Knowledge
Discovery Handbook , Springer, 2005.\\
D. T. Larose:  Discovering Knowledge in Data: An
Introduction to Data Mining , John Wiley, 2005.\\
M. Kantardzic, J. Zurada:  Next Generation of
Data-Mining Applications , Wiley-IEEE Press, 2005.\\
S. Mitra, T. Acharya:  Data Mining: Multimedia, Soft
Computing, and Bioinformatics , John Wiley, 2003.\\
Artigos recentes da área.

\newpage IDENTIFICAÇÃO

DISCIPLINA: Visão Computacional CÓDIGO:\\ 
DEPARTAMENTO: Unidade Acadêmica de Garanhuns\\
ÁREA: \\
CARGA HORÁRIA TOTAL: 60 HORAS NÚMERO DE CRÉDITOS: 04\\
CARGA HORÁRIA SEMANAL: TEÓRICAS: PRÁTICAS: TOTAL: \\
PRÉ-REQUISITOS: Reconhecimento de Padrões\\
CO-REQUISITOS: não há

EMENTA 

Introdução à visão computacional. Ferramentas de apoio.
Pré-processamento e Filtros. Segmentação. Rastreamento. Reconhecimento
e classificação de Padrões em Imagens e Vídeos. Avaliação de desempenho
de algoritmos de visão computacional. Seminário. Projeto.

BIBLIOGRAFIA 

Básica:\\
Richard Szeliski.Computer Vision: Algorithms and Applications. Springer,
2010.\\
Nicu Sebe, Michael S. Lew. Robust Computer Vision: Theory and
Applications. Kluwer / Springer, 2003.\\
Richard Hartley, Andrew Zisserman. Multiple View Geometry in Computer
Vision, 2nd Edition. Cambridge University Press, 2004.\\
Complementar:\\
Artigos recentes da área.

\newpage IDENTIFICAÇÃO

DISCIPLINA: Agentes Inteligentes CÓDIGO:\\ 
DEPARTAMENTO: Unidade Acadêmica de Garanhuns\\
ÁREA: \\
CARGA HORÁRIA TOTAL: 60 HORAS NÚMERO DE CRÉDITOS: 04\\
CARGA HORÁRIA SEMANAL: TEÓRICAS: PRÁTICAS: TOTAL: \\
PRÉ-REQUISITOS: Inteligência Artificial\\
CO-REQUISITOS: não há

EMENTA 

Sistemas Multiagentes, princípios gerais e aplicações. Agentes autônomos
e sistemas multiagentes. Introdução à resolução distribuída de
problemas. Cooperação, Coordenação e Negociação. Comunicação entre
agentes. Arquiteturas de comunicação. Linguagens de comunicação e
conteúdo. Protocolos de interação. Modelos e arquiteturas de agentes.
Taxonomia de Agentes. Agentes autônomos, reativos, deliberativos e
adaptativos. Ontologias.

BIBLIOGRAFIA 

Básica:\\
COPPIN, B. Inteligência Artificial. LTC, 2010.\\
RUSSELL, S. and NORVIG, P. Artificial Intelligence: a Modern Approach.
Prentice Hall, 2a edição, 2002.\\
WOOLDRIGE, M. Introduction to Multi-Agent Systems. Wiley, 2002.\\
Complementar:\\
FERBER, J. Multi-Agent Systems: An Introduction to Distributed
Artificial Intelligence. Addison-Wesley. 1999.\\
WEISS, G. Multiagent Systems: A Modern Approach to Distributed
Artificial Intelligence, Morgan Kauffmann, 2000.

\newpage IDENTIFICAÇÃO

DISCIPLINA: Aprendizagem de Máquina CÓDIGO:\\ 
DEPARTAMENTO: Unidade Acadêmica de Garanhuns\\
ÁREA: \\
CARGA HORÁRIA TOTAL: 60 HORAS NÚMERO DE CRÉDITOS: 04\\
CARGA HORÁRIA SEMANAL: TEÓRICAS: PRÁTICAS: TOTAL: \\
PRÉ-REQUISITOS: Reconhecimento de Padrões\\
CO-REQUISITOS: não há

EMENTA 

Introdução; Vetores de Características; Domínio dos atributos;
Aprendizagem Não Supervisionada; Algoritmos de Agrupamento;
Aprendizagem Supervisionada; Algoritmos de Classificação de padrões;
Algoritmos de regressão de funções; Avaliação de técnicas de
classificação, regressão, agrupamento e testes estatísticos; Tratamento
dos dados; Projeto.

BIBLIOGRAFIA 

Básica:\\
TAN, Pang-Ning; STEINBACH; Michael; KUMAR, Vipin. Introdução ao Data
Mining. Ciência Moderna. 2009.\\
Witten, I. H., Frank, E. (2005). Data mining : practical machine
learning tools and techniques. Second Edition. Elsevier.\\
Mitchell, T. (1997). Machine Learning. McGraw-Hill.\\
Theodoridis. S., Koutroumbas, K. (2009). Pattern Recognition. Fourth
Edition, Academic Press.\\
Duda, R.O., Hart, P.E., Stork, D.G. (2001) Pattern Classification.
Second Edition. Wiley..\\
Complementar:\\
Rezende, Solange O. (2003). Sistemas Inteligentes  Fundamentos e
aplicações. Barueri, SP. Editora Manole. 2003.\\
Marques, Jorge Salvador. (2005). Reconhecimento de Padrões  Métodos
Estatísticos e Neurais. IST Press.\\
Braga, A. P., Ludermir, T.B. e de Carvalho A. P. L. F. (2007). A Redes
Neurais Artificiais  Teoria e Aplicações. LTC.\\
Bishop, C. M. (2006). Pattern Recognition and Machine Learning.
Springer.\\
Haykin, Simon. Redes Neurais: Princípios e Prática, 2. edição. Bookman.
2001.\\
Beale, R.; Jackson, T. Neural Computing  An Introduction. Institute of Physics Publishing. 1990.

\newpage IDENTIFICAÇÃO

DISCIPLINA: Biometria CÓDIGO:\\ 
DEPARTAMENTO: Unidade Acadêmica de Garanhuns\\
ÁREA: \\
CARGA HORÁRIA TOTAL: 60 HORAS NÚMERO DE CRÉDITOS: 04\\
CARGA HORÁRIA SEMANAL: TEÓRICAS: PRÁTICAS: TOTAL: \\
PRÉ-REQUISITOS: Reconhecimento de Padrões\\
CO-REQUISITOS: não há

EMENTA 

Introdução a Biometria. Etapas de um sistema computacional para
econhecimento/verificação dos seguintes elementos: digitais, face, voz, íris, retina, veias, mão, pé, assinaturas e manuscritos. Sistemas de segurança biométricos. Seminário. Projeto.

BIBLIOGRAFIA 

Básica:\\
Jerzy Pejas, Andrzej Piegat. Enhanced Methods in Computer Security,
Biometric and Artificial Intelligence Systems. Springer, 2004.\\
John R. Vacca. Biometric Technologies and Verification Systems.
Butterworth-Heinemann, 2007.
\\Ted Dunstone, Neil Yager. Biometric System and Data Analysis: Design,
Evaluation, and Data Mining. Springer, 2008.\\
Complementar:\\
Artigos recentes da área.

\newpage IDENTIFICAÇÃO

DISCIPLINA: Sistemas de Recomendação CÓDIGO:\\ 
DEPARTAMENTO: Unidade Acadêmica de Garanhuns\\
ÁREA: \\
CARGA HORÁRIA TOTAL: 60 HORAS NÚMERO DE CRÉDITOS: 04\\
CARGA HORÁRIA SEMANAL: TEÓRICAS: PRÁTICAS: TOTAL: \\
PRÉ-REQUISITOS: Inteligência Artificial\\
CO-REQUISITOS: não há

EMENTA 

Função do Sistema de Recomendação. Coleta de Informações. Estratégias de
Recomendação. Descoberta do Conhecimento: tipos de descoberta;
classificação; agrupamento. Filtragem de informações; filtragem baseada
em conteúdo; filtragem colaborativa; filtragem híbrida. Técnicas de
Recomendação. Aplicação e avaliação. Sistemas de Recomendação e IHC.
Processamento de Dados, Classificação, Análise de Cluster, Regras
associativas de recomendação. Exemplos de Sistemas de Recomendação.
Projeto.

BIBLIOGRAFIA 

Básica:\\
RICCI, F., ROKACH, L., SHAPIRA, B., KANTOR, P. Recommender Systems
Handbook. Springer, 2010.\\JANNACH, D., ZANKER, M., FELFERING, A., FRIEDRICH, G. Recommender
Systems: An Introduction. Cambridge, 2011.\\RUSSEL, M. Mining the Social Web: Analyzing Data From Facebook, Twitter,
LinkedIn, and Other Social Media Sites. O'Reilly,
2011..\\
Complementar:\\
UCHYIGIT, G., MA, M. Y. Personalization Techniques and Recommender
Systems. Woeld Scientific Publishing Company.\\MUSIAL, K. Recommender System for online Social Network. Lambert
Academic Publishing, 2009.

\newpage IDENTIFICAÇÃO

DISCIPLINA: Reconhecimento de Padrões II CÓDIGO:\\ 
DEPARTAMENTO: Unidade Acadêmica de Garanhuns\\
ÁREA: \\
CARGA HORÁRIA TOTAL: 60 HORAS NÚMERO DE CRÉDITOS: 04\\
CARGA HORÁRIA SEMANAL: TEÓRICAS: PRÁTICAS: TOTAL: \\
PRÉ-REQUISITOS: Inteligência Artificial e Reconhecimento de Padrões\\
CO-REQUISITOS: não há

EMENTA 

Introdução; Extração de Características; Pré-processamento dos dados; Seleção de Características; Redução da dimensionalidade dos padrões; Algoritmos de Aprendizagem de Máquina; Reconhecimento de Imagens; Outras aplicações; Projeto.

BIBLIOGRAFIA 

Básica:\\
WITTEN, I. H., Frank, E. (2005). Data mining : practical machine
learning tools and techniques. Second Edition. Elsevier.\\
MITCHELL, T. (1997). Machine Learning. McGraw-Hill.\\
THEODORIDIS. S., Koutroumbas, K. (2009). Pattern Recognition. Fourth
Edition, Academic Press.\\
GONZALEZ, Rafael C.; WOODS, Richard E.. Processamento Digital De Imagens
(3 Ed). Pearson Education. 2010.\\
Complementar:\\
HAYKIN, Simon. Redes Neurais: Princípios e Prática, 2. edição. Bookman.
2001.\\
TAN, Pang-Ning; STEINBACH; Michael; KUMAR, Vipin. Introdução ao Data
Mining. Ciência Moderna. 2009.\\
BEALE, R.; Jackson, T. Neural Computing  An Introduction. Institute
of Physics Publishing. 1990.\\
REZENDE, Solange O. (2003). Sistemas Inteligentes  Fundamentos e
aplicações. Barueri, SP. Editora Manole. 2003.\\
BRAGA, A. P., Ludermir, T.B. e de Carvalho A. P. L. F. (2007). A Redes
Neurais Artificiais  Teoria e Aplicações. LTC.\\
MARQUES, Jorge Salvador. (2005). Reconhecimento de Padrões  Métodos
Estatísticos e Neurais. IST Press.\\
DUDA, R.O., Hart, P.E., Stork, D.G. (2001) Pattern Classification.
Second Edition. Wiley.\\
BISHOP, C. M. (2006). Pattern Recognition and Machine Learning.
Springer.\\
A. K. Jain. Fundamentals of digital image processing. Prentice-Hall
International Editions, Englewood Cliffs, NJ, 1989.\\
Bernd Jähne. Digital Image Processing 6th edition. Springer. 2005.

\newpage IDENTIFICAÇÃO

DISCIPLINA: Sistemas Inteligentes Híbridos CÓDIGO:\\ 
DEPARTAMENTO: Unidade Acadêmica de Garanhuns\\
ÁREA: \\
CARGA HORÁRIA TOTAL: 60 HORAS NÚMERO DE CRÉDITOS: 04\\
CARGA HORÁRIA SEMANAL: TEÓRICAS: PRÁTICAS: TOTAL: \\
PRÉ-REQUISITOS: Inteligência Artificial\\
CO-REQUISITOS: não há

EMENTA 

Introdução. Classificação dos sistemas inteligentes
híbridos. Sistemas Fuzzy. Algoritmos Genéticos. Raciocínio baseado em
casos. Sistemas NeuroFuzzy, Projeto Evolucionário de Redes Neurais.
Combinação de Raciocínio baseado em casos e Redes Neurais. Extração de
Regras de Redes Neurais. Projeto.

BIBLIOGRAFIA 

Básica:\\
BRAGA, A P, LUDEMIR, T.B. e CARVALHO, A. Redes Neurais Artificiais 
Teoria e Aplicações, LTC, 2007.\\
GOONATILAKE, S., KHEBBAL, S.  Intelligent Hybrid Systems. John Wiley
\& Sons, Inc. 1995.\\
Complementar:\\
Artigos recentes da área.

\newpage IDENTIFICAÇÃO

DISCIPLINA: Computação Forense CÓDIGO:\\ 
DEPARTAMENTO: Unidade Acadêmica de Garanhuns\\
ÁREA: \\
CARGA HORÁRIA TOTAL: 60 HORAS NÚMERO DE CRÉDITOS: 04\\
CARGA HORÁRIA SEMANAL: TEÓRICAS: PRÁTICAS: TOTAL: \\
PRÉ-REQUISITOS: Inteligência Artificial e Reconhecimento de Padrões\\
CO-REQUISITOS: não há

EMENTA 

Princípios básicos de Ciência Forense e áreas de atuação. Apresentar os
conceitos básicos da perícia criminal e cível. Novas tecnologias
disponíveis nas áreas de Computação Forense e ferramentas tecnológicas
para processamento e análise de evidências. Seminários e Projeto.

BIBLIOGRAFIA 

Básica:\\
DUDA, R.O., HART, P.E., STORK, D.G. Pattern Classification. Second
Edition. Wiley, 2001.\\MARQUES, J. S. Reconhecimento de Padrões  Métodos Estatísticos e
Neurais. IST Press, 2005.\\BISHOP, C. M. Pattern Recognition and Machine Learning. Springer, 2006.\\
Complementar:\\
Artigos recentes da área.

\newpage IDENTIFICAÇÃO

DISCIPLINA: Métodos Paramétricos de Aprendizagem de Máquina CÓDIGO:\\ 
DEPARTAMENTO: Unidade Acadêmica de Garanhuns\\
ÁREA: \\
CARGA HORÁRIA TOTAL: 60 HORAS NÚMERO DE CRÉDITOS: 04\\
CARGA HORÁRIA SEMANAL: TEÓRICAS: PRÁTICAS: TOTAL: \\
PRÉ-REQUISITOS: Reconhecimento de Padrões\\
CO-REQUISITOS: não há

EMENTA 

Reconhecimento de padrões. Revisão de Probabilidade e Estatística. Teoria da decisão Bayesiana. Máxima verossimilhança e estimação de parâmetros. Métodos paramétricos de classificação. Métodos paramétricos de regressão. Métodos paramétricos de agrupamento. Análise dos Componentes Principais. Discriminante linear de Fisher.

BIBLIOGRAFIA 

Básica:\\
Richard Szeliski.Computer Vision: Algorithms and Applications. Springer,
2010.\\
Nicu Sebe, Michael S. Lew. Robust Computer Vision: Theory and
Applications. Kluwer / Springer, 2003.\\
Richard Hartley, Andrew Zisserman. Multiple View Geometry in Computer
Vision, 2nd Edition. Cambridge University Press, 2004.\\
Complementar:\\
Theodoridis, S., Koutroumbas, K. Pattern Recognition. Fourth Edition,
Academic Press, 2009.

\newpage IDENTIFICAÇÃO

DISCIPLINA: Introdução á Bioinformática CÓDIGO:\\ 
DEPARTAMENTO: Unidade Acadêmica de Garanhuns\\
ÁREA: \\
CARGA HORÁRIA TOTAL: 60 HORAS NÚMERO DE CRÉDITOS: 04\\
CARGA HORÁRIA SEMANAL: TEÓRICAS: PRÁTICAS: TOTAL: \\
PRÉ-REQUISITOS: Inteligência Artificial e Reconhecimento de Padrões\\
CO-REQUISITOS: não há

EMENTA 

Introdução aos conceitos de Bio-informática e Biologia Molecular.
Aplicações envolvendo métodos computacionais, matemáticos e
estatísticos. Introdução aos diferentes tipos de dados biológicos,
Bancos de Dados Biológicos e ferramentas público. Alinhamento de
Sequências e Sequenciamento de DNA. Classificação e Anotação de
Sequências Biológicas. Estruturas de Dados Biológicos e Busca em
Cadeias. Transcrição, regulação e expressão gênica. Análise de dados de
microarray. Famílias de Proteínas e Predição de Estruturas. Modelagem
de Sistemas Biológicos. Árvores Evolucionárias e Filogenia.

BIBLIOGRAFIA 

Básica:\\
ALURU, S. Handbook of Computational Molecular Biology. Chapman \& Hall/CRC, 2006.\\
BAXEVANIS, A., OUELLETTE, F. Bioinformatics: A Practical Guide to the Analysis of Genes and Proteins. Wiley-Interscience, 1998.\\
SETUBAL, J. C., MEIDANIS, J. Introduction to Computational Molecular Biology. PSW Publ. Co., 1997.\\
ALBERTS, B. Essential Cell Biology: An Introduction to the Molecular Biology of the Cell. Garland Pub., 1997.\\
Complementar:\\
Artigos recentes da área.

\newpage IDENTIFICAÇÃO

DISCIPLINA: Tópicos Especiais em Bioinformática CÓDIGO:\\ 
DEPARTAMENTO: Unidade Acadêmica de Garanhuns\\
ÁREA: \\
CARGA HORÁRIA TOTAL: 60 HORAS NÚMERO DE CRÉDITOS: 04\\
CARGA HORÁRIA SEMANAL: TEÓRICAS: PRÁTICAS: TOTAL: \\
PRÉ-REQUISITOS: Bioinformática\\
CO-REQUISITOS: não há

EMENTA 

Estudo de técnicas avançadas em Bioinformática permitindo ao aluno conhecer o estado da arte nesta área de pesquisa.

BIBLIOGRAFIA 

Artigos recentes da área de pesquisa.

\newpage IDENTIFICAÇÃO

DISCIPLINA: Tópicos Especiais em Aprendizagem de Máquina. CÓDIGO:\\ 
DEPARTAMENTO: Unidade Acadêmica de Garanhuns\\ 
ÁREA: \\
CARGA HORÁRIA TOTAL: 60 HORAS NÚMERO DE CRÉDITOS: 04\\
CARGA HORÁRIA SEMANAL: TEÓRICAS: PRÁTICAS: TOTAL: \\
PRÉ-REQUISITOS: Aprendizagem de Máquina\\
CO-REQUISITOS: não há

EMENTA 

Estudo de técnicas avançadas em Aprendizagem de Máquina permitindo ao
aluno conhecer o estado da arte nesta área de pesquisa.

BIBLIOGRAFIA 

Básica:\\
IEEE Transactions on Pattern Analysis and Machine Intelligence\\
IEEE Transactions on Neural Networks\\
IEEE Transactions on Knowledge and Data Engineering (Print)\\
Complementar:\\
IEEE Transactions on Evolutionary Computation\\
Evolutionary Computation (Online)\\
ACM Computing Surveys\\
IEEE Transactions on Image Processing\\
Image and Vision Computing\\
Data Mining and Knowledge Discovery (Dordrecht. Online)\\
I.E.E.E. Transactions on Computers (Print)\\
IEEE Intelligent Systems\\
International Journal of Computer Vision\\

\newpage IDENTIFICAÇÃO

DISCIPLINA: Tópicos Especiais em Redes Neurais. CÓDIGO:\\\
DEPARTAMENTO: Unidade Acadêmica de Garanhuns\\ ÁREA: \\
CARGA HORÁRIA TOTAL: 60 HORAS NÚMERO DE CRÉDITOS: 04\\
CARGA HORÁRIA SEMANAL: TEÓRICAS: PRÁTICAS: TOTAL: \\
PRÉ-REQUISITOS: Redes Neurais\\
CO-REQUISITOS: não há

EMENTA 

Estudo de técnicas avançadas em Redes Neurais permitindo ao aluno
conhecer o estado da arte nesta área de pesquisa.

BIBLIOGRAFIA 

Básica:\\
IEEE Transactions on Neural Networks\\
IEEE Transactions on Knowledge and Data Engineering (Print)\\
IEEE Transactions on Pattern Analysis and Machine Intelligence\\
IEEE Transactions on Evolutionary Computation\\
Complementar:\\
Evolutionary Computation (Online)\\
ACM Computing Surveys\\
IEEE Transactions on Image Processing\\
Image and Vision Computing\\
Data Mining and Knowledge Discovery (Dordrecht. Online)\\
I.E.E.E. Transactions on Computers (Print)\\
IEEE Intelligent Systems\\
International Journal of Computer Vision\\

\newpage IDENTIFICAÇÃO

DISCIPLINA: Tópicos Especiais em Computação Evolucionária. CÓDIGO:\\ 
DEPARTAMENTO: Unidade Acadêmica de Garanhuns\\ 
ÁREA: \\
CARGA HORÁRIA TOTAL: 60 HORAS NÚMERO DE CRÉDITOS: 04\\
CARGA HORÁRIA SEMANAL: TEÓRICAS: PRÁTICAS: TOTAL: \\
PRÉ-REQUISITOS: Computação Evolucionária\\
CO-REQUISITOS: não há

EMENTA 

Estudo de técnicas avançadas em Computação Evolucionária permitindo ao
aluno conhecer o estado da arte nesta área de pesquisa.

BIBLIOGRAFIA 

Básica:\\
IEEE Transactions on Evolutionary Computation\\
Evolutionary Computation (Online)\\
ACM Computing Surveys\\
IEEE Transactions on Pattern Analysis and Machine Intelligence\\
IEEE Transactions on Neural Networks\\
Complementar:\\
IEEE Transactions on Knowledge and Data Engineering (Print)\\
IEEE Transactions on Image Processing\\
Image and Vision Computing\\
Data Mining and Knowledge Discovery (Dordrecht. Online)\\
I.E.E.E. Transactions on Computers (Print)\\
IEEE Intelligent Systems\\
International Journal of Computer Vision\\

\newpage IDENTIFICAÇÃO

IDENTIFICAÇÃO 

DISCIPLINA: Tópicos Especiais em Inteligência Artificial CÓDIGO:\\ 
DEPARTAMENTO: Unidade Acadêmica de Garanhuns\\
ÁREA: \\
CARGA HORÁRIA TOTAL: 60 HORAS NÚMERO DE CRÉDITOS: 04\\
CARGA HORÁRIA SEMANAL: TEÓRICAS: PRÁTICAS: TOTAL: \\
PRÉ-REQUISITOS: Inteligência Artificial.\\
CO-REQUISITOS: não há

EMENTA 

Estudo de técnicas avançadas em Inteligência Artificial permitindo ao
aluno conhecer o estado da arte nesta área de pesquisa.

BIBLIOGRAFIA 

Básica:\\
ACM Computing Surveys\\
IEEE Intelligent Systems\\
IEEE Transactions on Pattern Analysis and Machine Intelligence\\
IEEE Transactions on Neural Networks\\
IEEE Transactions on Knowledge and Data Engineering (Print)\\
Complementar:\\
Data Mining and Knowledge Discovery (Dordrecht. Online)\\
Evolutionary Computation (Online)\\
I.E.E.E. Transactions on Computers (Print)\\
IEEE Transactions on Evolutionary Computation\\
IEEE Transactions on Image Processing\\
Image and Vision Computing\\
International Journal of Computer Vision
\newpage

\section*{Área Temática: Matemática e Simulação Computacional}

IDENTIFICAÇÃO

DISCIPLINA: Análise Numérica CÓDIGO:\\ 
DEPARTAMENTO: Unidade Acadêmica de Garanhuns\\
ÁREA: \\
CARGA HORÁRIA TOTAL: 60 HORAS NÚMERO DE CRÉDITOS: 04\\
CARGA HORÁRIA SEMANAL: TEÓRICAS: PRÁTICAS: TOTAL: \\
PRÉ-REQUISITOS: Cálculo para Computação III\\
CO-REQUISITOS: não há

EMENTA 

Conceitos básicos de matemática computacional. Aproximações de funções.
Soluções para sistemas de equações lineares. Integração e derivação
numéricas: método de Gauss-Legendre e um método de cálculo aproximado
de autovalores e autovetores. Equações diferenciais ordinárias:
problemas de valores iniciais.

BIBLIOGRAFIA 

Básica:
PINA, Heitor. Métodos Numéricos. Lisboa. McGraw-Hill, 2004\\
ROSA, Mário. Tópicos de análise Numérica. Coimbra. Universidade de Coimbra.
Departamento de Matemática, 1991.\\
SCHEID, Francis. Análise Numérica. Lisboa. McGraw-Hill, 1991.\\
VALENÇA, Maria Raquel. Análise Numérica. Lisboa. Universidade Aberta, 1996.

\newpage IDENTIFICAÇÃO

DISCIPLINA: Cálculo Numérico e Computacional CÓDIGO:\\
DEPARTAMENTO: Unidade Acadêmica de Garanhuns\\ 
ÁREA:\\
CARGA HORÁRIA TOTAL: 60 HORAS NÚMERO DE CRÉDITOS: 04\\
CARGA HORÁRIA SEMANAL:04 TEÓRICAS: 04 PRÁTICAS: TOTAL: 04\\
PRÉ-REQUISITOS: Cálculo para Computação III\\
CO-REQUISITOS: não há

EMENTA 

Máquinas digitais: precisão, exatidão e erros. Aritmética de ponto
flutuante. Sistemas de enumeração. Sistemas lineares. Resolução
computacional de sistemas de equações lineares. Resolução de equações
transcendentes. Aproximação de funções: interpolação spline,
ajustamento de curvas, aproximação racional e por polinômios de
Chebyschev. Integração numérica: Newton-Cotes e quadratura Gaussiana.

BIBLIOGRAFIA 

Básica:\\
RUGGIERO, M. e LOPES, V.,  Cálculo Numérico: Aspectos
Teóricos e Computacionais , McGraw-Hill, 1996.\\
CLÁUDIO, D. M. e MARINS, J. M.,  Cálculo Numérico
Computacional , Teoria e Prática. São Paulo, Atlas,
1989.\\
CONTE, S. D.,  Elementos de Análise
Numérica , São Paulo, Globo:1977.\\
Complementar:\\
McCRACKEN, D. e DORN, W.,  Cálculo Numérico com
Estudos de Casos em FORTRAN IV , Rio de janeiro :
Campus, 1978.\\
BARROSO, L. C. et al.,  Cálculo Numérico (Com
Aplicações) 2..ed. , São Paulo, Harbra, 1987.

\newpage IDENTIFICAÇÃO

DISCIPLINA: Métodos de Otimização CÓDIGO:\\
DEPARTAMENTO: Unidade Acadêmica de Garanhuns\\
ÁREA: \\
CARGA HORÁRIA TOTAL: 60 HORAS NÚMERO DE CRÉDITOS: 04\\
CARGA HORÁRIA SEMANAL: TEÓRICAS: PRÁTICAS: TOTAL: \\
PRÉ-REQUISITOS: Teoria de Grafos e Análise de Algoritmos\\
CO-REQUISITOS: não há

EMENTA 

Programação Linear, Método Simplex. Programação Não-Linear: convexidade,
otimização sem restrições, otimização com restrições, condições de
otimalidade, métodos computacionais de otimização.

BIBLIOGRAFIA 

Básica:\\
Bazaraa,Sherali,Shetty  Nonlinear Programming, Theory and Algorithms
 John Wiley and Sons.\\
Luenberger,D.G.  Linear and Nonlinear Programming  Addison
Wesley,1984.\\
Dennis, Schnabel  Numerical Methods for Unconstrained Optimization
and Nonlinear Equations Siam.\\
Complementar:\\
Peressini, Sullivan, Uhl The Mathematics of Nonlinear Programming 
Springer Verlag.\\
Bregalda, Oliveira, Bornstein  Introdução à Programação Linear  Ed.
Campus.

\newpage IDENTIFICAÇÃO

DISCIPLINA: Modelagem Matemática CÓDIGO:\\ 
DEPARTAMENTO: Unidade Acadêmica de Garanhuns\\ 
ÁREA: \\
CARGA HORÁRIA TOTAL: 60 HORAS NÚMERO DE CRÉDITOS: 04\\
CARGA HORÁRIA SEMANAL: TEÓRICAS: PRÁTICAS: TOTAL: \\
PRÉ-REQUISITOS: Equações Diferenciais\\
CO-REQUISITOS: não há

EMENTA 

Princípios básicos (o que é um modelo, porque modelar, objetivos e requisitos); Metodologia: etapas (identificação, formulação e solução, modelos matemáticos (quantitativos e qualitativos), tipos de modelos (determinísticos, fuzzy, estatístico, estocástico),modelos discretos e contínuos,processos de modelagem; Noções de cálculo vetorial e tensorial, significado físico dos operadores gradiente, divergente, rotacional e laplaciano; Propriedades físicas; sistemas referências; leis de conservação, equações constitutivas; Exemplos envolvendo todas as etapas de modelagem (exceto a solução). 

BIBLIOGRAFIA 

Básica: \\
C.L. Dym \& E.S. Ivey  Principles of Mathematical Modeling, Academic
Press, 1980.\\
Karam F., J. e Almeida, R. C., Introdução à Modelagem Matemática, Notas
impressas Pós-Graduação, LNCC, 2003.\\
T.L. Saaty \& J.M. Alexander, Thinking with Models  Mathematical
Models in Physical, Biological and Social Sciences, Pergamon Press, 1981.\\
Complementar:\\
R.B. Bird, W.E. Stewart \& E.N. Lightfoot, Transport Phenomena, John
Wiley \& Sons, 1960.\\
Mathematical Modelling Techniques, Rutherford Aris, Dover, 1994.\\
Introduction to Continuum Mechanics, W. M. Lai, D. Rubin, E.
Krempl,Pergamon Press, 1974.

\newpage IDENTIFICAÇÃO

DISCIPLINA: Programação Matemática CÓDIGO:\\
DEPARTAMENTO: Unidade Acadêmica de Garanhuns\\
ÁREA: \\
CARGA HORÁRIA TOTAL: 60 HORAS NÚMERO DE CRÉDITOS: 04\\
CARGA HORÁRIA SEMANAL: TEÓRICAS: PRÁTICAS: TOTAL: \\
PRÉ-REQUISITOS: Álgebra Linear \\
CO-REQUISITOS: não há

EMENTA 

Modelagem sistêmica de problemas industriais. Modelos de Programação Linear Inteira Mista (PLIM) para apoio à tomada de decisão. Programação Linear (PL). Método primal simplex. Problema de transporte. Problema de designação. Dualidade. Método dual simplex. Análise de sensibilidade. Interpretação econômica da PL. Programação inteira. Programação inteira mista. Resolução de problemas de grande porte. Decomposição em PL e PLIM. Aplicações em sistemas produtivos.

BIBLIOGRAFIA 

Básica:
J. F. Benders. Partitioning procedures for solving mixed-variables
programming problems. Numerisch Mathematik, v. 4, p. 238-252, 1962. \\
G. B. Dantzig and Wolfe. Decomposition principle for linear programs.
Operations Research, v. 8, p. 101-111, 1960. \\
C. R. V. de Carvalho. Notas de Aula\\
C. R. V. de Carvalho. Une Proposition d'Integration de
la Planification et l'Ordonancement de Production:
Application de la Métode de Benders. PhD thesis, Université Blaise
Pascal, Clermont-Ferrand, França, 1998. \\
M. T. P. de Carvalho. Confecção de horários de aulas em instituições de
ensino privadas. Master's thesis, Programa de
Pós-Graduação em Engenharia de Produção da UFMG, 2002. \\
M. C. Goldbarg and H. P. L. Luna. Otimização Combinatória e Programação
Linear: Modelos e Algoritmos. Ed. Campus, 2000. \\
F. S. Hiller and G. J. Liberman. Introdução à Pesquisa Operacional. Ed.
Campus Ltda, Rio de Janeiro, 1989. \\
L. S. Lasdon. Optimization Theory for Large Systems. The Macmillan
Company, New York, 1972. \\
T. R. Neto. Uma metodologia para elaboração de planos de compras de
carvão em empresas siderúrgicas brasileiras. Dissertação de Mestrado.
Programa de Pós-Graduação em Engenharia de Produção da UFMG, 2003. \\
C. R. Oliveira. Planejamento da distribuição de produtos siderúrgicos
utilizando modelos de localização. Dissertação de Mestrado. Programa de
Pós-Graduação em Engenharia de Produção da UFMG, 2003. \\
H. M. Wagner. Pesquisa Operacional. Prentice-Hall do Brasil, Rio de
Janeiro, 1986.\\
R. Baker. Introduction to Sequencing and Scheduling. Wiley, 1974.\\
M. S. Bazaraa and J. J. Jarvis. Linear Programming and Network Flows.
John Wiley \& Sons, New York, 1977. \\

\newpage IDENTIFICAÇÃO

DISCIPLINA: Simulação de Sistemas CÓDIGO:\\ 
DEPARTAMENTO: Unidade Acadêmica de Garanhuns\\ ÁREA: \\
CARGA HORÁRIA TOTAL: 60 HORAS NÚMERO DE CRÉDITOS: 04\\
CARGA HORÁRIA SEMANAL: TEÓRICAS: PRÁTICAS: TOTAL: \\
PRÉ-REQUISITOS: Equações Diferenciais\\
CO-REQUISITOS: não há

EMENTA 

Introdução à simulação. Propriedades e classificação dos modelos de
simulação. Revisão de conceitos: estatística, probabilidade, processos
estocásticos. Exemplos de sistemas de simulação. Geração de números
aleatórios. Noções básicas em teoria dos números. Geração e teste.
Distribuições clássicas contínuas e discretas. Simulação de sistemas
discretos e de sistemas contínuos. Verificação e validação de modelos.
Técnicas estatísticas para análise de dados e de resultados de modelos
de simulação. Simulação de sistemas simples de filas. Simulação de
sistemas de computação. Estudos de caso.

BIBLIOGRAFIA 

Freitas Filho, Paulo José. Introdução a modelagem e simulação de
sistemas. Ed. VisualBooks, Florianópolis, 2001.\\
Law, A.M. e Kelton, W.D. Simulation Modeling and Analysis. Ed.
McGraw-Hill, USA, 1991.\\
Perin Filho, C. Introdução a simulação de Sistemas. Ed. da Unicamp,
Campinas, 1995.\\
Prado, Darci. Teoria das Filas e da Simulação. Editora DG, Belo
Horizonte (MG), 1999.\\
Prado, Darci. Usando o Arena em Simulação. Editora DG, Belo Horizonte
(MG), 1999.

\newpage IDENTIFICAÇÃO

DISCIPLINA: Tópicos em Matemática Computacional CÓDIGO:\\ 
DEPARTAMENTO: Unidade Acadêmica de Garanhuns\\ 
ÁREA:\\
CARGA HORÁRIA TOTAL: 60 HORAS NÚMERO DE CRÉDITOS: 04\\
CARGA HORÁRIA SEMANAL: TEÓRICAS: PRÁTICAS: TOTAL:\\
PRÉ-REQUISITOS: Calculo para Computação III\\
CO-REQUISITOS: não há

EMENTA 

Revisão de estatística: Conceitos clássicos de estatística
(probabilidade, distribuições de probabilidade, médias, variância).
Introdução ao pacote computacional STATISTICA. Revisão de álgebra
matricial: Operações com matrizes e vetores. Valores característicos e
vetores característicos de matrizes quadradas. Forma canônica de
matrizes. Solução de um sistema de equações lineares e análise da
solução. Introdução ao pacote computacional MATLAB. Revisão de métodos
numéricos: Métodos de busca de zeros de funções. Interpolação de dados
experimentais. Integração numérica. Resolução numérica de equações
diferenciais. Revisão de equações diferenciais: Solução analítica de
EDOs. Solução de EDOs de 1. e 2. ordem. Noções básicas de EDPs.

BIBLIOGRAFIA 

BUSSAB, W.; MORETTIN, P.A. Estatística Básica. 5. ed. São Paulo:
Saraiva, 2004.\\
CHAPMAN, S.J. Fortran 90/95 for Scientists and Engineers. Mc Graw Hill,
1997.\\
CHAPMAN, S.J. Matlab Programming for Engineers. 2nd ed. Brooks Cole,
2001.\\
CHAPRA, S.C.; CANALE, R.P. Numerical Methods for Engineers with
Programming and Software Applications. 4th Ed. Mc Graw Hill, 2002.\\
SPIEGEL, M.R. Manual de Fórmulas e Tabelas de Matemática. 2. Ed. Versão
rev. E ampl. Makron Books, 2002.\\
STEINBRUCH, A.; WINTERLE, P. Álgebra Linear. 2a Ed. São Paulo: McGraw
Hill, 1987.\\
WYLIE, C.R.; BARRET, L.C. Advanced Engineering Mathematics. 5th ed.
McGraw Hill, 1982.\\
ZILL, D.G.; CULLEN, M.R. Equações Diferenciais -- Vol. 1 e Vol. 2, 3.
Ed. São Paulo: Makron Books, 2001.\\
WONNACOTT, H. and J. WONNACOTT, J., Introdução à Estatística; Rio de
Janeiro, RJ: Livros Técnicos e Científicos; 1980.

\newpage IDENTIFICAÇÃO

DISCIPLINA: Cálculo para Computação III CÓDIGO:\\ 
DEPARTAMENTO: Unidade Acadêmica de Garanhuns\\
ÁREA: \\
CARGA HORÁRIA TOTAL: 60 HORAS NÚMERO DE CRÉDITOS: 04\\
CARGA HORÁRIA SEMANAL: TEÓRICAS: PRÁTICAS: TOTAL: \\
PRÉ-REQUISITOS: Cálculo para Computação II\\
CO-REQUISITOS: não há

EMENTA 

Conceitos introdutórios e classificação das equações diferenciais. Equações diferenciais de primeira ordem. Obtenção de soluções de equações lineares, separáveis, exatas, não exatas com fatores integrantes simples, etc. Algumas aplicações das equações de primeira ordem. Equações diferenciais de segunda ordem, propriedades gerais das soluções, soluções das homogêneas com coeficientes constantes. Equações lineares não homogêneas, método dos coeficientes a determinar e método da variação dos parâmetros. Estudo introdutório das oscilações lineares livres e forçadas. Transformada de Laplace, propriedades fundamentais, e utilização para resolução de equações diferenciais. Equação do calor. Método de separação de variáveis. Séries de Fourier, propriedades básicas e aplicações. Equação da onda, vibrações em uma corda elástica. Equaçãode Laplace.

BIBLIOGRAFIA 

Básica:\\
Equações Diferenciais Elementares e Problemas de Valores de
Contorno, de BOYCE \& DiPRIMA, Ed. Guanabara Dois (Sétima Edição).

\newpage IDENTIFICAÇÃO

DISCIPLINA: Tópicos em Pesquisa Operacional CÓDIGO:\\ 
DEPARTAMENTO: Unidade Acadêmica de Garanhuns\\ 
ÁREA: \\
CARGA HORÁRIA TOTAL: 60 HORAS NÚMERO DE CRÉDITOS: 04\\
CARGA HORÁRIA SEMANAL: TEÓRICAS: PRÁTICAS: TOTAL: \\
PRÉ-REQUISITOS: Teoria de Grafos\\
CO-REQUISITOS: não há

EMENTA 

Método dos quadrados mínimos: modelos de programação linear; problema da análise de atividades; problema da dieta; problema do transporte; problema da designação; solução gráfica; limitações da programação linear. Método Simplex. Algoritmos especiais: problema do transporte; problema da designação. Dualidade. Análise pós-otimização. Noções de algoritmos Genéticos. Funções de várias variáveis. Método de otimização sem restrição. Método de otimização com restrição de igualdade. Método de otimização com restrição de desigualdade. 

BIBLIOGRAFIA 

Básica:\\
R. J. Vanderbei, Linear Programming  Foundations and Extensions.
Kluwer Academics\\
Publishers, Boston 1996.\S. J. Wright, Primal-Dual Interior-Point Methods. SIAM Publications,
Philadelphia 1997. 
\newpage

\section*{Área Temática: Metodologia e Técnicas da Computação}

IDENTIFICAÇÃO

DISCIPLINA: Analise Quantitativa Estatística CÓDIGO: \\
DEPARTAMENTO: Unidade Acadêmica de Garanhuns \\
ÁREA: \\
CARGA HORÁRIA TOTAL: 60 horas NÚMERO DE CRÉDITOS: 04\\
CARGA HORÁRIA SEMANAL: TEÓRICAS: 50h PRÁTICAS: 10h TOTAL: 60h\\
PRÉ-REQUISITOS: Probabilidade e Estatística, Metodologia Científica\\
CO-REQUISITOS: não há

EMENTA 

Conceitos de Probabilidade, Análise Exploratória de Dados, Inferência
Estatística Paramétrica, Inferência Estatística Não Paramétrica,
Regressão Linear, Análise de Variância, Simulação Estocástica.

BIBLIOGRAFIA 

Básica:\\
Estatística Aplicada e Probabilidade para Engenheiros, Douglas C.
Montgomery e George C. Runger. Quarta Edição. Editora LTC. (LIVRO
TEXTO) \\
Estatística Básica. Probabilidade e Inferência. Volume Único. Luiz
Gonzaga Morettin. Person. (LIVRO DE EXERCÍCIOS)\\
Applied Nonparametric Statistics  Wayne W. Daniel, Duxbury Resource
Center, Second Edition.\\
Complementar:\\
The Art of Computer System Performance Analysis: Techniques for
Experimental Design, Measurement, Simulation and Modeling, Raj Jain,
John Wiley \& Sons, 1991, ISBN: 0-471-50336-3. \\
Statistics for Experimenters: An Introduction to Design, Data Analysis,
and Model Building , George E. P. Box, Wiliam G. Hunter, J. Stuart
Hunter, John Wiley \& Sons, Inc. 1978.

\newpage IDENTIFICAÇÃO

DISCIPLINA: Métodos de Pesquisa em Computação CÓDIGO: \\
DEPARTAMENTO: Unidade Acadêmica de Garanhuns \\
ÁREA: \\
CARGA HORÁRIA TOTAL: 60 horas NÚMERO DE CRÉDITOS: 04\\
CARGA HORÁRIA SEMANAL: TEÓRICAS: 50h PRÁTICAS: 10h TOTAL: 60h\\
PRÉ-REQUISITOS: Metodologia Científica\\
CO-REQUISITOS: não há

EMENTA 

A metodologia científica é o estudo de como se conduz/produz pesquisa
científica. Metodologia científica é necessária, entre outras razões,
para tornar os resultados da pesquisa mais confiáveis e possíveis de
serem reproduzidos, de forma independente, por outros pesquisadores.
Este curso irá apresentar estratégias e métodos para pesquisa em
computação desde a formulação do problema até a validação de uma
possível solução. Em particular, o curso irá focar em métodos
experimentais e explorar o papel da experimentação na pesquisa em
computação.

BIBLIOGRAFIA 

Magne Jørgensen's site:
http://simula.no/people/magnej/bibliography?b\_size:int=9999999\&b\_start:int=0\&-C=\\
ESERNET  the Experimental Software Engineering Network\\
Steve Easterbrook's course CSC2130S: Empirical Research
Methods in Software Engineering, at University of Toronto;\\
Susan Sim's course ICS 280: Research Methodology for
Software at UC Irvine;\\
Dewayne Perry's course EE382C Empirical Studies in
Software Engineering at U Texas;\\
Mary Shaw's course 17-939A What Makes Good Research in
Software Engineering at CMU;\\
Jim Herbsleb's course 17-810 Empirical Methods in
Software Engineering Research at CMU;\\
Wilhelm Hasselbring's course 2.01.261 Research Methods
in Software Engineering at U Oldenburg;\\
Philip Johnson's Readings in Empirical Evaluation for
Budding Software Engineering Researchers at U Hawaii;\\
Letizia Jaccheri's Empirical software engineering
course at NTNU Trondheim;\\
Andreas Zeller's course Empirical Software Engineering
at Saarland U.\\
Jonathan Sillito. Qualitative Research Methods in Software Engineering
(CPSC 601.23), Winter 2007, 2008 and 2009.\\
Ioannidis: Why Most Published Research Findings Are False.\\
Lung, J., Aranda, J., Easterbrook, S. M. and Wilson, G. V., On the
Difficulty of Replicating Human Subjects Studies in Software
Engineering. 30th ACM/IEEE International Conference on Software
Engineering (ICSE 2008), Leipzig, Germany, May 10-18,
2008.

\newpage IDENTIFICAÇÃO

DISCIPLINA: Análise de Dados Simbólicos CÓDIGO: \\
DEPARTAMENTO: Unidade Acadêmica de Garanhuns \\
ÁREA: \\
CARGA HORÁRIA TOTAL: 60 horas NÚMERO DE CRÉDITOS: 04\\
CARGA HORÁRIA SEMANAL: TEÓRICAS: 50h PRÁTICAS: 10h TOTAL: 60h\\
PRÉ-REQUISITOS: Probabilidade e Estatística, Metodologia Científica\\
CO-REQUISITOS: não há

EMENTA 

Introdução a Análise de Dados Simbólicos, Dados Simbólicos, Estatística
Descritiva para Dados Simbólicos, Similaridade e Dissimilaridade entre
Dados Simbólicos, Análise de Componentes Principais de Dados
Simbólicos, Regressão para Dados Simbólicos, Classificação para Dados
Simbólicos, Análise de Agrupamento para Dados Simbólicos, Previsão de
Séries Temporais de Dados Simbólicos.

BIBLIOGRAFIA 

Applied Multivariate Data Statistical Analysis, Richard Johnson and Dean
W. Wichern. Third Edition, Prentice Hall, 1992.\\
Analysis of Symbolic Data. H.-H. Bock and E. Diday (Eds.),
Springer-Verlag, 2000.\\
Symbolic Data Analysis. Conceptual Statistics and Data Mining. L.
Billard and E. Diday, Wiley, 2006\\
Symbolic Data Analysis and the SODAS Software. E. Diday and M.
Noirhomme-Fraiture, Wiley 2008.

\newpage IDENTIFICAÇÃO

DISCIPLINA: Métodos de Pesquisa Qualitativa CÓDIGO: \\
DEPARTAMENTO: Unidade Acadêmica de Garanhuns\\
ÁREA: \\
CARGA HORÁRIA TOTAL: 60 horas NÚMERO DE CRÉDITOS: 04\\
CARGA HORÁRIA SEMANAL: TEÓRICAS: 50h PRÁTICAS: 10h TOTAL: 60h\\
PRÉ-REQUISITOS: Metodologia Científica\\
CO-REQUISITOS: não há

EMENTA 

Pergunta de Pesquisa, Estudo de Caso, Protocolos de Pesquisa, Métodos de
Coleta de Dados, Análise e Síntese.

BIBLIOGRAFIA 

Kitchenham et al: Preliminary guidelines for empirical research in
software engineering, IEEE Transactions On Software Engineering, Vol.
28, No. 8, August 2002.\\
D. I. K. Sjøberg, T. Dybå, and M. Jørgensen. The Future of Empirical
Methods in Software Engineering Research, In: Future of Software
Engineering (FOSE 2007), 2007.\\
Hannay J E, Sjoberg D I K, Dyba T. A systematic review of theory use in
software engineering experiments. IEEE Transactions on Software
Engineering VOL. 33, NO. 2, FEBRUARY 2007.\\
M. Jørgensen, and D. I. K. Sjøberg. Generalization and Theory-Building
in Software Engineering Research, In: Empirical Assessment in Software
Engineering (EASE2004), ed. by unknown, pp. 29-36, IEEE Proceedings
(ISBN: 0 86341 435 4), 2004.

\newpage IDENTIFICAÇÃO

DISCIPLINA: Revisões Sistemáticas de Literatura CÓDIGO: \\
DEPARTAMENTO: Unidade Acadêmica de Garanhuns \\
ÁREA: \\
CARGA HORÁRIA TOTAL: 60 horas NÚMERO DE CRÉDITOS: 4\\
CARGA HORÁRIA SEMANAL: 4h TEÓRICAS: 50h PRÁTICAS: 10h TOTAL: 60h\\
PRÉ-REQUISITOS: Metodologia Científica\\
CO-REQUISITOS: não há

EMENTA 

Pergunta de Pesquisa, Estratégias de Busca, Termos de Pesquisa,
Protocolo da Revisão, Selecionando as Fontes, Mapeamento Sistemáticos,
Extração de Dados, Desenvolvendo os relatórios de revisão.

BIBLIOGRAFIA 

Arlene Fink. Conducting Research Literature Reviews. Sage Publications,
2010.\\Mark Petticrew and Helen Roberts. Systematics Review in the Social
Scienced. Blackwell Publishing, 2006.\\
B. Kitchenham, T. Dybå, and M. Jørgensen. Evidence-based Software
Engineering, In: International Conference on Software Engineering
(ICSE 2004), pp. 273-281, Edinburgh, IEEE Computer
Society, Washington DC, USA, 2004. \\
T. Dybå, B. Kitchenham, and M. Jørgensen. Evidence-based Software
Engineering for Practitioners, IEEE Software 22(1):58-65, 2005.\\
Magne Jørgensen, Tory Dyb, e Barbara Kitchenham Teaching Evidence-Based
Software Engineering to University Students.\\
Briony J Oates, Graham Capper. Using systematic reviews and
evidence-based software engineering with masters students. EASE 2009.\\
Barbara Kitchenham and Stuart Charters. Guidelines for performing
Systematic Literature Reviews in Software Engineering. Keele University
and Durham University Joint Report. 2007.\\
Travassos e Biolchini. Revisões Sistemáticas Aplicadas a Engenharia de
Software. Tutorial SBES 2007.\\
Andreas Jedlitschka and Marcus Ciolkowski. Reporting Experiments in
Software Engineering.\\
Basili and Rombach. The Goal Question Metric Paradigm.\\
Basili and Caldiera. Improving Software Qualitiy by Reusing Knowledge
and Experience.\\
Miller, J. Can results from software engineering experiments be safely
combined? Software Metrics Symposium, 1999. \\
Jedlitschka, A.; Ciolkowski, M. Towards evidence in software
engineering. Empirical Software Engineering, 2004.
ISESE 2004. Proceedings. 2004 International Symposium on
Volume , Issue , 19-20 Aug. 2004 Page(s): 261  270.\\
Pickard L.M.1;Kitchenham B.A.; Jones P.W. Combining empirical results in
software engineering. Authors: Information and Software Technology, vol
40, num. 14.\\
Tore Dyba, Vigdis By Kampenes, Dag I.K. Sjøberg. A systematic review of
statistical power in software engineering experiments.\\
Shull, F. J., Carver, J. C., Vegas, S. and Juristo, N., The role of
replications in Empirical Software Engineering. Empirical Software
Engineering 13 (2), p211-218. April 2008 \\
Kitchenham, B. The Role of Replications in Software Engineering  A
word of Warning. Empirical Software Engineering 13 (2), p219-221. April
2008.\\
Lung, J., Aranda, J., Easterbrook, S. M. and Wilson, G. V., On the
Difficulty of Replicating Human Subjects Studies in Software
Engineering. 30th ACM/IEEE International Conference on Software
Engineering (ICSE 2008), Leipzig, Germany, May 10-18,
2008.\\
Jedlitschka, A.; Pfahl, D. Reporting guidelines for controlled
experiments in software engineering. 2005 International Symposium on
Empirical Software Engineering.\\
Evaluating guidelines for empirical software engineering studies.
Proceedings of the 2006 ACM/IEEE international symposium on Empirical
software engineering.

\newpage IDENTIFICAÇÃO

DISCIPLINA: Tópicos Especiais em Metodologia de Pesquisa CÓDIGO: \\
DEPARTAMENTO: Unidade Acadêmica de Garanhuns \\
ÁREA: \\
CARGA HORÁRIA TOTAL: 60 horas NÚMERO DE CRÉDITOS: 04\\
CARGA HORÁRIA SEMANAL:04 horas TEÓRICAS: 60h PRÁTICAS: 0h TOTAL: 60h\\
PRÉ-REQUISITOS: Metodologia Cientifica\\
CO-REQUISITOS: não há

EMENTA 

Estudo de técnicas avançadas em na área multidisciplinar de Metodologia
de Pesquisa permitindo ao aluno conhecer o estado da arte nesta área de
pesquisa.

BIBLIOGRAFIA 

Artigos da area.
\newpage

\section*{Área Temática: Mídia e Interação}

IDENTIFICAÇÃO

DISCIPLINA: Processamento Digital de Imagens CÓDIGO:\\ 
DEPARTAMENTO: Unidade Acadêmica de Garanhuns\\
ÁREA: \\
CARGA HORÁRIA TOTAL: 60 HORAS NÚMERO DE CRÉDITOS: 04\\
CARGA HORÁRIA SEMANAL: TEÓRICAS: PRÁTICAS: TOTAL: \\
PRÉ-REQUISITOS: Matemática Discreta e Computação Gráfica\\
CO-REQUISITOS: não há

EMENTA 

Introdução. Amostragem. Processamento de histogramas. Filtragem
espacial. Filtragem no domínio da freqüência. Restauração. Modelos de
cor. Processamento de imagens coloridas. Processamento morfológico.
Segmentação. Representação e Descrição. Tópicos sobre formatos de
arquivo de imagens. Tópicos sobre compressão de imagens. Projeto.

BIBLIOGRAFIA 

Básica:\\
Gonzalez, R.C. \& Woods, R.E. Digital Image Processing, 3rd Ed.
Prentice Hall, USA, 2008.\\
Computação Gráfica: Teoria e Prática, AZEVEDO, E. e Conci, A. Editora
Campus, Elsevier, 2003. Rio de Janeiro.ISBN 8535212533.\\
Facon J. Morfologia Matemática: Teoria e
Exemplos, Ed. Univ. Champagnat, 1996.\\
Complementar:\\
Pratt, W. K.  Digital Image Processing, 4th ed. USA, Wiley
Interscience Pub., 2007.\\
Análise De Imagens Digitais: Principios, Algoritmos E Aplicaçoes,
PEDRINI, HELIO; SCHWARTZ; WILLIAM ROBSON THOMSON PIONEIRA, 2007. ISBN
8522105952.\\
Marques Filho O., Vieira Neto H., Processamento Digital
de Imagens, Brasport Livros, 1999.

\newpage IDENTIFICAÇÃO

DISCIPLINA: Processamento Digital de Imagens II CÓDIGO:\\ 
DEPARTAMENTO: Unidade Acadêmica de Garanhuns\\
ÁREA: \\
CARGA HORÁRIA TOTAL: 60 HORAS NÚMERO DE CRÉDITOS: 04\\
CARGA HORÁRIA SEMANAL: TEÓRICAS: PRÁTICAS: TOTAL: \\
PRÉ-REQUISITOS: Processamento Digital de Imagens\\
CO-REQUISITOS: não há

EMENTA 

Wavelets. Compressão de Imagens. Formatos de arquivos de imagens.
Segmentação. Representação e Descrição. Reconhecimento de objetos.
Tópicos Avançados. Projeto.

BIBLIOGRAFIA 

Básica:\\
James Murray D., William vanRyper. Encyclopedia of Graphics File Formats
2nd Edition. O'Reilly \& Associates, Inc. 1996
(http://www.fileformat.info/mirror/egff/index.htm)\\
Michael W. Frazier. An Introduction to Wavelets Through Linear Algebra.
Springer. 1999.\\
Witten, Moffat, Bell. Managing gigabytes compressing and indexing
documents and images, 2ed.,(1999).\\
Complementar:\\

Gonzales \& Woods. Processamento Digital De Imagens (3 Ed). Pearson
Education. 2010.\\
A. K. Jain. Fundamentals of digital image processing. Prentice-Hall
International Editions, Englewood Cliffs, NJ, 1989.\\
Bernd Jähne. Digital Image Processing 6th edition. Springer. 2005.

\newpage IDENTIFICAÇÃO

DISCIPLINA: Modelagem Geométrica CÓDIGO:\\ 
DEPARTAMENTO: Unidade Acadêmica de Garanhuns\\ 
ÁREA: \\
CARGA HORÁRIA TOTAL: 60 HORAS NÚMERO DE CRÉDITOS: 04\\
CARGA HORÁRIA SEMANAL: TEÓRICAS: PRÁTICAS: TOTAL: \\
PRÉ-REQUISITOS: Computação Gráfica\\
CO-REQUISITOS: não há

EMENTA 

Geometria projetiva. Curvas paramétricas, composição de curvas,
superfícies paramétricas, modelagem sólida, particionamento espacial,
malhas poligonais e níveis de detalhes.

BIBLIOGRAFIA 

Básica:\\
Farin, G.  Curves and Surfaces for CAGD. 4th edition. Academic
Press.\\
Farin, G. Nurbs: From Projective Geometry to Practical Use. 2nd
edition. AK Peters.\\
Gomes, J. e Velho, L. Fundamentos da Computação Gráfica  SBM.\\
Complementar:\\
Mortenson, M. E.  Geometric Modeling. 2nd edition. Wiley and Sons.\\
Luebke, D. et al  Level of Detail for 3D Graphics. Morgan Kaufmann.

\newpage IDENTIFICAÇÃO

DISCIPLINA: Realidade Virtual e Aumentada CÓDIGO:\\ 
DEPARTAMENTO: Unidade Acadêmica de Garanhuns\\ 
ÁREA: \\
CARGA HORÁRIA TOTAL: 60 HORAS NÚMERO DE CRÉDITOS: 04\\
CARGA HORÁRIA SEMANAL: TEÓRICAS: PRÁTICAS: TOTAL: \\
PRÉ-REQUISITOS: Computação Gráfica\\
CO-REQUISITOS: não há

EMENTA 

Introdução  Histórico. Aplicações: mercados tradicionais e
emergentes. Tecnologias Básicas. Definições e Caracterizações. Fatores
Humanos, Percepção Humana e Interação. Princípios Básicos de Computação
Gráfica aplicados a RV e RA. Princípios de Modelagem Geométrica
Aplicados a RV e RA. Modelagem de Ambientes Virtuais. Ferramentas de
Desenvolvimento de Ambientes Virtuais. Tópicos Especiais em Realidade
Virtual.

BIBLIOGRAFIA 

Básica:\\
Desenvolvimento De Jogos 3d E Aplicações Em Realidade Virtual, AZEVEDO,
EDUARDO; STELKO, MICHELLE; MEYER, HOMERO. CAMPUS, 2005. ISBN:
8535215697.\\
Grigore C. Burdea et al. Virtual Reality Technology, 2nd. edition,
Wiley-Interscience, 2003.\\
Oliver Bimber et al. Spatial Augmented Reality: Merging Real and Virtual
Worlds, A K Peters, 2005.\\
Complementar:\\
Oliver Bimber et al. Spatial Augmented Reality: Merging Real and Virtual
Worlds, A K Peters, 2005.\\
MCCARTH, Martin, et al. Reality Architecture: Building 3D Worlds with
Java and VRML. Hertfordshire: Prentice-Hall, 1998.

\newpage IDENTIFICAÇÃO

DISCIPLINA: Processamento Digital de Sinais CÓDIGO:\\
DEPARTAMENTO: Unidade Acadêmica de Garanhuns\\
ÁREA:\\
CARGA HORÁRIA TOTAL: 60 HORAS NÚMERO DE CRÉDITOS: 04\\
CARGA HORÁRIA SEMANAL: TEÓRICAS: PRÁTICAS: TOTAL:\\
PRÉ-REQUISITOS: Calculo II e Reconhecimento de Padrões\\
CO-REQUISITOS: não há

EMENTA 

Representação digital de sinais de áudio, imagens, e vídeo: amostragem,
quantização e  aliasing. Transformada Discreta de Fourier e FFT (1D, 2D e 3D). Outras
transformações: Transformada de Fourier (Contínua), Transformada do
Coseno Discreta, Transformada z, Transformada de Walsh-Hadamard,
Transformada de Haar. Convolução linear, circular e secionada. Filtros
lineares (FIR) e Filtros recursivos (IIR). Aplicações de filtros:
suavização, interpolação, realce, detecção de bordas e segmentação.
Espaço de transformação no tempo e no espaço, localização e efeitos no
espectro. Bancos de filtros e técnicas de análise-ressíntese.
Compressão: Predição Linear, compressão usando DCT, Compensação de
Movimento. Sinais aleatórios: Representação, Filtros de Wiener e de
Kalman.

BIBLIOGRAFIA 

Básica:\\
S. Allen Broughton and Kurt M. Bryan. Discrete Fourier Analysis and
Wavelets: Applications to Signal and Image Processing.
Wiley-Interscience, 2008.\\
John W. Woods. Multidimensional Signal, Image and Video Processing and
Coding. Academic Press, 2006.\\
Rafael C. Gonzales and Richard E. Woods. Digital Image Processing, 3rd
ed. Prentice Hall, 2007.\\
Complementar:\\
Charles L. Byrne. Signal Processing: a Mathematical Approach. A. K.
Peters Ltd., 2005.\\
Ronald N. Bracewell. Fourier Analysis and Imaging. Springer, 2004.\\
Richard W. Hamming. Digital Filters, 3rd ed. Dover Publications, 1997.\\
Alan. V. Oppenheim and Ronald W. Schafer. Discrete-Time Signal
Processing, 2nd ed. Prentice Hall, 1999.

\newpage IDENTIFICAÇÃO

DISCIPLINA: Processamento Digital de Áudio CÓDIGO:\\ 
DEPARTAMENTO: Unidade Acadêmica de Garanhuns\\
ÁREA: \\
CARGA HORÁRIA TOTAL: 60 HORAS NÚMERO DE CRÉDITOS: 04\\
CARGA HORÁRIA SEMANAL: TEÓRICAS: PRÁTICAS: TOTAL: \\
PRÉ-REQUISITOS: Calculo 2\\
CO-REQUISITOS: não há

EMENTA 

Conceitos básicos de som, acústica e psicoacústica. Sinais de áudio:
fundamentos, representação, processamento e síntese. Software para
processamento de áudio. Protocolos e padrões de áudio. Compressão de
áudio. Tópicos Avançados. Projetos.

BIBLIOGRAFIA 

Básica:\\
Curtis Roads. The Computer Music Tutorial. MIT Press, 1996\\
Ken C. Pohlman. Principles of Digital Audio 4th Edition. McGraw Hill.
2000.\\
Charles Dodge \& Thomas A. Jerse. Computer Music: Synthesis,
Composition, and Performance. Schirmer Books. 1997.\\
Complementar:\\
Charles L. Byrne. Signal Processing: a Mathematical Approach. A. K.
Peters Ltd., 2005.\\
Richard W. Hamming. Digital Filters, 3rd ed. Dover Publications, 1997.\\
Alan. V. Oppenheim and Ronald W. Schafer. Discrete-Time Signal
Processing, 2nd ed. Prentice Hall, 1999.

\newpage IDENTIFICAÇÃO

DISCIPLINA: Desenvolvimento de Jogos CÓDIGO:\\ 
DEPARTAMENTO: Unidade Acadêmica de Garanhuns\\ 
ÁREA: \\
CARGA HORÁRIA TOTAL: 60 HORAS NÚMERO DE CRÉDITOS: 04\\
CARGA HORÁRIA SEMANAL: TEÓRICAS: PRÁTICAS: TOTAL: \\
PRÉ-REQUISITOS: Computação Gráfica\\
CO-REQUISITOS: não há

EMENTA 

Introdução: história e categorias de jogos. Gerenciamento de equipes de
desenvolvimento de jogos. Projeto de jogos: roteiro, interface.
Conceitos gráficos: modelagem 2D e 3D. Técnicas e Ferramentas de
Implementação (2D e 3D). Conceitos: gráficos, sons, inteligência
artificial e redes em Jogos. Tópicos complementares. Projeto de
Desenvolvimento.

BIBLIOGRAFIA 

Básica:\\
Desenvolvimento De Jogos 3d E Aplicações Em Realidade Virtual, AZEVEDO,
EDUARDO; STELKO, MICHELLE; MEYER, HOMERO. CAMPUS, 2005. ISBN:
8535215697.\\
Game Design: Principles, Practice, And Techniques- The Ultimate,
THOMPSON, JIM; BERBANK-GREEN, BARNABY; CUSWORTH, NIC. JOHN WILEY
CONSUMER, 2007.ISBN 0471968943.\\
Rules of Play: Game Design Fundamentals, Salen, K; e E, Zimmerman.
Hardcover, 2003. ISBN-10: 9780262240451.\\
Complementar:\\
Game Developer Magazine\\
Foruns de Desenvolvimento: www.gamasutra.com, www.gamedev.net e outros\\
(vários autores...) série: Game Programming Gems vol 1-6.

\newpage IDENTIFICAÇÃO

DISCIPLINA: Princípios da Animação Gráfica 3D CÓDIGO:\\ 
DEPARTAMENTO: Unidade Acadêmica de Garanhuns\\ 
ÁREA: \\
CARGA HORÁRIA TOTAL: 60 HORAS NÚMERO DE CRÉDITOS: 04\\
CARGA HORÁRIA SEMANAL: TEÓRICAS: PRÁTICAS: TOTAL: \\
PRÉ-REQUISITOS: Computação Gráfica
CO-REQUISITOS: não há

EMENTA 

Introdução a animação gráfica 3D. O processo de produção e princípios de
planejamento de Animações e Jogos 3D.Ferramentas de modelagem 3D.
Geração de animações em 3D. Ferramentas para desenvolvimento de Jogos
3D (Ogre 3D e XNA). Projeto de Desenvolvimento.

BIBLIOGRAFIA 

Básica:\\
Desenvolvimento De Jogos 3d E Aplicações Em Realidade Virtual, AZEVEDO,
EDUARDO; STELKO, MICHELLE; MEYER, HOMERO. CAMPUS, 2005. ISBN:
8535215697.\\
Allan Brito; Blender 3D -- Guia do Usuário (4. Edição), 2010 Ed.Novatec,
ISBN: 978-85-7522-258-4.\\
Reinicke, J. F.; Modelando Personagens com o Blender 3D, 2008 Ed.
Novatec, 978-85-7522-144-0. \\
Complementar:\\
Tutoriais e Fóruns de discussão.

\newpage IDENTIFICAÇÃO

DISCIPLINA: Wavelets CÓDIGO:\\ 
DEPARTAMENTO: Unidade Acadêmica de Garanhuns\\ ÁREA: \\
CARGA HORÁRIA TOTAL: 60 HORAS NÚMERO DE CRÉDITOS: 04\\
CARGA HORÁRIA SEMANAL: TEÓRICAS: PRÁTICAS: TOTAL: \\
PRÉ-REQUISITOS: Álgebra Linear; Algoritmos e Estrutura de Dados II
\\CO-REQUISITOS: não há

EMENTA 

Revisão em Álgebra Linear. Transformada Discreta de Fourier. Wavelets.

BIBLIOGRAFIA 

Básica:\\
Michael W. Frazier. An Introduction to Wavelets Through Linear Algebra.
Springer. 1999.\\
Christian Blatter. Wavelets: A Primer. A K Peters/CRC Press. 1999.\\
James S. Walker. A Primer on Wavelets and Their Scientific Applications.
CRC Press. 1999.\\
Complementar:\\
Gonzales \& Woods. Processamento Digital De Imagens (3 Ed). Pearson
Education. 2010. \\
A. K. Jain. Fundamentals of digital image processing. Prentice-Hall
International Editions, Englewood Cliffs, NJ, 1989.\\
Bernd Jähne. Digital Image Processing 6th edition. Springer. 2005.

\newpage IDENTIFICAÇÃO

DISCIPLINA: Introdução à Computação Quântica CÓDIGO:\\ 
DEPARTAMENTO: Unidade Acadêmica de Garanhuns\\ 
ÁREA: \\
CARGA HORÁRIA TOTAL: 60 HORAS NÚMERO DE CRÉDITOS: 04\\
CARGA HORÁRIA SEMANAL: 04 TEÓRICAS: 03 PRÁTICAS: 01 TOTAL: 04\\
PRÉ-REQUISITOS: não há\\
CO-REQUISITOS: não há

EMENTA 

Espaços de Hilbert sobre corpo Complexo. Elementos da Teoria da
Computação Clássica contendo Circuitos Booleanos. Elementos da Teoria Quântica. Elementos da
Computação quântica: modelos teóricos, portas lógicas quânticas. Algoritmos
quânticos do tipo Oráculo (Deutsch-Josza, Grover). Algorítmos quânticos do tipo Transformada de
Fourier (Simon, Shor). Simuladores e Linguagens de Programação Quânticas. Noções de
complexidade de computação: Classe NP, Algoritmos Probabilísticos e a Classe BPP.

BIBLIOGRAFIA 

Noson S. Yanofsky; Mirco A. Mannucci: Quantum Computing for Computer
Scientists.\\
Cambridge University Press, 2008, ISBN 978-0-521-87996-5 \\
David McMahon: Quantum Computing Explained. Wiley-Interscience,
Hoboken, New Jersey, USA, 2008, ISBN 978-0-470-09699-4\\
N. David Mermin: Quantum Computer Science  An Introduc-tion.
Cambridge University Press, New York, USA, 2007, ISBN 978-0-521-87658-2\\
Alexei Yu. Kitaev, Alexander H. Shen e Mikhail N. Vyalyi: Classical
and Quantum Computation. Graduate Studies in Mathematics, vol 47, AMS, 2002. ISBN
0-8218-3229-8\\
Michael A. Nielsen e Isaac L. Chuang: Computação Quântica e
Informação Quântica. 1a. Edição, Editora Bookman, 2005, ISBN 8536305541.\\
Artigos recentes da área de pesquisa.

\newpage IDENTIFICAÇÃO

DISCIPLINA: Tópicos Especiais em Processamento de Sinais CÓDIGO:\\ 
DEPARTAMENTO: Unidade Acadêmica de Garanhuns\\
ÁREA: \\
CARGA HORÁRIA TOTAL: 60 HORAS NÚMERO DE CRÉDITOS: 04\\
CARGA HORÁRIA SEMANAL: TEÓRICAS: PRÁTICAS: TOTAL: \\
PRÉ-REQUISITOS: Calculo II; Álgebra Linear\\
CO-REQUISITOS: não há

EMENTA 

Revisão em Álgebra Linear. Transformada Discreta de Fourier. Wavelets.
Introdução a analise de sinais Biométricos. 

BIBLIOGRAFIA 

Básica:\\
S. Allen Broughton and Kurt M. Bryan. Discrete Fourier Analysis and
Wavelets: Applications to Signal and Image Processing.
Wiley-Interscience, 2008.\\
John W. Woods. Multidimensional Signal, Image and Video Processing and
Coding. Academic Press, 2006.\\
Robert B. Northrop. Signals and Systems Analysis In Biomedical
Engineering, Second Edition  Hardcover, 2010.\\
Complementar:\\
Charles L. Byrne. Signal Processing: a Mathematical Approach. A. K.
Peters Ltd., 2005.\\
I. Gath, G.F. Inbar. Advances in Processing and Pattern Analysis of
Biological Signals (Advances in Experimental Medicine \& Biology.
Hardcover, 1996.\\
Ronald N. Bracewell. Fourier Analysis and Imaging. Springer, 2004.\\
Alan. V. Oppenheim and Ronald W. Schafer. Discrete-Time Signal
Processing, 2nd ed. Prentice Hall, 1999.

\newpage IDENTIFICAÇÃO

DISCIPLINA: Tópicos Especiais em Processamento de Imagens CÓDIGO:\\ 
DEPARTAMENTO: Unidade Acadêmica de Garanhuns\\ 
ÁREA: \\
CARGA HORÁRIA TOTAL: 60 HORAS NÚMERO DE CRÉDITOS: 04\\
CARGA HORÁRIA SEMANAL: TEÓRICAS: PRÁTICAS: TOTAL: \\
PRÉ-REQUISITOS: Processamento de Imagens.\\
CO-REQUISITOS: não há

EMENTA 

Estudo de técnicas avançadas em Processamento de Imagens permitindo ao
aluno conhecer o estado da arte nesta área de pesquisa.

BIBLIOGRAFIA 

Artigos recentes da área de pesquisa.

\newpage IDENTIFICAÇÃO

DISCIPLINA: Tópicos Especiais em Processamento de Áudio CÓDIGO:\\ 
DEPARTAMENTO: Unidade Acadêmica de Garanhuns\\ 
ÁREA: \\
CARGA HORÁRIA TOTAL: 60 HORAS NÚMERO DE CRÉDITOS: 04\\
CARGA HORÁRIA SEMANAL: TEÓRICAS: PRÁTICAS: TOTAL: \\
PRÉ-REQUISITOS: Processamento de Áudio\\
CO-REQUISITOS: não há

EMENTA 

Estudo de técnicas avançadas em Processamento de Áudio permitindo ao
aluno conhecer o estado da arte nesta área de pesquisa.

BIBLIOGRAFIA 

Artigos recentes da área de pesquisa.

\newpage IDENTIFICAÇÃO

DISCIPLINA: Tópicos Especiais em Mídia e Interação CÓDIGO:\\ 
DEPARTAMENTO: Unidade Acadêmica de Garanhuns\\
ÁREA: \\
CARGA HORÁRIA TOTAL: 60 HORAS NÚMERO DE CRÉDITOS: 04\\
CARGA HORÁRIA SEMANAL: TEÓRICAS: PRÁTICAS: TOTAL: \\
PRÉ-REQUISITOS: Não há.\\
CO-REQUISITOS: não há

EMENTA 

Estudo de técnicas avançadas em de interação e representação das
diversas mídias existentes, permitindo ao aluno conhecer o estado da
arte nesta área de pesquisa.

BIBLIOGRAFIA 

Artigos recentes da área de pesquisa.

\newpage IDENTIFICAÇÃO

DISCIPLINA: Tópicos Especiais em Jogos CÓDIGO:\\ 
DEPARTAMENTO: Unidade Acadêmica de Garanhuns\\
ÁREA: \\
CARGA HORÁRIA TOTAL: 60 HORAS NÚMERO DE CRÉDITOS: 04\\
CARGA HORÁRIA SEMANAL: TEÓRICAS: PRÁTICAS: TOTAL: \\
PRÉ-REQUISITOS: não há\\
CO-REQUISITOS: não há

EMENTA 

Estudo de técnicas para o gerenciamento de produção de jogos.
Ferramentas (Engineer) de desenvolvimento de jogos 3D. Projeto.

BIBLIOGRAFIA 

Rabin, Steve, 2a, Introduction to Game Development. Charles River Media
2010\\
Penton, Ron., 1a, Data structures for game programmers. The Premier
Press 2003\\
Artigos recentes da área de pesquisa.

\newpage IDENTIFICAÇÃO

DISCIPLINA: Tópicos Especiais em Realidade Virtual e Aumentada CÓDIGO:\\ 
DEPARTAMENTO: Unidade Acadêmica de Garanhuns\\
ÁREA: \\
CARGA HORÁRIA TOTAL: 60 HORAS NÚMERO DE CRÉDITOS: 04\\
CARGA HORÁRIA SEMANAL: TEÓRICAS: PRÁTICAS: TOTAL: \\
PRÉ-REQUISITOS: não há\\
CO-REQUISITOS: não há

EMENTA 

Estudo de técnicas avançadas em de interação e representação das
diversas mídias existentes, permitindo ao aluno conhecer o estado da
arte nesta área de pesquisa.

BIBLIOGRAFIA 

Grigore C. Burdea et al. Virtual RealityTechnology, 2nd. edition,
Wiley-Interscience, 2003.\\
Desenvolvimento De Jogos 3d E Aplicações Em Realidade Virtual, AZEVEDO,
EDUARDO; STELKO, MICHELLE; MEYER, HOMERO. CAMPUS, 2005. ISBN:
8535215697.\\
Oliver Bimber et al. Spatial Augmented Reality: Merging Real and Virtual
Worlds, A K Peters, 2005.\\
Artigos recentes da área de pesquisa.

\section*{Área Temática: Redes e Sistemas Distribuídos}

\newpage IDENTIFICAÇÃO

DISCIPLINA: Gerenciamento de Redes CÓDIGO:\\ 
DEPARTAMENTO: Unidade Acadêmica de Garanhuns\\
ÁREA: \\
CARGA HORÁRIA TOTAL: 60 HORAS NÚMERO DE CRÉDITOS: 04\\
CARGA HORÁRIA SEMANAL: TEÓRICAS: 3 PRÁTICAS: TOTAL: 1 \\
PRÉ-REQUISITOS: Redes de Computadores\\
CO-REQUISITOS: não há

EMENTA 

Princípios, organização e métodos de administração de rede; Tecnologias
para operação e gerência de rede; Rede de gerência de telecomunicações
TMN; Recursos humanos para administração de rede; Plataformas de
gerência de redes; e Aplicações de gerência de rede.

BIBLIOGRAFIA 

Básica:\\
Stallings, William. SNMP, SNMPv2, SNMPv3, and RMON 1 and RMON 2. Third
Edition. Pearson, 1999.\\
Burgess, Mark. Princípios de Administração de redes e sistemas. 3. ed.
São Paulo: LTC, 2006.\\
Kurose, James F.; Ross, Keith W. Redes de computadores e a internet: uma
abordagem top-down. 3. Edição. São Paulo: Pearson Addison Wesley,
2008.\\
Complementar:\\
Tanenbaum, Andrew S. Redes de computadores. 4. Edição. Rio de Janeiro:
Editora Campus, 2002. \\
COMER, Douglas E. Interligação de Redes com TCP/IP, Volumes I 5. Edição.
Prentice Hall, 2006.\\
COMER, Douglas E. Interligação de Redes com TCP/IP, Volumes II 5. Edição. Prentice Hall, 2006.

\newpage IDENTIFICAÇÃO

DISCIPLINA: Segurança de Redes CÓDIGO:\\ 
DEPARTAMENTO: Unidade Acadêmica de Garanhuns\\ 
ÁREA: \\
CARGA HORÁRIA TOTAL: 60 HORAS NÚMERO DE CRÉDITOS: 04\\
CARGA HORÁRIA SEMANAL: TEÓRICAS: 3 PRÁTICAS: TOTAL: 1 \\
PRÉ-REQUISITOS: Redes de Computadores\\
CO-REQUISITOS: não há

EMENTA 

Introdução à criptografia; Criptografia de chave simétrica; Criptografia
de chave pública; Assinatura Digital; Certificados Digitais;
Gerenciamento da Segurança da Informação; Pragas Virtuais; SPAM;
Mecanismos de Segurança;Firewall; e Redes privativas virtuais.

BIBLIOGRAFIA 

Básica:\\
Nakamura, Emilio T.; Geus, Paulo L. Segurança de Redes em Ambientes
Coorporativos. São Paulo: Editora Novatec , 2007. \\
Burnett, Steve; Paine, Stephen. Criptografia e Segurança. O Guia Oficial
RSA. 1. ed. Rio de Janeiro: Editora Campos, 2002.\\
Stallings, William. Cryptography and network security: principles and
practices. 14. ed. London: Pearson Prentice Hall, 2006.\\
Complementar:\\
Kurose, James F.; Ross, Keith W. Redes de computadores e a internet: uma
abordagem top-down. 3. Edição. São Paulo: Pearson Addison Wesley, 2008.\\
Tanenbaum, Andrew S. Redes de computadores. 4. Edição. Rio de Janeiro:
Editora Campus, 2002. \\
COMER, Douglas E. Interligação de Redes com TCP/IP, Volumes I,II 5.
Edição. Prentice Hall, 2006.

\newpage IDENTIFICAÇÃO

DISCIPLINA: Tópicos Avançados em Redes e Sistemas Distribuídos CÓDIGO:\\ 
DEPARTAMENTO: Unidade Acadêmica de Garanhuns\\ 
ÁREA: \\ 
CARGA HORÁRIA TOTAL: 60 HORAS NÚMERO DE CRÉDITOS: 04\\
CARGA HORÁRIA SEMANAL: TEÓRICAS: 3 PRÁTICAS: TOTAL: 1 \\
PRÉ-REQUISITOS: Sistemas Distribuídos\\
CO-REQUISITOS: não há

EMENTA 

Estudo de técnicas avançadas em Redes de Computadores e Sistemas
Distribuídos permitindo ao aluno conhecer o estado da arte nesta área
de pesquisa.

BIBLIOGRAFIA 

Básica:\\
Ad Hoc NetworksACM Transactions on Computer SystemsDistributed
Computing Distributed Computing (Internet)\\
Complementar:\\
Computational \& Applied MathematicsComputer Communication Review;
Computer Communications\\
\newpage 

\section*{Área Temática: Tecnologia Educacional}

IDENTIFICAÇÃO

DISCIPLINA: Desenvolvimento de Software Educacional CÓDIGO:\\
DEPARTAMENTO: Unidade Acadêmica de Garanhuns\\
ÁREA: \\
CARGA HORÁRIA TOTAL: 60 HORAS NÚMERO DE CRÉDITOS: 04\\
CARGA HORÁRIA SEMANAL: TEÓRICAS: PRÁTICAS: TOTAL: \\
PRÉ-REQUISITOS: Engenharia de Software\\
CO-REQUISITOS: não há

EMENTA 

O computador no processo ensino-aprendizagem; Concepções de
conhecimento e Teorias Cognitivas. Qualidade de software
(produto) e qualidade no desenvolvimento (processo). Modelos de
avaliação de softwares educacionais. Etapas para o
desenvolvimento de um software educacional (ciclo de vida).

BIBLIOGRAFIA 

Básica:\\
BARBOSA, S. D. J.; Silva, B. S. Interação Humano-Computador. Rio de
Janeiro: Elservier, 2010.\\
PREECE, Rogers e Sharp; Design de Interação  Além da Interação
Homem-computador. São Paulo: Bookman, 2005.\\
MAYER. R, Cambridge Handbook of Multimedia Learning. New York: Cambridge
University Press, 2005\\
Complementar:\\
COX, Kenia. Informática na Educação Escolar, Autores Associados, 2003.\\
LEVY, Pierre. Cibercultura, Editora 34, 1999.\\
RBIE -- Revista Brasileira de Informática na Educação ISSN 1414-5685\\
RENOTE -- Revista Novas Tecnologias na Educação ISSN 1679-1916.

\newpage IDENTIFICAÇÃO

DISCIPLINA: Informática na Educação CÓDIGO:\\ 
DEPARTAMENTO: Unidade Acadêmica de Garanhuns\\
ÁREA: \\
CARGA HORÁRIA TOTAL: 60 HORAS NÚMERO DE CRÉDITOS: 04\\
CARGA HORÁRIA SEMANAL: TEÓRICAS: PRÁTICAS: TOTAL:\\
PRÉ-REQUISITOS: Introdução a Ciência da Computação\\
CO-REQUISITOS: não há

EMENTA 

Motivação e discussão crítica sobre o uso da informática na educação,
incluindo conhecimento sobre assuntos atuais; Histórico da informática
na educação; Ambientes educacionais baseados em computador; As
implicações pedagógicas e sociais do uso da informática na educação;
Informática na educação especial, na educação à distância e no
aprendizado cooperativo.

BIBLIOGRAFIA 

Básica:\\
FRANCO, Sergio R.K. Informática na Educação: Estudos Interdisciplinares.
Editora da UFRGS, 2004\\
LEVY, Pierre -- As tecnologias da Inteligência- O futuro do pensamento
na era da informática. São Paulo: Editora 34, 2004, 13a. Edição.\\
PAPERT, Seymour. A máquina das crianças: repensando a escola na era da
informática. Porto Alegre: Artes Médicas, 1994.\\
Complementar:\\
COX, Kenia. Informática na Educação Escolar, Autores Associados, 2003.\\
LEVY, Pierre. Cibercultura, Editora 34, 1999.\\
RBIE -- Revista Brasileira de Informática na Educação ISSN 1414-5685\\
RENOTE -- Revista Novas Tecnologias na Educação ISSN 1679-1916.

\newpage IDENTIFICAÇÃO

DISCIPLINA: Linguagem Brasileira de Sinais CÓDIGO:\\ 
DEPARTAMENTO: Unidade Acadêmica de Garanhuns\\
ÁREA: \\
CARGA HORÁRIA TOTAL: 60 HORAS NÚMERO DE CRÉDITOS: 04\\
CARGA HORÁRIA SEMANAL: TEÓRICAS: PRÁTICAS: TOTAL: \\
PRÉ-REQUISITOS: não há\\
CO-REQUISITOS: não há

EMENTA 

Abordagem social e lingüística da Surdez e do mundo dos Surdos com
ênfase nos aspectos específicos da Língua Brasileira de Sinais --
LIBRAS. Surdez, linguagem e sociedade, contextos sociais, históricos e
políticos, Comunidade, cultura, identidade, representação e alteridade Surda 
Estrutura lingüística e gramatical da LIBRAS, com produção de diálogos e
textos.

BIBLIOGRAFIA 

Básica:\\
BRITO, Lucinda Ferreira (Org.). Brasil, Secretaria de Educação Especial.
Língua Brasileira de Sinais. Volume III. Série Atualidades Pedagógicas,
n. 4. Brasília EC/SCESP, 1998\\
FELIPE, Tanya Amara. Libras em contexto -- Curso Básico. Rio de Janeiro
MEC/FNDE/SEESP, 1999.\\
FERNANDES,Eulália. Linguagem e surdez. Porto Alegre: Artmed, 2003\\
Complementar:\\
BERNARDINO, Elidéa Lúcia. Absurdo ou lógica? A produção lingüística do
surdo. Belo Horizonte , Editora Profetizando vidas, 2000\\
OARES, Maria Aparecida Leite. Educação de Surdos no Brasil.
SãoPaulo:EDUSF ; Autores Associados, 1999\\
SKILIAR, Carlos (Org.). Educação e Exclusão: Abordagens Sócio 
Antropológicas em Educação Especial. Porto Alegre: Mediação, 1999\\
SOARES, Maria Aparecida Leite. Educação de Surdos no Brasil. São Paulo:
EDUSF, Autores Associados, 1999.

\newpage IDENTIFICAÇÃO

DISCIPLINA: Projeto de Sistemas Educacionais CÓDIGO:\\ 
DEPARTAMENTO: Unidade Acadêmica de Garanhuns\\
ÁREA: \\
CARGA HORÁRIA TOTAL: 60 HORAS NÚMERO DE CRÉDITOS: 04\\
CARGA HORÁRIA SEMANAL: TEÓRICAS: PRÁTICAS: TOTAL: \\
PRÉ-REQUISITOS: Engenharia de Software\\
CO-REQUISITOS: não há

EMENTA 

Revisão de tópicos da Engenharia de Software; Introdução à
Sistemas e Ciclo de Vida de um Sistema de Informação;  
Necessidades Básicas para o Desenvolvimento de Sistemas Educacionais; Elicitação e Validação de Requisitos para Desenvolvimento de
Sistemas Educacionais.

BIBLIOGRAFIA 

Básica:\\
SOMMERVILLE, Ian. Engenharia de Software. ISBN.: 8588639076. Addison --
Wesley.\\
BARBOSA, S. D. J.; Silva, B. S. Interação Humano-Computador. Rio de
Janeiro: Elservier, 2010.\\
MAYER. R, Cambridge Handbook of Multimedia Learning. New York: Cambridge
University Press, 2005\\
Complementar:\\
COX, Kenia. Informática na Educação Escolar, Autores Associados, 2003.\\
LEVY, Pierre. Cibercultura, Editora 34, 1999.\\
RBIE -- Revista Brasileira de Informática na Educação ISSN 1414-5685\\
RENOTE -- Revista Novas Tecnologias na Educação ISSN 1679-1916.\\

\newpage IDENTIFICAÇÃO

DISCIPLINA: Tecnologia Educacional e Cognição CÓDIGO:\\ 
DEPARTAMENTO: Unidade Acadêmica de Garanhuns\\
ÁREA: \\
CARGA HORÁRIA TOTAL: 60 HORAS NÚMERO DE CRÉDITOS: 04\\
CARGA HORÁRIA SEMANAL: TEÓRICAS: PRÁTICAS: TOTAL: \\
PRÉ-REQUISITOS: Introdução a Ciência da Computação\\
CO-REQUISITOS: não há

EMENTA 

Diferentes concepções de Aprendizagem; Teoria behaviorista, humanistas e
cognitivistas da aprendizagem e suas implicações para Informática na
Educação.

BIBLIOGRAFIA 

Básica:\\
COLL, Cesar. PALACIOS, Jesus. MARCHESI, Alvaro (org.). Desenvolvimento
psicológico e educação: psicologia da educação. Porto Alegre: Artes
Medicas, 1996. 460p. v. 1\\
COLL, Cesar. PALACIOS, Jesus. MARCHESI, Alvaro (org.). Desenvolvimento
psicológico e educação: psicologia da educação. Porto Alegre: Artes
Medicas, 1996. 460p. v. 2.\\
MOREIRA, Marco. Teorias de aprendizagem. São Paulo: EPU, 1999.\\
Complementar:\\
COX, Kenia. Informática na Educação Escolar, Autores Associados, 2003.\\
LEVY, Pierre. Cibercultura, Editora 34, 1999.\\
RBIE -- Revista Brasileira de Informática na Educação ISSN 1414-5685\\
RENOTE -- Revista Novas Tecnologias na Educação ISSN 1679-1916.

\newpage IDENTIFICAÇÃO

DISCIPLINA: Tópicos Especiais em Tecnologia Educacional CÓDIGO:\\ 
DEPARTAMENTO: Unidade Acadêmica de Garanhuns\\ 
ÁREA: \\
CARGA HORÁRIA TOTAL: 60 HORAS NÚMERO DE CRÉDITOS: 04\\
CARGA HORÁRIA SEMANAL: TEÓRICAS: PRÁTICAS: TOTAL: \\
PRÉ-REQUISITOS: Introdução a Ciência da Computação\\
CO-REQUISITOS: não há

EMENTA 

Estudo de teorias de aprendizagem multimídia. Elaboração de projetos
pedagógicos utilizando recursos digitais. Novas tecnologias da
informação aplicadas ao ensino formal e informal.

BIBLIOGRAFIA 

Básica:\\
FRANCO, Sergio R.K. Informática na Educação: Estudos Interdisciplinares.
Editora da UFRGS, 2004\\
LEVY, Pierre -- As tecnologias da Inteligência- O futuro do pensamento
na era da informática. São Paulo: Editora 34, 2004, 13a. Edição.\\
PAPERT, Seymour. A máquina das crianças: repensando a escola na era da
informática. Porto Alegre: Artes Médicas, 1994.\\
Complementar:\\
ARAUJO, Alberto.E.P. Apostila de Informática na Educação I. UFRPE.
2008.\\
COX, Kenia. Informática na Educação Escolar, Autores Associados, 2003.\\
LEVY, Pierre. Cibercultura, Editora 34, 1999.\\
RBIE -- Revista Brasileira de Informática na Educação ISSN 1414-5685\\
RENOTE -- Revista Novas Tecnologias na Educação ISSN 1679-1916.
\newpage

\section*{Área Temática: Tecnologias da Informação}

IDENTIFICAÇÃO

DISCIPLINA: Gestão da Tecnologia da Informação CÓDIGO:\\ 
DEPARTAMENTO: Unidade Acadêmica de Garanhuns\\
ÁREA: \\
CARGA HORÁRIA TOTAL: 60 HORAS NÚMERO DE CRÉDITOS: 04\\
CARGA HORÁRIA SEMANAL: TEÓRICAS: PRÁTICAS: TOTAL: \\
PRÉ-REQUISITOS: Sistemas de Informação e Tecnologias\\
CO-REQUISITOS: não há

EMENTA 

Estratégia empresarial . Alinhamento estratégico; Processo de aquisição
e de implementação de sistemas; Desenvolvimento interno, outsourcing,
Software as a Service; Acordos de nível de serviço. Arquitetura
orientada a serviço; Tópicos em segurança da informação, falhas em SIs;
continuidade do negócio; Risco em TI; Uso de tecnologias de
Virtualização, datacenter; cloud computing; Valor da TI; Retorno sobre
investimento. Impactos da TI nos indivíduos, organizações. Questões
éticas.

BIBLIOGRAFIA 

Básica:\\
TURBAN, E.; MCLEAN, E; WETHERBE, J. Tecnologia da Informação para
Gestão. Transformando os Negócios na Economia Digital. Tradução de
Renate Schinke. Revisão técnica de Ângela F. Brodbeck. Porto Alegre:
Bookman, 2004.\\
MAGALHÃES, Ivan Luizio; PINHEIRO, Walfrido Brito. Gerenciamento de
Serviços de TI na Prática  Uma abordagem com base na ITIL. São
Paulo: Novatec, 2007.\\
Complementar:\\
Artigos de conferências e periódicos da área.

\newpage IDENTIFICAÇÃO

DISCIPLINA: Gestão de Processos de Negócio CÓDIGO:\\ 
DEPARTAMENTO: Unidade Acadêmica de Garanhuns\\
ÁREA: \\
CARGA HORÁRIA TOTAL: 60 HORAS NÚMERO DE CRÉDITOS: 04\\
CARGA HORÁRIA SEMANAL: TEÓRICAS: PRÁTICAS: TOTAL: \\
PRÉ-REQUISITOS: Sistemas de Informação e Tecnologias\\
CO-REQUISITOS: não há

EMENTA 

Contextualizando o gerenciamento de processos. Conceitos de processos.
Engenharia de Processos de Negócios: Desenho, Ferramentas,
Metodologias, Suporte de TI para Engenharia de Processos. Sistemas de
Informação e os processos organizacionais. A prática da modelagem de
processos. BPM e BPMS;Conceitos básicos: ciclo de vida de BPM; BPMS e
serviços Web; Modelagem de processos: BPMN.

BIBLIOGRAFIA 

Básica:\\
BALDAM R. et al. Gerenciamento de Processos de Negócio. São Paulo:
Érica, 2008, 240 p.\\
DAVENPORT, H. T. Reengenharia de processos: como inovar na empresa
através de tecnologia de informação. Rio de Janeiro, Campus, 1994, 391
p.\\
PAIM, R.,et aL. Gestão de Processos: Pensar, Agir e Aprender, Porto
Alegre: Bookman, 2009, 327 p.\\
Complementar:\\
VALLE, R.; OLIVEIRA, B. S. Análise e Modelagem de Processos de Negócio
-- Foco na notação BPMN. São Paulo: Atlas, 2009, 207 p.\\
BURLTON, Roger. Business Process Management: profiting from process.
Indianapolis: Sams Publishing, 2001.

\newpage IDENTIFICAÇÃO

DISCIPLINA: Gestão de Serviços em TI CÓDIGO:\\ 
DEPARTAMENTO: Unidade Acadêmica de Garanhuns\\
ÁREA: \\
CARGA HORÁRIA TOTAL: 60 HORAS NÚMERO DE CRÉDITOS: 04\\
CARGA HORÁRIA SEMANAL: TEÓRICAS: PRÁTICAS: TOTAL: \\
PRÉ-REQUISITOS: Sistemas de Informação e Tecnologias\\
CO-REQUISITOS: não há

EMENTA 

Conceitos do que é um Serviço. Características do Serviço. Governança de
TI e Gerenciamento de Serviços. Gestão estratégica e tática de serviços
de TI. Acordos de nível de serviço. Gerenciamento de serviços com base
no conjunto de melhores práticas baseado no ITIL (Information
Technology Infrastructure Library Biblioteca de Infra-estrutura de
Tecnologia da Informação) que identifica o relacionamento das diversas
atividades necessárias para entrega e suporte dos serviços de TI.
Ferramentas de apoio ao gerenciamento de serviços. Elaboração de
Projeto.

BIBLIOGRAFIA 

Básica:\\
JANDER, M. , STURM , R. e MORRIS , W. Service Level Management 
Fundamentos Do Gerenciamento De Níveis De Serviço. São Paulo: Campus,
2000.\\
REIS, A. Alinhamento da estratégia de Ti com a estratégia corporativa;
São Paulo: Atlas, 2003.\\
MAGALHÃES, Ivan Luizio; PINHEIRO, Walfrido Brito. Gerenciamento de
Serviços de TI na Prática  Uma abordagem com base na ITIL. São
Paulo: Novatec, 2007.\\
Complementar:\\
Artigos de conferências e periódicos da área.

\newpage IDENTIFICAÇÃO

DISCIPLINA: Governança em Tecnologia da Informação CÓDIGO:\\ 
DEPARTAMENTO: Unidade Acadêmica de Garanhuns\\ 
ÁREA: \\
CARGA HORÁRIA TOTAL: 60 HORAS NÚMERO DE CRÉDITOS: 04\\
CARGA HORÁRIA SEMANAL: TEÓRICAS: PRÁTICAS: TOTAL: \\
PRÉ-REQUISITOS: Sistemas de Informação e Tecnologias\\
CO-REQUISITOS: não há

EMENTA 

Planejamento Estratégico; Alinhamento estratégico. Decisões Estratégicas
de TI. Governança Corporativa e Governança de TI. Arquétipos de TI para
alocação de direitos decisórios. Mecanismos para implantar a Governança
de TI. Tipos de governança. Associação da Estratégia, da Governança e o
Desempenho. Princípios de Liderança para governança de TI. Normas,
processos e indicadores de desempenho para a área de TI. Modelos de
apoio para Governança de TI: COBIT (Control Objectives for Information
and Related Technology); ITIL (Information Technology Infrastructure
Library); BSC (Balanced Scorecard). Estruturação de um plano de
implantação de um modelo de governança de TI. A norma ISO 20000.

BIBLIOGRAFIA 

Básica:\\
WEILL, Peter; ROSS, Jeanne W. Governança de TI: Tecnologia da
Informação. São Paulo: Makron Books, 2006.\\
MAGALHÃES, Ivan Luizio; PINHEIRO, Walfrido Brito. Gerenciamento de
Serviços de TI na Prática  Uma abordagem com base na ITIL. São
Paulo: Novatec, 2007.\\
FERNANDES, Aguinaldo Aragon; ABREU, Vladimir Ferraz. Implantando a
Governança de TI -- da Estratégia à Gestão dos Processos e Serviços.
Rio de Janeiro: Brasport, 2006.\\
Complementar:\\
Artigos de conferências e periódicos da área.

\newpage IDENTIFICAÇÃO

DISCIPLINA: Contabilidade e Administração de Empresas CÓDIGO:\\ 
DEPARTAMENTO: Unidade Acadêmica de Garanhuns\\
ÁREA: \\
CARGA HORÁRIA TOTAL: 60 HORAS NÚMERO DE CRÉDITOS: 04\\
CARGA HORÁRIA SEMANAL: TEÓRICAS: PRÁTICAS: TOTAL: \\
PRÉ-REQUISITOS: não há\\
CO-REQUISITOS: não há

EMENTA 

Empreendedorismo; Características; Oportunidades; Desenvolvimento de
Atitudes Empreendedoras. Novos Paradigmas. Administração do Crescimento
da Empresa. Prospecção Empresarial. Plano de Negócio. Inovação e
Criatividade. Modelagem Organizacional. Pesquisa de Mercado. Técnicas
de Venda. Técnicas de Negociação. Qualidade. Formação de Preços.
Ferramentas Gerenciais.

BIBLIOGRAFIA 

Básica:
DOLABELA, Fernando. Oficina do empreendedor. São Paulo: Cultura, 1999.\\
MATTAR, Fauze Najib. Pesquisa de marketing. 2. ed., São Paulo: Atlas,
2000.\\
MONTGOMERY, Cynthia A.; PORTER, Michael (Org.). Estratégia: a busca da
vantagem competitiva. 5. ed., Rio de Janeiro: Campus, 2000.
Complementar:\\
BEEMER, C. Britt; SHOOK, Robert L. Marketing estratégico: tudo o que
mega e micro empresários devem saber para conquistar novos clientes.
São Paulo: Futura, 1998.\\
ÂNGELO, Cláudio Felisoni de (Coord.). Varejo: modernização e
perspectivas. São Paulo: Atlas, 1994.

\newpage IDENTIFICAÇÃO

DISCIPLINA: Contabilidade e Administração Financeira CÓDIGO:\\ 
DEPARTAMENTO: Unidade Acadêmica de Garanhuns\\
ÁREA: \\
CARGA HORÁRIA TOTAL: 60 HORAS NÚMERO DE CRÉDITOS: 04\\
CARGA HORÁRIA SEMANAL: TEÓRICAS: PRÁTICAS: TOTAL: \\
PRÉ-REQUISITOS: não há\\
CO-REQUISITOS: não há

EMENTA 

Administração de capital de giro. Administração de caixa e títulos
negociáveis. Administração de duplicatas a receber e do estoque. Fontes
de Financiamento em curto prazo. Fontes de financiamento com garantia
de curto prazo.

BIBLIOGRAFIA 

ASSAF, Neto Alexandre. Administração de Capital de Giro. São Paulo:
Atlas, 1999.\\
ASSEF, Roberto. Guia prático de administração financeira. RJ: Campus,
2003.\\
BRAGA, Roberto. Fund. e técnicas de administração financeira. SP:Atlas,
1998.\\
BRIGHAM, Eugene F. \& HOUSTON, Joel F. Fundamentos da moderna
administração financeira. Rio de Janeiro : Campus, 1999.\\
CROSWFORD, Richard. Na era do capital Humano: o talento, a inteligência
e o conhecimento como forças econômicas, seu impacto nas empresas e nas
decisões de investimento. São Paulo: Atlas, 2000.

\newpage IDENTIFICAÇÃO

DISCIPLINA: Tópicos Especiais em Gestão de Projetos CÓDIGO:\\ 
DEPARTAMENTO: Unidade Acadêmica de Garanhuns\\ ÁREA: \\
CARGA HORÁRIA TOTAL: 60 HORAS NÚMERO DE CRÉDITOS: 04\\
CARGA HORÁRIA SEMANAL: TEÓRICAS: PRÁTICAS: TOTAL: \\
PRÉ-REQUISITOS: Engenharia de Software\\
CO-REQUISITOS: não há

EMENTA 

Breve histórico da evolução do gerenciamento de projetos; Modelos de
gerenciamento de projetos: Revisão PMBoK e dos conceitos de
metodologias ágeis; Categorização de projetos.Modelos de Maturidade.
Inovação, complexidade e incerteza em projetos. Sustentabilidade em
projetos. Projetos globais. Ferramentas de gerenciamento de projetos.
Apresentação das áreas de pesquisa em projetos.

BIBLIOGRAFIA 

Básica:\\
Artigos de conferências e periódicos da área.\\
Complementar:\\
Periódicos.

\end{apendicesenv}
% ---


% ----------------------------------------------------------
% Anexos
% ----------------------------------------------------------

% ---
% Inicia os anexos
% ---
% \begin{anexosenv}

% % % Imprime uma página indicando o início dos anexos
% \partanexos

% % % ---
% \chapter{Creditação da Carga Horária de Extensão}
% \label{anexo:creditacao_extensao}
% % % ---

% \section*{Normativa UFAPE/BCC n. 1, de 19 de outubro de 2022, para creditação da carga horária de extensão.}

% \hfill\begin{minipage}[t]{8cm}
%   Dispõe sobre as Ações de extensão como componente curricular no âmbito do curso de Bacharelado em Ciência da Computação, da Universidade Federal do Agreste de Pernambuco, conforme Res. UFAPE/CONSEPE n.~7, de 22 de agosto de 2022.
% \end{minipage}

% \vspace{12pt}

% O Colegiado do curso de Bacharelado em Ciência da Computação, da Universidade Federal do Agreste de Pernambuco, reunido em sessão extraordinária, no dia 19 de outubro de 2022, no uso de suas atribuições legais.

% \vspace{12pt}

% \noindent RESOLVE:

% \vspace{12pt}

% \noindent \textbf{Art. 1º}. Regulamentar as Ações de Extensão como componente curricular do curso de Bacharelado em Ciência da Computação, da Universidade Federal do Agreste de Pernambuco, sendo indispensável o seu integral cumprimento para a conclusão do curso, e consequente colação de grau.

% \section*{Capítulo I. Das disposições preliminares}

% \noindent \textbf{Art. 2º}. A creditação da carga horária de extensão no âmbito do curso de Bacharelado em Ciência da Computação, da Universidade Federal do Agreste de Pernambuco, corresponde ao mínimo de 10\% (dez por cento) da carga horária total do curso, sendo obrigatória, conforme estabelece a Res. UFAPE/CONSEPE n. 7, de 22 de agosto de 2022, em consonância com a Res. CNE/CES/MEC n. 7, de 18 de dezembro de 2018.

% \section*{Capítulo II. Da forma de creditação}

% \noindent \textbf{Art. 3º}. A carga horária de extensão a ser creditada no currículo do curso de Bacharelado em Ciência da Computação deve ser desenvolvida de forma a destacar o discente como protagonista da ação de extensão, totalizando o mínimo de 320 (trezentos e vinte horas).

% \noindent \textbf{Art. 4º}. O curso promoverá a creditação da carga horária de extensão considerando as modalidades previstas no Art. 11, Capítulo II, da Res. UFAPE/CONSEPE n. 7, de 22 de agosto de 2022:

% \noindent I - ACEX I – participação do/a estudante como membro da equipe executora dos Projetos ou Programas de Extensão institucionalizados na UFAPE, com ou sem bolsa e com ou sem parceria com outras instituições público-privadas;

% \noindent II - ACEX II – participação do/a estudante como membro organizador e/ou ministrante de Cursos, Palestras e Eventos ou na Prestação de Serviço, que tenham sido institucionalizados e vinculados a Projetos ou Programas de Extensão, conforme artigo 3°; e

% \noindent III - ACEX III – participação do/a estudante como membro da equipe executora dos Projetos ou Programas de Extensão em outras Instituições de Ensino Superior (IES) pública ou privada, dentro ou fora do Brasil. 

% \noindent \textbf{§ 1º}. O pedido de aproveitamento, juntamente com a documentação comprobatória (comprovação de cadastro junto à Pró-Reitoria de Extensão e Cultura, da Universidade Federal do Agreste de Pernambuco, e declaração/certificado de participação) que se refere o Art. 4° deverá ser encaminhado à Coordenação de Curso e avaliado pelo Colegiado de Coordenação Didática do Curso de Bacharelado em Ciência da Computação, da Universidade Federal do Agreste de Pernambuco.

% \section*{Capítulo V. Das disposições finais}

% \noindent \textbf{Art. 5º}. Os documentos comprobatórios para fins de creditação da carga horária de extensão deverão ser encaminhados à Coordenação do Curso.

% \noindent \textbf{Art. 6º}. Somente terão validade, para fins de creditação da carga horária de extensão, a participação em Programas e Projetos, durante o período de matrícula no curso de Bacharelado em Ciência da Computação, da Universidade Federal do Agreste de Pernambuco.

% \noindent \textbf{Art. 7º}. Os casos omissos deverão ser apresentados e discutidos no Núcleo Docente Estruturante para posterior encaminhamento e deliberação do Colegiado de Coordenação Didática do Curso de Bacharelado em Ciência da Computação, da Universidade Federal do Agreste de Pernambuco.

% % % ---
% % \chapter{Cras non urna sed feugiat cum sociis natoque penatibus et magnis dis
% % parturient montes nascetur ridiculus mus}
% % % ---

% % \lipsum[31]

% % % ---
% % \chapter{Fusce facilisis lacinia dui}
% % % ---

% % \lipsum[32]

% \end{anexosenv}

%---------------------------------------------------------------------
% INDICE REMISSIVO
%---------------------------------------------------------------------
% \phantompart
% \printindex
%---------------------------------------------------------------------

\end{document}